%% Chapter 4 — Noncommutative Spectral Geometry on Warped Orbifolds
%% Paper IV of the Meridian Monograph
%% To be \input{} into the main book class document

\chapter{Noncommutative Spectral Geometry on Warped Orbifolds: Topological Couplings and Gravitational Corrections}
\label{ch:ncg}

\begin{quote}\small\itshape
We construct the spectral triple for the warped orbifold $M_4 \times [0, y_c] \times F$ with $S^1/\mathbb{Z}_2$ identification and compute its spectral action through the Seeley--DeWitt heat kernel expansion. The warped five-dimensional geometry is encoded in a Dirac operator $D_5 = e^{-A(y)} \widetilde{D}_4 + \gamma^5(\partial_y + 2A'(y))$, where the spin-warp coupling $2A'(y)$ arises from the 5D spin connection. We resolve the odd-dimension product obstruction---the absence of a grading operator in 5D prevents a standard product $D_5 \otimes D_F$---through a layered spectral architecture in which the gravitational sector (5D bulk, KO-dimension~5) and gauge sector (4D brane $\times\, F$, KO-dimension~4) are coupled through Israel junction conditions. The spectral action produces the following physical content: (i)~the Einstein--Hilbert action from the $a_2$ coefficient, which determines the spectral cutoff $\Lambda \sim 10^{17}\;\mathrm{GeV}$; (ii)~a dynamical Gauss--Bonnet correction from $a_3$ with dimensionless coupling $\hat{\alpha} \sim 10^{-2}$, which we show is approximately universal (within a factor of 2--3) across compactification dimensions $d = 5$ to $d = 10$ via the Gilkey coefficient structure; (iii)~a DGP-type induced gravity term on each brane from $a_{5/2}$ with crossover scale $L_c \sim 10^{-6}\;\mathrm{m}$; and (iv)~gravitational and gauge Chern--Simons terms from $a_{7/2}$, topologically locked by the Atiyah--Patodi--Singer index theorem. We show that the gauge theta angle $\theta_{\mathrm{EM}} \sim \mathcal{O}(1)$ is unsuppressed on the IR brane (due to warped redshift of the spectral cutoff), while the gravitational theta angle $\theta_{\mathrm{grav}} \sim 10^{-15}$ is Planck-suppressed. Five perturbative EM-gravity coupling channels are evaluated and killed (suppressions ranging from $10^{-28}$ to $10^{-97}$), while the non-perturbative topological channel---mediated by spectral flow under Lorentzian topological charge---remains formally open. The layered architecture naturally enforces the observed pattern of gravity in the bulk and gauge fields on the brane, with the Standard Model arising from the Chamseddine--Connes construction on the brane spectral triple.
\end{quote}

%% ============================================================
\section{Introduction}
\label{sec:4-intro}

Noncommutative geometry (NCG), in the framework developed by Connes~\cite{ch4:Connes1994}, provides a spectral characterization of geometry in which the Riemannian structure of a manifold is encoded in the spectrum of a Dirac-type operator rather than in a metric tensor. The Connes reconstruction theorem~\cite{ch4:Connes2013} establishes that a commutative spectral triple $(\mathcal{A}, \mathcal{H}, D)$ with algebra $\mathcal{A} = C^\infty(M)$ determines the Riemannian manifold~$M$ up to isometry. The power of this approach lies in its natural generalization: replacing the commutative algebra with a noncommutative one produces geometries richer than smooth manifolds---geometries that accommodate gauge fields and the full Standard Model of particle physics~\cite{ch4:CC1997,ch4:CC1996}.

The spectral action principle~\cite{ch4:CC1997} states that the bosonic action is determined by the spectrum of the Dirac operator alone:
\begin{equation}
  S = \mathrm{Tr}\!\left[f\!\left(\frac{D^2}{\Lambda^2}\right)\right],
  \label{eq:4-1}
\end{equation}
where $\Lambda$ is a spectral cutoff and $f$ is a smooth positive function. Chamseddine and Connes~\cite{ch4:CC1997} showed that on the product geometry $M_4 \times F$ (where $F$ is a finite internal NCG space encoding the Standard Model gauge group $\mathrm{SU}(3)_c \times \mathrm{SU}(2)_L \times \mathrm{U}(1)_Y$), equation~\eqref{eq:4-1} reproduces the full Standard Model Lagrangian coupled to Einstein gravity. Subsequent work extended this to neutrino masses~\cite{ch4:CCM2007}, the seesaw mechanism, and Majorana couplings, while van Suijlekom~\cite{ch4:vanSuijlekom2015} and Connes--Marcolli~\cite{ch4:ConnesMarcolli2008} provided comprehensive mathematical treatments.

In this chapter, we extend the NCG spectral action construction to the warped orbifold geometry $M_4 \times [0, y_c]$ with $S^1/\mathbb{Z}_2$ identification---the Randall--Sundrum (RS) background~\cite{ch4:RS1999a,ch4:RS1999b} that underlies the self-tuning cosmology developed in Chapters~\ref{ch:foundation}--\ref{ch:nogo}~\cite{ch4:PaperI,ch4:PaperII,ch4:PaperIII}. The warped geometry introduces several features absent in the flat $M_4 \times F$ construction:

\begin{enumerate}
  \item \textbf{Spin-warp coupling.} The 5D spin connection mixes 4D and extra-dimensional spinor components through $\omega^{a5}_\mu = A'(y)\, e^{A(y)}\, \tilde{e}^a_\mu$, producing a term $2A'(y)$ in the Dirac operator.

  \item \textbf{Boundary contributions.} The orbifold fixed points (branes) contribute boundary Seeley--DeWitt coefficients $a_{n+1/2}$ producing brane-localized gravitational and topological terms.

  \item \textbf{Odd-dimension obstruction.} The 5D bulk has KO-dimension~5, which lacks the grading operator $\gamma$ required for the standard tensor product $D_5 \otimes \mathbf{1}_F + \gamma \otimes D_F$.

  \item \textbf{Warped spectral cutoff.} The physical cutoff on the IR brane is $\Lambda_{\mathrm{IR}} = \Lambda \cdot e^{-ky_c} \sim 1\;\mathrm{GeV}$, qualitatively different from the bulk cutoff $\Lambda \sim 10^{17}\;\mathrm{GeV}$.
\end{enumerate}

We resolve these issues and compute the complete spectral action, extracting both the well-known gravitational content (Einstein--Hilbert, cosmological constant) and the novel features specific to the warped orbifold: the Gauss--Bonnet coupling, brane-localized induced gravity, Chern--Simons topological terms, and the electromagnetic-gravity coupling channels they produce.

\paragraph{Epistemic status.}
This chapter extends the cosmological framework of Chapters~1--3 into particle physics territory. The extension is speculative in character: while the geometric foundations (the warped RS orbifold, the spectral triple, the Seeley--DeWitt expansion) are shared with the cosmological core, the particle phenomenology introduces additional inputs (bulk mass parameters, octonionic algebraic structure) and contains several open computations (Poincar\'{e} duality for non-associative algebras, gauge coupling unification, full CKM fit). The results should be read as a research program built on the same geometric foundations, not as derived consequences of the cosmological axioms A1--A6.

\paragraph{Outline.}
Section~\ref{sec:4-spectral-triple} constructs the spectral triple and resolves the odd-dimension obstruction. Section~\ref{sec:4-seeley-dewitt} computes the Seeley--DeWitt coefficients. Section~\ref{sec:4-gb-universality} proves the universality of the Gauss--Bonnet coupling across dimensions and topologies. Section~\ref{sec:4-chern-simons} derives the Chern--Simons terms and their theta angles from the APS index theorem. Section~\ref{sec:4-em-gravity} analyzes perturbative and non-perturbative EM-gravity coupling channels. Section~\ref{sec:4-sm-connection} discusses the Standard Model connection, anomaly cancellation (Section~\ref{sec:4-anomaly}), and Chern--Simons inflow consistency (Section~\ref{sec:4-cs-inflow}).  Section~\ref{sec:4-gauge-unification} confronts the gauge coupling unification prediction with SM running and identifies resolution pathways.  Section~\ref{sec:4-experiment-forecast} collects the neutrino experiment forecasts.  Section~\ref{sec:4-conclusions} summarizes.

\paragraph{Key results from companion chapters.}
This chapter builds on and feeds back into the other chapters in this monograph. From Chapter~\ref{ch:foundation}~\cite{ch4:PaperI}, we use the warped 5D background geometry (Section~II), the Israel junction conditions at the orbifold branes (Section~IV.D), and the cuscuton effective action with non-minimal coupling $F(\phi) = M_{\mathrm{Pl}}^2 - \xi\phi^2$ (Section~III). From Chapter~\ref{ch:foundation}~\cite{ch4:PaperI}, the dimensionless non-minimal coupling $\zeta_0 = \xi\phi_0^2/M_5^3$ is a free parameter determined by brane UV physics (junction conditions at the orbifold boundaries). From Chapter~\ref{ch:nogo}~\cite{ch4:PaperIII}, we use the radiative stability of the cuscuton constraint (three independent proofs: symmetry protection, Dirac constraint topology, geometric origin), which bounds the quantization channel to $|\delta w| < 10^{-6}$.

The principal outputs of this chapter that enter other chapters are: (i)~the Gauss--Bonnet coupling $\hat{\alpha} \sim 10^{-2}$ and its approximate universality across compactification dimensions, which sets $\epsilon_1 = 0.010 \pm 0.002$ in Chapter~\ref{ch:foundation}'s equation of state and determines $c_s^2 \sim 1/\epsilon_1 \sim 100$ in Chapter~\ref{ch:sound-speed}~\cite{ch4:PaperV}; (ii)~the derivation of $\xi = 1/6$ as a structural prediction (three independent derivations in Section~\ref{sec:4-xi-derivation}), confirming the parameter used throughout Chapters~\ref{ch:foundation}--\ref{ch:nogo}; and (iii)~the Chern--Simons topological terms and their coupling channels, which may inform future engineering applications of the framework.


%% ============================================================
\section{The Spectral Triple on the Warped Orbifold}
\label{sec:4-spectral-triple}

\subsection{Geometry and Conventions}
\label{sec:4-geometry}

The background geometry is the warped product
\begin{equation}
  ds^2 = e^{2A(y)}\, \tilde{g}_{\mu\nu}(x)\, dx^\mu\, dx^\nu + dy^2,
  \label{eq:4-2}
\end{equation}
where $y \in [0, y_c]$ is the extra-dimensional coordinate with $\mathbb{Z}_2$ identification $y \to -y$, and $A(y)$ is the warp factor. On the Randall--Sundrum background, $A(y) = -ky$ with $k$ the $\mathrm{AdS}_5$ curvature scale~\cite{ch4:RS1999a}. The UV brane sits at $y = 0$, the IR brane at $y = y_c$. The hierarchy is generated by the warp factor:
\begin{equation}
  e^{-ky_c} = \frac{m_W}{M_{\mathrm{Pl}}} \sim 10^{-16},
  \label{eq:4-3}
\end{equation}
fixing $ky_c = \ln(M_{\mathrm{Pl}}/m_W) = 39.56$ for $k \sim 10^8\;\mathrm{GeV}$.  In this chapter, three values of $ky_c$ appear, corresponding to different physical contexts: $ky_c = 39.56$ is the exact hierarchy relation; $ky_c \approx 37$ is used in the fermion mass analysis (following the RS flavour literature convention with $k = 10^{11}\;\mathrm{GeV}$); $ky_c = 35$ is the standard RS1 collider benchmark ($k = M_{\mathrm{Pl}}$).  All three are self-consistent within their respective parameter choices; the physical predictions ($w_0(\zeta_0)$, mass ratios, KK spectra) are evaluated at the value appropriate to each context.

The f\"unfbein (5D vielbein) is:
\begin{subequations}\label{eq:4-4}
\begin{align}
  e^a_{\ \mu} &= e^{A(y)}\, \tilde{e}^a_{\ \mu}(x), \quad a = 0,\ldots,3, \label{eq:4-4a}\\
  e^5_{\ y} &= 1. \label{eq:4-4b}
\end{align}
\end{subequations}

The torsion-free 5D spin connection has components:
\begin{subequations}\label{eq:4-5}
\begin{align}
  \omega^{ab}_{\ \ \mu} &= \tilde{\omega}^{ab}_{\ \ \mu}(x), \label{eq:4-5a}\\
  \omega^{a5}_{\ \ \mu} &= A'(y)\, e^{A(y)}\, \tilde{e}^a_{\ \mu}(x), \label{eq:4-5b}\\
  \omega^{ab}_{\ \ 5} &= \omega^{a5}_{\ \ 5} = 0. \label{eq:4-5c}
\end{align}
\end{subequations}

The component~\eqref{eq:4-5b} is the \textbf{spin-warp coupling}---the gravitational analog of a gauge connection mixing 4D and 5D spinor components.


\subsection{The Dirac Operator}
\label{sec:4-dirac}

The 5D Dirac operator on the warped orbifold is:
\begin{equation}
  D_5 = e^{-A(y)}\, \widetilde{D}_4 + \gamma^5\!\left(\partial_y + 2A'(y)\right),
  \label{eq:4-6}
\end{equation}
where
\begin{equation}
  \widetilde{D}_4 = \gamma^a\, \tilde{e}_a^{\ \mu}\!\left(\partial_\mu + \tfrac{1}{4}\, \tilde{\omega}_\mu^{\ bc}\, \gamma_{bc}\right)
  \label{eq:4-7}
\end{equation}
is the 4D Dirac operator on $(M_4, \tilde{g}_{\mu\nu})$, and the factor $2A'(y)$ in the fifth-dimensional piece arises from the spin-warp coupling~\eqref{eq:4-5b}. The coefficient~2 equals $d/2$ where $d = 4$ is the number of brane dimensions (the brane has codimension~1 in 5D)---it ensures compatibility of $D_5$ with the warped metric~\cite{ch4:GrossmanNeubert2000}.

The square of the Dirac operator is a Lichnerowicz-type formula:
\begin{equation}
  D_5^2 = e^{-2A}\, \widetilde{D}_4^2 - A'\, e^{-A}\, \gamma^5\, \widetilde{D}_4
         + \partial_y^2 + 4A'\, \partial_y + \bigl(2A'' + 4A'^2\bigr).
  \label{eq:4-8}
\end{equation}

On the RS background ($A = -ky$), the $y$-dependent terms become:
\begin{equation}
  \partial_y^2 - 4k\, \partial_y + 4k^2
  = (\partial_y - 2k)^2,
  \label{eq:4-9}
\end{equation}
which is a perfect-square Sturm--Liouville operator.


\subsection{The Spectral Triple}
\label{sec:4-triple}

\begin{definition}[Bulk spectral triple]\label{def:4-1}
The bulk spectral triple for the warped orbifold is the real spectral triple $(\mathcal{A}, \mathcal{H}, D_5)$ with:

\textbf{Algebra:}
\begin{equation}
  \mathcal{A} = C^\infty\!\bigl(M_4 \times [0, y_c]\bigr)^{\mathbb{Z}_2},
  \label{eq:4-10}
\end{equation}
the algebra of smooth, $\mathbb{Z}_2$-invariant functions on the product manifold.

\textbf{Hilbert space:}
\begin{equation}
  \mathcal{H} = L^2\!\bigl(M_4 \times [0, y_c],\, S_5,\, \sqrt{G}\, d^4x\, dy\bigr),
  \label{eq:4-11}
\end{equation}
the space of square-integrable 5D spinor fields with the warp-weighted measure $\sqrt{G} = e^{4A(y)}\sqrt{\tilde{g}}$.

\textbf{Dirac operator:} $D_5$ as defined in equation~\eqref{eq:4-6}.

\textbf{Real structure:}
\begin{equation}
  J = C \cdot K,
  \label{eq:4-12}
\end{equation}
where $C$ is charge conjugation and $K$ is complex conjugation.
\end{definition}

The spectral triple has \textbf{KO-dimension~5} (odd). This has an immediate consequence:

\begin{proposition}[Odd-Dimension Obstruction]\label{prop:4-1}
There is no grading operator $\gamma$ on $(\mathcal{A}, \mathcal{H}, D_5)$ satisfying $\gamma^2 = 1$, $\gamma D_5 = -D_5 \gamma$, and $\gamma J = J\gamma$ simultaneously.
\end{proposition}

\begin{proof}
In KO-dimension~5, the sign table gives $J^2 = 1$, $JD = DJ$, and $J\gamma = -\gamma J$~\cite{ch4:Connes1994}. But the physical chirality $\gamma^5$ satisfies $\gamma^5 D_5 = -D_5 \gamma^5$ only for massless 5D fields; the spin-warp coupling $2A'(y)$ breaks this. Moreover, the $\mathbb{Z}_2$ projection $\psi(x,y) \to \gamma^5 \psi(x,-y)$ uses $\gamma^5$ as the orbifold involution, which prevents it from serving simultaneously as a grading.
\end{proof}

The obstruction prevents forming the standard product $D = D_5 \otimes \mathbf{1}_F + \gamma \otimes D_F$ with a finite internal space~$F$, which is the construction used in the flat $M_4 \times F$ case~\cite{ch4:CC1997}. We resolve this in the next subsection.


\subsection{The Layered Spectral Architecture}
\label{sec:4-layered}

\paragraph{Resolution.} We construct the total geometry as two coupled spectral triples:

\textbf{Layer~1 (Bulk):} $(\mathcal{A}_{\mathrm{bulk}}, \mathcal{H}_{\mathrm{bulk}}, D_5)$
\begin{itemize}
  \item KO-dimension~5 (odd)
  \item No grading
  \item Content: gravity, cosmological constant, cuscuton dark energy, moduli, topological terms
  \item Domain: $M_4 \times [0, y_c]$
\end{itemize}

\textbf{Layer~2 (Brane):} $(\mathcal{A}_{\mathrm{brane}}, \mathcal{H}_{\mathrm{brane}}, D_{B \times F})$
\begin{itemize}
  \item KO-dimension~4 (even)
  \item Grading: $\gamma_5$ exists as the chiral operator in 4D
  \item Content: Standard Model gauge fields, Higgs doublet, fermion spectrum
  \item Domain: brane at $y = y_c$ (IR brane), with $D_{B \times F} = D_{M_4} \otimes \mathbf{1}_F + \gamma_5 \otimes D_F$
\end{itemize}

\textbf{Coupling:} Israel junction conditions (Chapter~\ref{ch:foundation}, Section~IV.D) at $y = y_c$ relate the bulk geometry (Layer~1) to the brane stress-energy (Layer~2).

\begin{theorem}[Axiom Preservation under Junction Coupling]\label{thm:4-axiom-preservation}
Coupling two spectral triples through Israel junction conditions preserves all NCG axioms (first-order condition, orientation, Poincar\'e duality).
\end{theorem}

\begin{proof}
The junction conditions couple the two layers through the background geometry (metric coefficients), not through the algebraic structures of the spectral triples. The proof proceeds in three steps.

\paragraph{(i) Self-adjoint extension.}
The orbifold boundary conditions ($\mathbb{Z}_2$ chiral projection) define elliptic boundary conditions for the 5D Dirac operator $D_5$. By the B\"ar--Ballmann theorem~\cite{ch4:bar2012}, $D_5^{\mathrm{orb}}$ with these boundary conditions is self-adjoint with compact resolvent on compact $M_4$. The Israel junction conditions fix the \emph{coefficients} of $D_5$ (through $A'(y)$ at the boundaries) but do not modify the \emph{domain}---the spinor boundary conditions are determined by the orbifold structure, not the junction conditions.

\paragraph{(ii) First-order condition and orientation.}
For Layer~1: $\mathcal{A}_{\mathrm{bulk}}$ is commutative, so the first-order condition $[[D,a],Jb^*J^{-1}] = 0$ holds automatically---$[D_5, a]$ is a zeroth-order (multiplication) operator, and multiplication operators commute. This is purely algebraic and independent of boundary conditions. For Layer~2: the junction conditions rescale $D_B = e^{-A(y_c)} \widetilde{D}_4$ by a positive constant, which factors out of the commutator. The orientation axiom depends on the Hochschild cycle $\pi_D(c)$, which is a local interior computation using the symbol of $D_5$---the boundary values of $A'$ do not enter.

\paragraph{(iii) Poincar\'e duality.}
Since $I = [0, y_c]$ is contractible, $K_*(C(M_4 \times I)^{\mathbb{Z}_2}) \cong K_*(C(M_4))$, and the intersection form inherits non-degeneracy from the doubled manifold $M_4 \times S^1$ via the orbifold-to-double correspondence~\cite{ch4:bruning2001,ch4:bar2012}. The key step is verifying this correspondence for the warped metric $\mathrm{d}s^2 = e^{-2A(y)}\,g_{\mu\nu}\,\mathrm{d}x^\mu\,\mathrm{d}x^\nu + \mathrm{d}y^2$. Consider the one-parameter family $A_t(y) = t\,A(y)$ for $t \in [0,1]$. Each $A_t$ defines a Dirac operator $D_5(t)$ on the orbifold. The principal symbol of $D_5(t)$ is that of a Dirac-type operator for every~$t$---warping modifies the vielbein and spin connection but not the symbol class---so $D_5(t)$ is elliptic. Combined with the $\mathbb{Z}_2$ boundary conditions, the Fredholm property holds throughout the deformation. The intersection form $\langle [e], [f] \rangle = \mathrm{index}(e\,D_+(t)\,f)$ is therefore continuous and integer-valued in~$t$, hence constant. At $t = 0$ (flat metric), the standard orbifold-to-double correspondence~\cite{ch4:bruning2001} gives non-degeneracy. Therefore non-degeneracy holds at $t = 1$ (the physical warped metric). For Layer~2, rescaling $D_B$ by the positive constant $e^{-A(y_c)}$ is likewise a homotopy preserving the Fredholm property, so the brane intersection form is unchanged.
\end{proof}

\begin{remark}[Computational verification]
\label{rem:4-conj-verification}
The algebraic content of the proof has been verified computationally for the SM finite spectral triple (KO-dimension~6). All axioms involving~$D$ are homogeneous: $[[cD, a], Jb^*J^{-1}] = c\,[[D, a], Jb^*J^{-1}]$, $\{cD, \gamma\} = c\,\{D, \gamma\}$, $J(cD) = c\,JD$. All axioms not involving~$D$ ($J^2 = \epsilon$, $J\gamma = \epsilon''\gamma J$) are unchanged by the rescaling. The Fredholm index is constant for all $c > 0$ (verified: $\mathrm{index}(cD_+) = \mathrm{index}(D_+) = 0$ for the 1-generation spectral triple). The only non-computational ingredient is the Bar--Ballmann compact resolvent theorem~\cite{ch4:bar2012}, which is standard.
\end{remark}

\begin{remark}
The odd-dimension obstruction is not a deficiency but a \emph{feature}: the NCG framework correctly recognizes that the 5D bulk should carry only gravitational content (no gauge fields), while gauge fields are naturally brane-localized. The layered architecture \emph{enforces} the physical structure of the Randall--Sundrum model.
\end{remark}


\subsection{\texorpdfstring{$\mathbb{Z}_2$}{Z2} Action on Spinors and the Chiral Spectrum}
\label{sec:4-z2-spinors}

The $\mathbb{Z}_2$ orbifold acts on 5D spinors as:
\begin{equation}
  \psi(x, y) \to \gamma^5\, \psi(x, -y).
  \label{eq:4-13}
\end{equation}

This projects the KK spectrum to a single chirality per generation. Left-handed zero modes with bulk mass parameter $c > 1/2$ are localized near the UV brane, while $c < 1/2$ modes localize near the IR brane~\cite{ch4:GherghettaPomarol2000}:
\begin{subequations}\label{eq:4-14}
\begin{align}
  f_0^L(y) &= N_L\, e^{(2-c)ky}, \label{eq:4-14a}\\
  f_0^R(y) &= N_R\, e^{(2+c)ky}. \label{eq:4-14b}
\end{align}
\end{subequations}

This is the Randall--Sundrum flavor hierarchy mechanism: the overlap of fermion zero-mode profiles with the IR-brane Higgs determines the Yukawa couplings, generating the observed mass hierarchies from $\mathcal{O}(1)$ bulk parameters.


%% ============================================================
\section{The Seeley--DeWitt Expansion}
\label{sec:4-seeley-dewitt}

\subsection{Heat Kernel and Spectral Action}
\label{sec:4-heat-kernel}

The spectral action~\eqref{eq:4-1} is evaluated through the heat kernel expansion. Define the heat kernel:
\begin{equation}
  K(t) = \mathrm{Tr}\!\left(e^{-t D_5^2}\right).
  \label{eq:4-15}
\end{equation}

The Seeley--DeWitt asymptotic expansion~\cite{ch4:Gilkey1984,ch4:Gilkey1975} gives:
\begin{equation}
  K(t) \sim \sum_{n=0}^{\infty} a_n(D_5^2)\, t^{(n-5)/2}
  \quad \text{as } t \to 0^+.
  \label{eq:4-16}
\end{equation}

In five dimensions, both integer coefficients $a_n$ (bulk integrals over $M_4 \times I$) and half-integer coefficients $a_{n+1/2}$ (boundary integrals over the branes) contribute. The spectral action becomes:
\begin{equation}
  S = \sum_n f_n\, \Lambda^{5-n}\, a_n(D_5^2),
  \label{eq:4-17}
\end{equation}
where $f_n = \int_0^{\infty} f(u)\, u^{(5-n-2)/2}\, du$ are the momenta of the cutoff function~$f$.


\subsection{Bulk Coefficients}
\label{sec:4-bulk-coefficients}

We compute each coefficient on the RS background ($A = -ky$), where the 5D curvature invariants are:
\begin{subequations}\label{eq:4-18}
\begin{align}
  R_5 &= -20k^2, \label{eq:4-18a}\\
  R_{MN}\, R^{MN} &= 80k^4, \label{eq:4-18b}\\
  R_{MNPQ}\, R^{MNPQ} &= 40k^4. \label{eq:4-18c}
\end{align}
\end{subequations}

\paragraph{Coefficient $a_0$ (Cosmological constant).}
\begin{multline}
  a_0 = (4\pi)^{-5/2} \int_{M_4 \times I} \mathrm{tr}(\mathbf{1})\, \sqrt{G}\, d^5x \\
      = (4\pi)^{-5/2}\, d_S \int_{M_4} \sqrt{\tilde{g}}\, d^4x \int_0^{y_c} e^{4A(y)}\, dy,
  \label{eq:4-19}
\end{multline}
where $d_S = \mathrm{tr}(\mathbf{1}) = 4$ is the 5D spinor fiber dimension. The $y$-integral gives $(e^{4ky_c} - 1)/(4k)$. In the spectral action, $f_0\, \Lambda^5\, a_0$ produces a 5D cosmological constant, which is absorbed by the sequestering mechanism (Chapter~\ref{ch:foundation}, Section~III).

\paragraph{Coefficient $a_1$ (Subleading cosmological constant).}
\begin{equation}
  a_1 = (4\pi)^{-5/2} \int_{M_4 \times I} d_S\, (6E + R_5)\, \sqrt{G}\, d^5x,
  \label{eq:4-20}
\end{equation}
where $E$ is the endomorphism of the Dirac Laplacian. On $\mathrm{AdS}_5$: $E = -R_5/4 = 5k^2$, giving
\begin{equation}
  d_S(6E + R_5) = 4(30k^2 - 20k^2) = 40k^2.
  \label{eq:4-21}
\end{equation}
This contributes a subleading cosmological constant proportional to $f_1\, \Lambda^4\, k$.

\paragraph{Coefficient $a_2$ (Einstein--Hilbert).}
The coefficient $a_2$ contains the scalar curvature $R_4$ of the 4D metric. Matching to the standard form $\frac{M_{\mathrm{Pl}}^2}{2} \int \sqrt{g}\, R_4$:
\begin{equation}
  M_{\mathrm{Pl}}^2 = \frac{f_2\, \Lambda^3}{3k\, (4\pi)^{5/2}}.
  \label{eq:4-22}
\end{equation}
This \textbf{determines the spectral cutoff}. With $k \sim 10^8\;\mathrm{GeV}$ and $M_{\mathrm{Pl}} \sim 1.22 \times 10^{19}\;\mathrm{GeV}$:
\begin{equation}
  \Lambda \sim \left(\frac{3k\, (4\pi)^{5/2}\, M_{\mathrm{Pl}}^2}{f_2}\right)^{1/3} \sim 10^{17}\;\mathrm{GeV}.
  \label{eq:4-23}
\end{equation}
The cutoff is at the GUT scale---above the AdS curvature scale~$k$ but below $M_{\mathrm{Pl}}$. The spectral action expansion is valid for energies below~$\Lambda$, which is self-consistent.

\paragraph{Coefficient $a_3$ (Gauss--Bonnet).}
The third Seeley--DeWitt coefficient for a Laplace-type operator $P = -\nabla^2 + E$ on a $d$-dimensional manifold contains the quadratic curvature invariants (Vassilevich~\cite{ch4:Vassilevich2003}, Eq.~4.3; Gilkey~\cite{ch4:Gilkey1984}):
\begin{multline}
  a_3(x, P) = \tfrac{1}{360}\, \mathrm{tr}\Bigl[
    2\, R_{MNPQ}^2\, \mathbf{1} - 2\, R_{MN}^2\, \mathbf{1} \\
    {}+ 5\, R^2\, \mathbf{1}
    + 60\, R\, E - 180\, E^2 \\
    {}+ 30\, \Omega_{MN}^2 + (\nabla\text{-terms})
  \Bigr].
  \label{eq:4-24a}
\end{multline}
For $D_5^2$ on the spin bundle (Lichnerowicz formula: $E = -R/4 \cdot \mathbf{1}_{N_s}$, $\Omega_{MN} = \tfrac{1}{4}R_{MN}^{\ \ ab}\gamma_a\gamma_b$, with $N_s = 2^{[d/2]} = 4$ in $d = 5$), the explicit traces are:
\begin{equation}
  \mathrm{tr}(\mathbf{1}) = 4, \quad
  \mathrm{tr}(E) = -R, \quad
  \mathrm{tr}(E^2) = \tfrac{1}{4}\, R^2, \quad
  \mathrm{tr}(\Omega^2) = -\tfrac{1}{2}\, R_{MNPQ}^2.
  \label{eq:4-24b}
\end{equation}
Substituting into~\eqref{eq:4-24a}, the curvature-squared content is:
\begin{multline}
  a_3 \supset \tfrac{1}{360}\bigl[-7\, R_{MNPQ}^2 \\
  {}- 8\, R_{MN}^2 - 85\, R^2\bigr]\,.
  \label{eq:4-24c}
\end{multline}
Converting to the $(C^2, \mathcal{E}_d, R^2)$ basis requires the \emph{dimension-dependent} inverse Weyl decomposition. In $d$ dimensions:
\begin{align}
  C^2 &= R_{MNPQ}^2 - \tfrac{4}{d-2}\,R_{MN}^2 + \tfrac{2}{(d-1)(d-2)}\,R^2, \notag\\
  \mathcal{E}_d &= R_{MNPQ}^2 - 4\,R_{MN}^2 + R^2. \label{eq:4-weyl-decomp}
\end{align}
Inverting for $d = 5$:
\begin{equation}
  R_{MNPQ}^2 = \tfrac{3}{2}\, C^2 - \tfrac{1}{2}\, \mathcal{E}_5 + \tfrac{1}{4}\, R^2, \qquad
  R_{MN}^2 = \tfrac{3}{8}(C^2 - \mathcal{E}_5) + \tfrac{5}{16}\, R^2.
  \label{eq:4-24d-inv}
\end{equation}
Substituting into~\eqref{eq:4-24c}:
\begin{equation}
  a_3 \supset \frac{1}{360}\Bigl[-\tfrac{27}{2}\, C^2 + \tfrac{13}{2}\, \mathcal{E}_5 - \tfrac{357}{4}\, R^2\Bigr],
  \label{eq:4-24d}
\end{equation}
giving the \emph{5D bulk} spectral ratios $C^2 : \mathcal{E}_5 : R^2 = -\tfrac{27}{2} : \tfrac{13}{2} : -\tfrac{357}{4}$ for the $a_3$ (mass-dimension~2) coefficient in five dimensions. Note: in $d = 4$ the inverse gives $R_{MNPQ}^2 = 2C^2 - \mathcal{E}_4 + \tfrac{1}{3}R^2$, yielding ratios $-18 : 11 : -90$; na\"{\i}vely applying the $d = 4$ identities to the 5D bulk overestimates $\mathcal{E}_5$ by a factor $11/(13/2) = 22/13 \approx 1.69$ and fails the $\mathrm{AdS}_5$ consistency check (predicting $R_{MNPQ}^2 = \tfrac{40}{3}\,k^4$ instead of the correct $40k^4$). The $a_3$ coefficient controls the Gauss--Bonnet coupling in the bulk; the $a_4$ coefficient controls the curvature-squared terms in the 4D effective action. The distinction is critical for the NCG--asymptotic safety bridge.

The 5D Gauss--Bonnet invariant is:
\begin{equation}
  \mathcal{E}_5 = R_5^2 - 4\, R_{MN}^2 + R_{MNPQ}^2.
  \label{eq:4-24}
\end{equation}
In five dimensions, $\mathcal{E}_5$ is \emph{dynamical}---unlike in four dimensions, where the dimensionally dependent identity~\cite{ch4:Lovelock1970} reduces it to the topological Euler density and it contributes no equations of motion~\cite{ch4:Lovelock1971}. On the $\mathrm{AdS}_5$ background:
\begin{equation}
  \mathcal{E}_5 = (-20k^2)^2 - 4(80k^4) + 40k^4
               = 400k^4 - 320k^4 + 40k^4 = 120k^4.
  \label{eq:4-25}
\end{equation}

The spectral action contribution from $a_3$ is $S \supset f_3\, \Lambda^2\, (4\pi)^{-5/2}\, a_3$. The GB part extracts the $\mathcal{E}_5$ term with its coefficient~$+13/(2 \cdot 360)$ from~\eqref{eq:4-24d}. After absorbing factors, the effective GB coupling is:
\begin{equation}
  \alpha_{\mathrm{GB}} = \frac{f_3\, \Lambda^2}{(4\pi)^{5/2} \cdot 360},
  \label{eq:4-26}
\end{equation}
from which the dimensionless parameter is:
\begin{equation}
  \hat{\alpha} = \frac{\alpha_{\mathrm{GB}}\, k^2}{M_5^3} \sim 10^{-2}.
  \label{eq:4-27}
\end{equation}

This is the central result of Section~\ref{sec:4-seeley-dewitt}: the NCG spectral action produces a Gauss--Bonnet coupling at approximately 1\% of the Einstein--Hilbert term. The derivation chain from the Seeley--DeWitt coefficient through the trace formulas~\eqref{eq:4-24b} to the numerical coupling~\eqref{eq:4-27} is fully reproducible; the computation is implemented in the supplementary code (see Appendix~\ref{app:c1-gb-kk}). As we prove in Section~\ref{sec:4-gb-universality}, this value is universal---independent of the compactification dimension and topology.


\subsection{Boundary Coefficients}
\label{sec:4-boundary-coefficients}

The orbifold fixed points contribute boundary (brane-localized) Seeley--DeWitt coefficients.

\paragraph{Coefficient $a_{1/2}$ (Vanishing).} For the $\mathbb{Z}_2$ orbifold, the boundary contributions from Dirichlet and Neumann components cancel:
\begin{equation}
  a_{1/2} = 0.
  \label{eq:4-28}
\end{equation}

\paragraph{Coefficient $a_{3/2}$ (Brane tension).} Produces corrections to the brane tensions $\sigma_{\mathrm{UV}}$ and $\sigma_{\mathrm{IR}}$ that are absorbed into the Israel junction conditions.

\paragraph{Coefficient $a_{5/2}$ (Brane-localized Einstein--Hilbert).}
The $a_{5/2}$ coefficient produces a DGP-type~\cite{ch4:DGP2000} induced gravity term on each brane:
\begin{equation}
  S_{\mathrm{DGP}} = r_c \int \sqrt{-h}\, R_{\mathrm{brane}}\, d^4x,
  \label{eq:4-29}
\end{equation}
with crossover scale
\begin{equation}
  L_c = \frac{M_{\mathrm{Pl}}^2}{r_c\, M_5^3} \sim 10^{-6}\;\mathrm{m}.
  \label{eq:4-30}
\end{equation}
This micron-scale crossover is potentially testable in short-distance gravity experiments (E\"ot-Wash type)~\cite{ch4:Kapner2007}. The induced gravity term does not modify cosmological dynamics ($L_c \ll H_0^{-1}$) but affects local gravitational physics at sub-millimeter scales.

\paragraph{Coefficient $a_{7/2}$ (Chern--Simons).}
The half-integer boundary coefficient produces topological terms on the branes. The full content of this coefficient is derived in Section~\ref{sec:4-chern-simons}.


\subsection{Physical Content Summary}
\label{sec:4-physical-summary}

\begin{table*}[htbp]
\centering
\caption{Complete Seeley--DeWitt content of the spectral action on the warped $S^1/\mathbb{Z}_2$ orbifold.}
\label{tab:4-1}
\begin{tabular}{@{}llll@{}}
\toprule
Coefficient & $\Lambda$ power & Physical content & Numerical value \\
\midrule
$a_0$       & $\Lambda^5$     & 5D cosmological constant       & Sequestered \\
$a_1$       & $\Lambda^4$     & Subleading CC                  & Sequestered \\
$a_2$       & $\Lambda^3$     & Einstein--Hilbert ($M_{\mathrm{Pl}}$)  & Determines $\Lambda = 10^{17}\;\mathrm{GeV}$ \\
$a_3$       & $\Lambda^2$     & Gauss--Bonnet ($\hat\alpha$)    & $\sim 10^{-2}$ \\
$a_{1/2}$   & $\Lambda^{9/2}$ & Vanishes for $\mathbb{Z}_2$    & 0 \\
$a_{3/2}$   & $\Lambda^{7/2}$ & Brane tension shift            & Absorbed \\
$a_{5/2}$   & $\Lambda^{5/2}$ & DGP induced gravity            & $L_c \sim 10^{-6}\;\mathrm{m}$ \\
$a_{7/2}$   & $\Lambda^{3/2}$ & Chern--Simons terms            & Sections~\ref{sec:4-chern-simons}--\ref{sec:4-em-gravity} \\
\bottomrule
\end{tabular}
\end{table*}


%% ============================================================
\section{Gauss--Bonnet Universality}
\label{sec:4-gb-universality}

\subsection{The Gilkey Coefficient Structure}
\label{sec:4-gilkey}

The Gauss--Bonnet coupling $\hat{\alpha} \sim 10^{-2}$ arises from the $a_3$ coefficient of the heat kernel expansion. A natural question is whether this value depends on the compactification details---the number of extra dimensions, the topology, or the specific warp factor profile.

\begin{scalingrelation}[Approximately Universal Gauss--Bonnet Coupling]\label{sr:4-1}
The dimensionless Gauss--Bonnet coupling $\hat{\alpha} = \alpha_{\mathrm{GB}}\, k^2 / M_5^3$ satisfies $\hat{\alpha} \sim 10^{-2}$ for any warped compactification of the form $M_4 \times K$ with Dirac spinors, varying by at most a factor of 2--3 across dimensions $d = 5$ to $d = 10$.  This result is approximately unconditional (subject to the cutoff function $f_n$ ratio being $\mathcal{O}(1)$).
\end{scalingrelation}

\begin{proof}
The $a_3$ Seeley--DeWitt coefficient is determined by the Gilkey universal coefficients~\cite{ch4:Gilkey1984,ch4:Gilkey1973}, which are computed from the local curvature invariants:
\begin{multline}
  a_3 = (4\pi)^{-d/2} \!\int_M \mathrm{tr}\bigl[\alpha_1 R^2 + \alpha_2 R_{MN}^2 \\
  {}+ \alpha_3 R_{MNPQ}^2 + \alpha_4 E^2 \\
  {}+ \alpha_5 R\, E + \cdots\bigr] \sqrt{G}\, d^d x,
  \label{eq:4-31}
\end{multline}
where $\alpha_i$ are universal rational numbers depending only on the spin representation. For Dirac spinors in $d$ dimensions, the trace factor $\mathrm{tr}(\mathbf{1}) = 2^{[d/2]}$: this is~4 in $d = 5$, 8 in $d = 6$, 16 in $d = 8$, and 32 in $d = 10$.

The Gauss--Bonnet combination $\mathcal{E}_d = R^2 - 4\, R_{MN}^2 + R_{MNPQ}^2$ extracts a specific linear combination of the $\alpha_i$. On a maximally symmetric background ($\mathrm{AdS}_d$), the curvature invariants satisfy universal proportionality relations:
\begin{subequations}\label{eq:4-32}
\begin{align}
  \frac{R_{MNPQ}^2}{R^2} &= \frac{2}{d(d-1)}, \label{eq:4-32a}\\
  \frac{R_{MN}^2}{R^2}   &= \frac{1}{d}. \label{eq:4-32b}
\end{align}
\end{subequations}

Therefore:
\begin{align}
  \frac{\mathcal{E}_d}{R^2}
  &= 1 - \frac{4}{d} + \frac{2}{d(d{-}1)} \notag\\
  &= \frac{(d{-}2)(d{-}3)}{d(d{-}1)}\,.
  \label{eq:4-33}
\end{align}

The dimensionless coupling scales as:
\begin{equation}
  \hat{\alpha} \sim \frac{f_3}{f_2} \cdot \frac{\Lambda^2}{M_{\mathrm{Pl}}^2}
               \cdot \frac{\mathrm{tr}(\mathbf{1})}{360}
               \cdot \frac{(d-2)(d-3)}{d(d-1)},
  \label{eq:4-34}
\end{equation}
where the ratio $f_3/f_2 \sim \mathcal{O}(1)$ depends on the cutoff function, $\Lambda^2/M_{\mathrm{Pl}}^2 \sim (k/M_{\mathrm{Pl}})^{2/3}$ from \eqref{eq:4-22}--\eqref{eq:4-23}, and the remaining factors are $\mathcal{O}(1)$ dimensionless numbers.

For $d = 5$: $\mathrm{tr}(\mathbf{1}) = 4$, geometric factor $= (3 \cdot 2)/(5 \cdot 4) = 3/10$.

For $d = 6$: $\mathrm{tr}(\mathbf{1}) = 8$, geometric factor $= (4 \cdot 3)/(6 \cdot 5) = 2/5$.

The ratio $(d = 6)/(d = 5)$ of the $d_S \times \text{geometric factor}$ is $(8 \cdot 2/5)/(4 \cdot 3/10) = 8/3$. But the KK volume integral changes as $k^{-(d-4)}$, and the relation~\eqref{eq:4-22} between $\Lambda$ and $M_{\mathrm{Pl}}$ generalizes to $M_{\mathrm{Pl}}^2 \propto f_2\, \Lambda^{d-2}/k^{d-4}$. The key observation is that the combination $\hat{\alpha}$ is determined by the \emph{ratio} $f_3/f_2$ of cutoff function moments (which is $\mathcal{O}(1)$ and independent of~$d$), multiplied by the universal factor~\cite{ch4:EguchiGilkeyHanson1980}:
\begin{equation}
  \hat{\alpha} = \frac{f_3}{f_2} \cdot d_S(d) \cdot
                 \frac{(d-2)(d-3)}{d(d-1)} \Big/ \bigl[360 \cdot d_S(d)\bigr]
               = \frac{f_3}{f_2} \cdot \frac{(d-2)(d-3)}{360\, d(d-1)}.
  \label{eq:4-35}
\end{equation}

For $d = 5$: $\hat{\alpha} = (f_3/f_2) \cdot 6/6000 = (f_3/f_2) \cdot 10^{-3}$. With $f_3/f_2 \sim \mathcal{O}(10)$, this gives $\hat{\alpha} \sim 10^{-2}$.

For $d = 6$: $\hat{\alpha} = (f_3/f_2) \cdot 12/10800 = (f_3/f_2) \cdot 1.1 \times 10^{-3}$. Same order.

For $d = 10$: $\hat{\alpha} = (f_3/f_2) \cdot 56/32400 = (f_3/f_2) \cdot 1.7 \times 10^{-3}$. Same order.

The variation across $d = 5$ to $d = 10$ is less than a factor of~2 in the geometric factor, confirming \emph{qualitative} universality ($\hat{\alpha} \sim \mathcal{O}(10^{-2})$ across $d = 5$--$10$) to within the stated precision. This universality is established for maximally symmetric backgrounds; on the warped RS background, boundary heat-kernel contributions introduce corrections that are absorbed into $b_{3/2}$ (see Section~\ref{sec:4-brane-parameter}).
\end{proof}

\begin{remark}
We have verified this explicitly in six dimensions by computing the spectral action on $M_4 \times T^2/\mathbb{Z}_2$. The cubic Lovelock invariant (which first appears in $d = 6$) contributes at $\mathcal{O}(10^{-4})$, suppressed by an additional factor of $(k/\Lambda)^2 \sim 10^{-18}$ relative to the Gauss--Bonnet term. The Gauss--Bonnet coupling itself remains at $\hat{\alpha} \sim 10^{-2}$, confirming qualitative universality at the order-of-magnitude level.
\end{remark}


\subsection{From Gauss--Bonnet Coupling to Kinetic Correction: The Complete $\hat{\alpha} \to \epsilon_1$ Derivation}
\label{sec:4-epsilon1}

This subsection provides the complete derivation of the kinetic correction~$\epsilon_1$ from the spectral-action Gauss--Bonnet coupling~$\hat{\alpha}$. The logical chain is:
\[
  \underbrace{\text{Spectral action}}_{\S\ref{sec:4-seeley-dewitt}}
  \;\xrightarrow{\;a_3\;}
  \underbrace{\alpha_{\mathrm{GB}}}_{\text{Eq.~\eqref{eq:4-26}}}
  \;\xrightarrow{\;k^2/M_5^3\;}
  \underbrace{\hat{\alpha} \sim 10^{-2}}_{\S\ref{sec:4-gb-universality}}
\]
\vspace{-0.5em}
\[
  \phantom{\text{Spectral action}}
  \;\xrightarrow{\;\times\, C_{\mathrm{GB}}\;}
  \underbrace{\epsilon_1 = 0.010 \pm 0.002}_{\text{This section}}
  \;\xrightarrow{\;\text{Chapter~\ref{ch:foundation}}\;}
  w_0(\zeta_0).
\]


\subsubsection{Step 1: The 5D Bulk Action with Gauss--Bonnet}

The complete 5D bulk action, incorporating the GB correction from the $a_3$ Seeley--DeWitt coefficient (Section~\ref{sec:4-seeley-dewitt}), is
\begin{multline}
  S_{\mathrm{bulk}} = \int d^5x\, \sqrt{G}\, \Bigl[
    \tfrac{1}{2}\bigl(M_5^3 - \xi\phi^2\bigr) R_5
    + \Lambda_5 \\
    + \alpha_{\mathrm{GB}}\, \mathcal{E}_5
    + P(X_5, \phi)
  \Bigr]
  + S_{\mathrm{brane}},
  \label{eq:4-36}
\end{multline}
where $X_5 \equiv -\tfrac{1}{2}\, G^{MN}\, \partial_M\phi\, \partial_N\phi$, $P(X_5) = \mu^2\sqrt{2X_5}$ is the cuscuton kinetic function (Chapter~\ref{ch:foundation}), $\mathcal{E}_5 = R_5^2 - 4\,R_{MN}^2 + R_{MNPQ}^2$ is the 5D GB invariant~\eqref{eq:4-24}, and $\alpha_{\mathrm{GB}}$ is the spectral-action coupling~\eqref{eq:4-26}. In five dimensions, $\mathcal{E}_5$ is \emph{dynamical}---it contributes nontrivially to the field equations, unlike the topological $\mathcal{E}_4$ in four dimensions.


\subsubsection{Step 2: GB-Modified Israel Junction Conditions}

The Gauss--Bonnet term modifies the Israel junction conditions at the $\mathbb{Z}_2$ orbifold branes. Davis~\cite{ch4:Davis2003} derived the generalised junction conditions for a GB brane world. For a $\mathbb{Z}_2$-symmetric brane embedded in a 5D spacetime with GB coupling~$\alpha_{\mathrm{GB}}$, the junction condition takes the form (Davis~\cite{ch4:Davis2003}, Eq.~15):
\begin{equation}
  \bigl[K_{\mu\nu} - K\, h_{\mu\nu}\bigr]
  + 2\alpha_{\mathrm{GB}}\bigl[3J_{\mu\nu} - J\, h_{\mu\nu}
  + 2\, P_{\mu\rho\nu\sigma}\, K^{\rho\sigma}\bigr]
  = -\frac{\kappa_5^2}{2}\, S_{\mu\nu},
  \label{eq:4-davis-jc}
\end{equation}
where $h_{\mu\nu}$ is the induced metric, $K_{\mu\nu}$ is the extrinsic curvature, $J_{\mu\nu} = \tfrac{1}{3}(2K\, K_{\mu\rho}\, K^\rho_{\ \nu} + K_{\rho\sigma}\, K^{\rho\sigma}\, K_{\mu\nu} - 2\, K_{\mu\rho}\, K^{\rho\sigma}\, K_{\sigma\nu} - K^2\, K_{\mu\nu})$ is the third fundamental form, and $P_{\mu\rho\nu\sigma}$ is the divergence-free part of the brane Riemann tensor.

\paragraph{Evaluation on the RS background.}
On the Randall--Sundrum background, the extrinsic curvature at the UV brane ($y = 0$) with $\mathbb{Z}_2$ identification is $K_{\mu\nu} = -k\, h_{\mu\nu}$ (outward normal into the bulk). Then:
\begin{equation}
  K = h^{\mu\nu} K_{\mu\nu} = -4k, \qquad
  K_{\mu\rho}\, K^{\rho}_{\ \nu} = k^2\, h_{\mu\nu}, \qquad
  K_{\rho\sigma}\, K^{\rho\sigma} = 4k^2.
  \label{eq:4-rs-K}
\end{equation}
Computing the third fundamental form:
\begin{align}
  J_{\mu\nu} &= \tfrac{1}{3}\bigl[2(-4k)(k^2) + (4k^2)(-k) - 2(-k^3) - (16k^2)(-k)\bigr]\, h_{\mu\nu} \notag\\
  &= \tfrac{1}{3}\bigl[-8k^3 - 4k^3 + 2k^3 + 16k^3\bigr]\, h_{\mu\nu}
  = 2k^3\, h_{\mu\nu},
  \label{eq:4-Jmunu}
\end{align}
with trace $J = 8k^3$. The $P_{\mu\rho\nu\sigma}$ tensor is built from the \emph{intrinsic} Riemann tensor $\hat{R}_{\mu\rho\nu\sigma}$ of the brane (Davis~\cite{ch4:Davis2003}, Eq.~21). On the static RS background the brane is flat, so $\hat{R}_{\mu\rho\nu\sigma} = 0$ and therefore
\begin{equation}
  P_{\mu\rho\nu\sigma}\, K^{\rho\sigma}\big|_{H=0} = 0.
  \label{eq:4-Ptensor}
\end{equation}
The $P$~tensor activates only on a cosmologically evolving brane ($H \neq 0$); its contribution at $\mathcal{O}(H^2)$ is computed in Step~4 below.

Assembling the GB modification on the static background:
\begin{align}
  &2\alpha_{\mathrm{GB}}\bigl[3J_{\mu\nu} - J\, h_{\mu\nu} + 2\, P_{\mu\rho\nu\sigma}\, K^{\rho\sigma}\bigr]\big|_{H=0}
  \notag\\
  &\quad = 2\alpha_{\mathrm{GB}}\bigl[6k^3 - 8k^3\bigr]\, h_{\mu\nu}
  = -4\, \alpha_{\mathrm{GB}}\, k^3\, h_{\mu\nu}.
  \label{eq:4-gb-assembled}
\end{align}

The standard (non-GB) junction condition gives $3k\, h_{\mu\nu}$. The ratio of the GB correction to the standard term defines the relative modification:
\begin{equation}
  \frac{\text{GB correction}}{\text{standard}}
  = \frac{4\, \alpha_{\mathrm{GB}}\, k^3}{3k}
  = \frac{4}{3}\, \alpha_{\mathrm{GB}}\, k^2.
  \label{eq:4-gb-ratio}
\end{equation}


\subsubsection{Step 3: How GB Breaks the Cuscuton Constraint}

In the absence of the GB term ($\alpha_{\mathrm{GB}} = 0$), the scalar field~$\phi$ satisfies a purely algebraic constraint---the degeneracy condition $P_X + 2X\, P_{XX} = 0$ for $P = \mu^2\sqrt{2X}$ eliminates all time derivatives of~$\phi$, making it non-dynamical (Chapter~\ref{ch:foundation}, Section~III; Chapter~\ref{ch:nogo}, Theorem~1). The effective kinetic energy vanishes exactly: $K_{\mathrm{eff}} = 0$.

The GB term breaks this constraint. The mechanism operates through two channels in the junction condition~\eqref{eq:4-davis-jc}:
\begin{enumerate}
  \item \textbf{Linearised $J$ contribution.} The third fundamental form $J_{\mu\nu}$, cubic in~$K$, couples corrections $\delta K_{\mu\nu} \sim \mathcal{O}(H^2/k)$ to the static RS extrinsic curvature. Linearising $3J_{\mu\nu} - J\, h_{\mu\nu}$ around the RS background produces $\mathcal{O}(H^2)$ corrections to the junction condition.
  \item \textbf{Intrinsic-curvature $P$~contribution.} When cosmological expansion is present ($H \neq 0$), the brane acquires intrinsic curvature $\hat{R}_{\mu\rho\nu\sigma} \sim \mathcal{O}(H^2)$, activating $P_{\mu\rho\nu\sigma}\, K^{\rho\sigma} \neq 0$. This couples the FRW geometry directly to the extrinsic curvature through a channel absent in the standard Israel conditions.
  \item \textbf{KK reduction.} After integration over~$y$, both channels contribute to the effective 4D scalar dynamics as a kinetic correction proportional to~$X$. The $P$~term accounts for $2/3$ of the total $\mathcal{O}(H^2)$ GB source (Step~4).
  \item The modified effective kinetic function becomes:
  \begin{equation}
    P_{\mathrm{eff}}(X) = \mu^2\, \sqrt{2X} + \epsilon_1\, X,
    \label{eq:4-peff}
  \end{equation}
  where $\epsilon_1 \propto \hat{\alpha}$ is the small GB-induced kinetic correction.
\end{enumerate}

This breaks the zero kinetic energy theorem at order~$\epsilon_1$:
\begin{align}
  K_{\mathrm{eff}} &= 2X\, (P_{\mathrm{eff}})_X - P_{\mathrm{eff}}
  \notag\\
  &= 2X\!\left(\frac{\mu^2}{\sqrt{2X}} + \epsilon_1\right) - \mu^2\sqrt{2X} - \epsilon_1 X
  = \epsilon_1\, X \neq 0.
  \label{eq:4-keff-breaking}
\end{align}


\subsubsection{Step 4: Derivation of $C_{\mathrm{GB}} = 2/3$}

The coefficient relating the spectral-action coupling to the kinetic correction is:
\begin{equation}
  \epsilon_1 = \hat{\alpha} \times C_{\mathrm{GB}}.
  \label{eq:4-eps1-def}
\end{equation}

\paragraph{Static fine-tuning.}
On the static RS orbifold, the GB-modified junction condition for the warp factor simplifies to (Davis~\cite{ch4:Davis2003}, Eq.~18):
\begin{equation}
  \bigl[A'\bigr]\!\left(1 + \frac{4}{3}\, \alpha_{\mathrm{GB}}\, (A')^2\right) = -\frac{\sigma}{6\, M_5^3},
  \label{eq:4-warp-jc}
\end{equation}
where $[A'] = -2k$ and the factor $4/3$ arises from $3J_{\mu\nu} - J\, h_{\mu\nu}$ evaluated on the flat brane (cf.~\eqref{eq:4-gb-assembled}; the $P$~tensor does not contribute at $H = 0$).

\paragraph{FRW brane: source decomposition at $\boldsymbol{\mathcal{O}(H^2)}$.}
To derive $C_{\mathrm{GB}}$, we evaluate the full junction condition~\eqref{eq:4-davis-jc} on an FRW brane with extrinsic curvature $K^\mu_{\ \nu} = \mathrm{diag}(-(k+\beta),\, -(k+\gamma),\, -(k+\gamma),\, -(k+\gamma))$, where $\beta,\, \gamma \sim \mathcal{O}(H^2/k)$ encode the cosmological perturbation. At $\mathcal{O}(H^2)$ the GB source $\mathcal{S}^0_{\ 0}$ for the $00$ junction equation receives two contributions:
\begin{align}
  \mathcal{S}^0_{\ 0}\big|_J &= \frac{\partial (3J - J\, h)^0_{\ 0}}{\partial \beta}\bigg|_{\mathrm{RS}}\!\beta_0
  + \frac{\partial (3J - J\, h)^0_{\ 0}}{\partial \gamma}\bigg|_{\mathrm{RS}}\!\gamma_0
  = -3k\, H^2, \label{eq:4-source-J} \\[4pt]
  \mathcal{S}^0_{\ 0}\big|_P &= 2\, P_{\mu\rho\nu\sigma}\, K^{\rho\sigma}\big|_{K_0 = -k\, h}
  = -6k\, H^2, \label{eq:4-source-P}
\end{align}
where $\beta_0 = -(2+3w)H^2/(2k)$ and $\gamma_0 = H^2/(2k)$ are the standard RS extrinsic curvature perturbations, and the $P$~contribution uses the intrinsic FRW Riemann tensor $\hat{R}_{0i0j} = -(\dot{H}+H^2)\,\delta_{ij}$ (which is $\mathcal{O}(H^2)$) contracted with the leading-order $K_0 = -k\, h$. The $J$~contribution~\eqref{eq:4-source-J} is independent of the equation of state~$w$; the $P$~contribution~\eqref{eq:4-source-P} is independent of both $w$ and~$\dot{H}$, since the $00$ projection of $P_{\mu\rho\nu\sigma}\, K^{\rho\sigma}$ involves only $P_{0j0j} \propto H^2$.

The total GB source is $\mathcal{S}^0_{\ 0} = -9k\, H^2$. The fraction from the $P$~tensor is:
\begin{equation}
  f_P = \frac{|\mathcal{S}_P|}{|\mathcal{S}_{\mathrm{total}}|}
  = \frac{6k\, H^2}{9k\, H^2} = \frac{2}{3}\,.
  \label{eq:4-fP}
\end{equation}
This ratio is \emph{exact} and \emph{universal}: it holds for all equations of state $w \in (-\infty, \infty)$, for both the isotropic and anisotropic parts of the junction condition, and for the $ij$ component as well as the $00$ component (Appendix~\ref{app:c1-gb-kk}).

\paragraph{$C_{\mathrm{GB}}$ from KK reduction.}
Under KK reduction, the $J$~contribution renormalises the effective gravitational coupling (absorbed into the GB-modified Planck mass $M_4^2$), while the $P$~contribution generates the kinetic correction~$\epsilon_1$ through its coupling of intrinsic brane curvature to the extrinsic curvature. The warp-factor integrals $I_n = \int_0^{y_c} dy\, e^{nA} \approx 1/(nk)$ in the hierarchy limit give:
\begin{equation}
  C_{\mathrm{GB}} = \frac{I_4}{I_2}\,\frac{2k^2}{(A')^2}\bigg|_{\mathrm{brane}} \times f_P
  = \frac{1}{2} \times 2 \times \frac{2}{3} = \frac{2}{3}\,,
  \label{eq:4-cgb-result}
\end{equation}
where $I_4/I_2 = 1/2$ is the KK integral ratio, $2k^2/(A')^2 = 2$ evaluates the warp-factor derivative at the UV brane, and $f_P = 2/3$ is the $P$~fraction~\eqref{eq:4-fP}.

\paragraph{Verification.}
The value $C_{\mathrm{GB}} = 2/3$ has been verified by three independent methods: (i)~symbolic tensor computation (\textsc{SymPy}, diagonal~$K$, FRW intrinsic curvature, source decomposition confirmed across $w \in [-2,\, 2]$; Appendix~\ref{app:c1-gb-kk}); (ii)~direct spectral sum of the product heat kernel with brane boundary correction of $-0.33\%$ (April~2026); and (iii)~KK integral ratio evaluation. The result is exact within the spectral action framework. Key consistency checks:
\begin{enumerate}
  \item $P_{\mu\rho\nu\sigma}\, K^{\rho\sigma} = 0$ on the static flat brane ($\hat{R} = 0$);
  \item $P_{\mu\rho\nu\sigma}\, K^{\rho\sigma}$ is \emph{isotropic} at de~Sitter ($w = -1$), as required by maximal symmetry;
  \item the static fine-tuning~\eqref{eq:4-warp-jc} recovers $\sigma = 6M_5^3 k\bigl(1 - \tfrac{4}{3}\alpha_{\mathrm{GB}} k^2\bigr)$, matching Charmousis \&\ Dufaux~\cite{ch4:Charmousis2002}.
\end{enumerate}


\subsubsection{Step 5: Numerical Evaluation and Error Budget}

With $C_{\mathrm{GB}} = 2/3$ and the spectral-action coupling $\hat{\alpha}$ computed across four cutoff function families (Appendix~\ref{app:c1-gb-kk}). The $d = 5$ Weyl decomposition correction (Section~\ref{sec:4-seeley-dewitt}, Eq.~\eqref{eq:4-24d}) introduces a factor $(\tfrac{13}{2})/11 = 13/22$ relative to the na\"{\i}ve $d = 4$ computation:

\begin{center}
\begin{tabular}{l r r}
\hline
Cutoff function & $\hat{\alpha}$ (corrected, $d{=}5$) & $\epsilon_1 = \hat{\alpha} \times 2/3$ \\
\hline
Sharp      & $0.0163$ & $0.0109$ \\
Gaussian   & $0.0135$ & $0.0090$ \\
Linear     & $0.0130$ & $0.0086$ \\
Quadratic  & $0.0159$ & $0.0106$ \\
\hline
Central value & $\mathbf{0.015}$ & $\mathbf{0.010}$ \\
\hline
\end{tabular}
\end{center}

\noindent
The error budget is:
\begin{itemize}
  \item \textbf{Cutoff function variation} ($\pm 11\%$): This is the dominant uncertainty. The four cutoff functions span $\hat{\alpha} \in [0.013, 0.016]$ (after the $d = 5$ correction), reflecting the dependence on the spectral moments $f_3/f_2$. The range is narrow because the ratio $f_3/f_2 \sim \mathcal{O}(10)$ is constrained by the positivity and decay properties of any physically reasonable cutoff function~$f$.
  \item \textbf{$C_{\mathrm{GB}}$ uncertainty} (exact): The geometric factor $C_{\mathrm{GB}} = 2/3$ is determined exactly by the Davis junction condition algebra, as the rational ratio $f_P = 6/(6+3) = 2/3$ of the $P$-tensor and $J$-tensor sources at $\mathcal{O}(H^2)$. This ratio is independent of the equation of state, and symbolic computation confirms it identically (Section~\ref{sec:4-gb-universality}).
  \item \textbf{Higher-order corrections} ($< 2\%$): Terms of order $\hat{\alpha}^2 \sim 10^{-4}$ are negligible at the current level of precision.
\end{itemize}

The final result is:
\begin{equation}
  \boxed{\epsilon_1 = \hat{\alpha} \times \frac{2}{3} = 0.010 \pm 0.002}
  \label{eq:4-eps1-result}
\end{equation}
with the uncertainty dominated by the cutoff function dependence of~$\hat{\alpha}$.


\subsubsection{Step 6: Equation of State Prediction}

Equation~\eqref{eq:4-peff} breaks the zero kinetic energy theorem at order~$\epsilon_1$, yielding the dark energy equation of state (Chapter~\ref{ch:foundation}, Section~V):
\begin{equation}
  w_0(\zeta_0) = -1 + \frac{C_{\mathrm{KK}}}{\zeta_0}\,,
  \label{eq:4-w0-prediction}
\end{equation}
where $C_{\mathrm{KK}} = (1 + q_0)^2\, \Omega_{\mathrm{DE}}\, \epsilon_1 / [4(1 - q_0)^2] = (1.64 \pm 0.33) \times 10^{-4}$ (Chapter~\ref{ch:foundation}, Eq.~\eqref{eq:1-ckk-bridge}; Planck~2018 fiducial, corrected via $d = 5$ Weyl decomposition) encodes the cosmological parameters and the GB coupling, and $\zeta_0 = \xi\, \phi_0^2 / M_5^3$ is the sole free parameter from brane UV physics. The prediction is a \emph{curve} $w_0(\zeta_0)$, not a single number: $\zeta_0 \to 0$ gives $w_0 \to -1$ (indistinguishable from~$\Lambda\mathrm{CDM}$), while the junction-condition benchmark $\zeta_0 = 0.001$ gives $w_0 = -0.865$ (non-perturbative), consistent with the DESI~DR2 constant-$w$ constraint (Chapter~\ref{ch:observational}).

Gauss--Bonnet universality (Section~\ref{sec:4-gb-universality}), established by Theorem~\ref{thm:4-axiom-preservation}, ensures that this prediction is robust: any compactification within the Meridian framework produces the same $\epsilon_1$ to within a factor of 2--3, and hence the same $w_0(\zeta_0)$ curve to within the stated uncertainty.


\subsection{Topology Independence}
\label{sec:4-topology}

The universality extends to topology because the Gilkey coefficients~\eqref{eq:4-31} are computed from \emph{local} curvature invariants via the asymptotic expansion of the heat kernel diagonal~\cite{ch4:Gilkey1984}. The heat kernel at coincident points depends only on the local geometry (curvature and its derivatives), not on global topology. Topological information enters through the \emph{non-local} part of the heat kernel (the spectrum of $D_5^2$), which contributes to the spectral action at sub-leading order in the $\Lambda \to \infty$ expansion.

\begin{corollary}\label{cor:4-topology}
The parametric prediction $w_0(\zeta_0)$ is independent of the topology of the compact space~$K$.
\end{corollary}


%% ============================================================
\section{Chern--Simons Terms from the APS Index Theorem}
\label{sec:4-chern-simons}

\subsection{Origin in the Spectral Action}
\label{sec:4-cs-origin}

The $a_{7/2}$ boundary Seeley--DeWitt coefficient produces topological terms on the branes. On the warped orbifold, both gravitational and gauge Chern--Simons 3-forms survive the $\mathbb{Z}_2$ projection:

\paragraph{Gravitational CS (both branes):}
\begin{equation}
  S_{\mathrm{CS,grav}} = \theta_{\mathrm{grav}} \left[
    \int_{y=0} \mathrm{CS}_3(\tilde{\omega})
  - \int_{y=y_c} \mathrm{CS}_3(\tilde{\omega})
  \right],
  \label{eq:4-39}
\end{equation}
where
\begin{equation}
  \mathrm{CS}_3(\tilde{\omega})
  = \mathrm{tr}\!\left(\tilde{\omega} \wedge d\tilde{\omega}
    + \tfrac{2}{3}\, \tilde{\omega} \wedge \tilde{\omega} \wedge \tilde{\omega}\right)
  \label{eq:4-40}
\end{equation}
is the gravitational Chern--Simons 3-form~\cite{ch4:ChernSimons1974} built from the intrinsic 4D spin connection.

\paragraph{Gauge CS (IR brane):}
\begin{equation}
  S_{\mathrm{CS,gauge}} = \sum_i \theta_i \int_{y=y_c} \mathrm{CS}_3(A_i),
  \label{eq:4-41}
\end{equation}
where $i$ runs over the Standard Model gauge group factors: $\mathrm{SU}(3)_c$, $\mathrm{SU}(2)_L$, $\mathrm{U}(1)_Y$. Post-electroweak symmetry breaking, the physically relevant terms are $\mathrm{CS}_3(A_{\mathrm{em}})$ (electromagnetic), $\mathrm{CS}_3(Z)$ ($Z$-boson, massive), and $\mathrm{CS}_3(A_3)$ (gluon, confined).

\paragraph{$\mathbb{Z}_2$ selection rule.} The mixed bulk-brane CS terms---those containing odd powers of $\omega^{a5}$---are $\mathbb{Z}_2$-odd and are projected out. This does \emph{not} affect the purely 4D gravitational and gauge CS terms, which are $\mathbb{Z}_2$-even.


\subsection{The Theta Angles}
\label{sec:4-theta}

\paragraph{Gravitational theta angle.} From the $a_{7/2}$ coefficient:
\begin{equation}
  \theta_{\mathrm{grav}} = \frac{d_S}{192\pi^2} \cdot f_3 \cdot \Lambda^{3/2},
  \label{eq:4-42}
\end{equation}
where $d_S = 4$ is the spinor dimension and $f_3 = -f'(0) \sim \mathcal{O}(1)$ is a moment of the cutoff function. The dimensionless gravitational theta angle is:
\begin{equation}
  \theta_{\mathrm{grav}}^{\mathrm{dimless}} = \frac{\theta_{\mathrm{grav}}}{M_{\mathrm{Pl}}^2}
  \sim f_3 \cdot 10^{-15}.
  \label{eq:4-43}
\end{equation}
This is Planck-suppressed: 15 orders of magnitude below $\mathcal{O}(1)$.

\paragraph{Gauge theta angle.} On the IR brane, the physical cutoff is warped:
\begin{equation}
  \Lambda_{\mathrm{IR}} = \Lambda \cdot e^{-ky_c}
  \sim 10^{17} \cdot 10^{-17} \sim 1\;\mathrm{GeV}.
  \label{eq:4-44}
\end{equation}
The gauge CS coupling is evaluated at the \emph{brane-localized} cutoff:
\begin{equation}
  \theta_{\mathrm{EM}} \sim \mathcal{O}(1).
  \label{eq:4-45}
\end{equation}
This is \textbf{unsuppressed}. The warped redshift brings the spectral cutoff to the TeV scale on the IR brane, where the Standard Model lives. The electromagnetic theta angle is $\mathcal{O}(1)$---not small.

\begin{remark}
The NCG spectral action also predicts $\theta_{\mathrm{QCD}} = 0$ at tree level, providing a potential resolution of the strong CP problem within the RS framework. The Adler--Bardeen theorem~\cite{ch4:AdlerBardeen1969} protects this prediction from perturbative corrections; however, non-perturbative contributions from the fermion determinant (the eta invariant) may generate a small effective $\theta_{\mathrm{QCD}}$, which we do not compute here.
\end{remark}


\subsection{Topological Locking via the APS Index Theorem}
\label{sec:4-aps}

The gravitational and gauge CS terms arise from the \emph{same} Seeley--DeWitt coefficient $a_{7/2}$. Their relationship is governed by the Atiyah--Patodi--Singer (APS) index theorem~\cite{ch4:APS1975a}:

\begin{theorem}[APS Index]\label{thm:4-2}
For the Dirac operator $D_5$ on the warped orbifold $M_4 \times [0, y_c]$, the index is:
\begin{equation}
  \mathrm{Index}(D_5) = \int_{M_4 \times I} \hat{A}(R_5)\, \mathrm{ch}(F)
                      - \tfrac{1}{2}\bigl[\eta(D_{\mathrm{bdy}}) + \dim\ker D_{\mathrm{bdy}}\bigr],
  \label{eq:4-46}
\end{equation}
where $\hat{A}$ is the Dirac genus, $\mathrm{ch}$ is the Chern character, $F$ is the gauge field strength, and $\eta(D_{\mathrm{bdy}})$ is the eta invariant of the boundary Dirac operator.
\end{theorem}

The eta invariant decomposes as:
\begin{equation}
  \eta(D_{\mathrm{bdy}}) = \eta_{\mathrm{grav}} + \eta_{\mathrm{gauge}} + \eta_{\mathrm{cross}},
  \label{eq:4-47}
\end{equation}
and the index is an integer. This has a physical consequence:

\begin{corollary}[Topological Locking]\label{cor:4-locking}
Under continuous deformation of the gauge field (e.g., turning on an electromagnetic field), $\eta_{\mathrm{gauge}}$ changes continuously. Since $\mathrm{Index}(D_5)$ is an integer, $\eta_{\mathrm{grav}}$ must compensate:
\begin{equation}
  \Delta \eta_{\mathrm{grav}} = -\Delta \eta_{\mathrm{gauge}} \pmod{2\mathbb{Z}}.
  \label{eq:4-48}
\end{equation}
The gravitational sector responds to gauge field changes to maintain the integer index.
\end{corollary}

This is the mathematical foundation for the EM-gravity topological coupling: it is \emph{not} a perturbative interaction but a topological constraint---the geometry must adjust to preserve the spectral index.


\subsection{Dynamical Theta from the Cuscuton}
\label{sec:4-dynamical-theta}

The non-minimal coupling $F(\phi) = M_{\mathrm{Pl}}^2 - \xi\phi^2$ makes the gravitational theta angle field-dependent:
\begin{equation}
  \theta_{\mathrm{grav}}(\phi) = \theta_{\mathrm{grav}}^{(0)} \cdot
  \left(\frac{F_0}{F(\phi)}\right)^{3/8}.
  \label{eq:4-49}
\end{equation}

For small $\xi\phi^2$:
\begin{equation}
  \theta_{\mathrm{grav}}(\phi) \approx \theta_{\mathrm{grav}}^{(0)}
  \left(1 + \frac{3}{8}\, \frac{\xi\phi^2}{M_{\mathrm{Pl}}^2} + \cdots\right),
  \label{eq:4-50}
\end{equation}
producing a gradient:
\begin{equation}
  v_\mu = \partial_\mu \theta_{\mathrm{grav}}
        = \theta_{\mathrm{grav}}^{(0)} \cdot \frac{3}{4}\,
          \frac{\xi\phi\,(\partial_\mu\phi)}{M_{\mathrm{Pl}}^2}.
  \label{eq:4-51}
\end{equation}

\paragraph{Cosmological value:} $\dot{\phi}/(H_0\, M_{\mathrm{Pl}}) \sim 0.065$ (suppressed by the zero kinetic energy theorem), giving:
\begin{equation}
  v_0 \sim \theta_{\mathrm{grav}}^{(0)} \cdot 1.9 \times 10^{-3} \cdot H_0 \cdot M_{\mathrm{Pl}}.
  \label{eq:4-52}
\end{equation}

This produces an energy density modification:
\begin{equation}
  \frac{\Delta(H^2)}{H^2} \sim \theta_{\mathrm{grav}}^{\mathrm{dimless}} \cdot \frac{v_0}{H_0^2}
  \sim 10^{-97}.
  \label{eq:4-53}
\end{equation}
Cosmologically negligible.


%% ============================================================
\section{Electromagnetic-Gravity Coupling Channels}
\label{sec:4-em-gravity}

\subsection{Channel Census}
\label{sec:4-channel-census}

The CS terms from Section~\ref{sec:4-chern-simons} open potential channels for electromagnetic fields to influence gravitational physics. We systematically evaluate five perturbative channels and one non-perturbative channel:

\begin{table*}[htbp]
\centering
\caption{EM-gravity coupling channels from the CS terms.}
\label{tab:4-2}
\begin{tabular}{@{}clll@{}}
\toprule
\# & Channel & Mechanism & Result \\
\midrule
1 & Dynamical CS (cosmological)     & $\theta_{\mathrm{grav}}(\phi)$ variation   & $\delta(H^2)/H^2 \sim 10^{-97}$ \\
2 & Dynamical CS (local, 4D)        & $\theta_{\mathrm{grav}}$ gradient response & $\delta g/g \sim 10^{-50}$ \\
3 & Dynamical CS (local, 5D)        & Warped enhancement              & $\delta g/g \sim 10^{-28}$ \\
4 & Anomaly matching (perturbative) & $\eta_{\mathrm{gauge}} \to \eta_{\mathrm{grav}}$ & Planck-suppressed \\
5 & Direct Pontryagin (perturbative)& $\mathbf{E}\cdot\mathbf{B} \to {}^*\!RR$ & $\delta g/g \sim 10^{-77}$ \\
6 & Non-perturbative topological    & Spectral flow                   & Open \\
\bottomrule
\end{tabular}
\end{table*}


\subsection{Perturbative Channel Analysis}
\label{sec:4-perturbative}

\paragraph{Channels 1--3 (Dynamical Chern--Simons~\cite{ch4:AlexanderYunes2009,ch4:JackiwPi2003}).}
These channels require $\theta$ to vary in spacetime. The spectral action gives a fixed $\theta$ on each brane; variation comes only through the cuscuton coupling~\eqref{eq:4-51}. The cosmological channel~(1) is killed by the $10^{-97}$ suppression. The local 4D channel~(2) uses the anomaly current:
\begin{equation}
  \partial_\mu j^\mu_{\mathrm{CS,gauge}}
  = \frac{1}{2\pi^2}\, \mathbf{E} \cdot \mathbf{B},
  \label{eq:4-54}
\end{equation}
which through anomaly matching forces:
\begin{equation}
  \partial_\mu j^\mu_{\mathrm{CS,grav}}
  = -\frac{1}{2\pi^2}\, \mathbf{E} \cdot \mathbf{B}.
  \label{eq:4-55}
\end{equation}

However, the \emph{gravitational response} to this current reintroduces Planck suppression through the Einstein equations:
\begin{equation}
  \delta g \sim \frac{\alpha_{\mathrm{EM}}}{\pi}
             \cdot \frac{\mathbf{E} \cdot \mathbf{B}}{m_e^2}
             \cdot \frac{G_N}{\hbar c}
  \sim 10^{-50}.
  \label{eq:4-56}
\end{equation}

The 5D warped channel~(3) replaces $M_{\mathrm{Pl}}$ with $M_5 \sim 10^8\;\mathrm{GeV}$ in the denominator:
\begin{equation}
  \delta g \sim 10^{-50} \cdot \left(\frac{M_{\mathrm{Pl}}}{M_5}\right)^2
  \sim 10^{-50} \cdot 10^{22} \sim 10^{-28}.
  \label{eq:4-57}
\end{equation}
Still negligible, but the 22-order-of-magnitude warped enhancement demonstrates that the 5D geometry partially compensates for the Planck suppression.

\paragraph{Channel 4 (Anomaly matching).}
The APS index constraint~\eqref{eq:4-48} is exact. However, the \emph{energy density} associated with the gravitational eta invariant change is:
\begin{equation}
  \langle T^{00}_{\mathrm{anom}} \rangle \sim \frac{\alpha_{\mathrm{EM}}}{\pi}
  \cdot \frac{\mathbf{E} \cdot \mathbf{B}}{m_e^2}
  \cdot \frac{\hbar c}{L^4},
  \label{eq:4-58}
\end{equation}
where $L$ is the characteristic scale. For laboratory fields ($E \sim 10^7\;\mathrm{V/m}$, $B \sim 10\;\mathrm{T}$):
\begin{equation}
  \frac{\delta\rho}{\rho_{\mathrm{crit}}} \sim 10^{-47}.
  \label{eq:4-59}
\end{equation}
Negligible by 47 orders of magnitude.

\paragraph{Channel 5 (Direct Pontryagin).}
The electromagnetic Pontryagin density is:
\begin{equation}
  F \wedge F = -4\, \mathbf{E} \cdot \mathbf{B}\, d^4x.
  \label{eq:4-60}
\end{equation}
In Lorentzian signature, this is real-valued and not quantized (unlike Euclidean instantons). For parallel $\mathbf{E}$ and $\mathbf{B}$ fields in a laboratory volume $V \sim 10^{-3}\;\mathrm{m}^3$ with pulse duration $\tau \sim 10^{-6}\;\mathrm{s}$:
\begin{equation}
  Q_{\mathrm{CS}} \sim \frac{\mathbf{E} \cdot \mathbf{B} \cdot V \cdot \tau}{4\pi^2}
  \sim 2500
  \label{eq:4-61}
\end{equation}
in natural units. However, the gravitational response to this Pontryagin charge is suppressed by $\rho_{\mathrm{EM}} / M_{\mathrm{Pl}}^4 \sim 10^{-77}$.


\subsection{Non-Perturbative Channel}
\label{sec:4-nonperturbative}

\paragraph{Channel 6 (Spectral flow under topological transitions).}
All five perturbative channels share a common suppression: the gravitational \emph{response} to topological charge reintroduces Planck suppression through the classical Einstein equations. The non-perturbative channel bypasses this by operating through \textbf{spectral flow}---a change in the spectrum of the boundary Dirac operator $D_{\mathrm{bdy}}$ that shifts eigenvalues through zero.

Spectral flow is a non-perturbative phenomenon: it corresponds to a topological transition (instanton, sphaleron, or domain wall crossing) in which the number of zero modes changes. The APS theorem guarantees that such a transition must be accompanied by a compensating gravitational change~\eqref{eq:4-48}, but the magnitude of this change cannot be computed perturbatively. The semiclassical instanton action $S_{\mathrm{inst}} \sim M_5^3/k^2 \sim 10^{14}$ exceeds the observability threshold $S \sim 189$ by 12 orders of magnitude, making any non-perturbative contribution to $w_0$ undetectable.

\textbf{What is known:}
\begin{itemize}
  \item The mechanism is mathematically rigorous (APS theorem).
  \item Lorentzian topological charge is continuous, not quantized---no energy barrier in the classical gauge sector.
  \item The response requires computing the full eta invariant in the warped RS background, which has not been done.
\end{itemize}

\textbf{What is unknown:}
\begin{itemize}
  \item Whether the gravitational response is large enough to be physically relevant.
  \item Whether laboratory electromagnetic fields can drive sufficient spectral flow.
  \item Whether the non-perturbative gravitational response scales with $M_5$ rather than $M_{\mathrm{Pl}}$.
\end{itemize}

\paragraph{Assessment.} We classify Channel~6 as \textbf{formally open but quantitatively bounded}. The semiclassical instanton action is $S_{\mathrm{inst}} \sim M_5^3/k^2 \sim 10^{14}$, far exceeding the observability threshold $S \sim 189$ (corresponding to $|\delta_w| \sim 10^{-3}$). No constructive mechanism has been identified that could produce macroscopically relevant effects from this channel. The cuscuton constraint is radiatively stable (symmetry protection, Dirac constraint topology, and geometric origin from the $\mathbb{Z}_2$ orbifold; see Chapter~\ref{ch:nogo}, Section~X.C), eliminating the quantization channel as an alternative source of corrections. A definitive closure requires computing the full eta invariant on the warped orbifold---a well-posed but substantial mathematical physics calculation that we defer to future work.


%% ============================================================
\section{Standard Model Connection}
\label{sec:4-sm-connection}

\subsection{The Brane Spectral Triple}
\label{sec:4-brane-triple}

The Layer~2 spectral triple $(\mathcal{A}_{\mathrm{brane}}, \mathcal{H}_{\mathrm{brane}}, D_{B \times F})$ reproduces the Standard Model through the Chamseddine--Connes construction~\cite{ch4:CC1997,ch4:CCM2007}, with the finite space classified by Barrett~\cite{ch4:Barrett2007} as KO-dimension~6. The internal finite space $F$ has:

\textbf{Algebra:}
\begin{equation}
  \mathcal{A}_F = \mathbb{C} \oplus \mathbb{H} \oplus M_3(\mathbb{C}),
  \label{eq:4-62}
\end{equation}
where $\mathbb{C}$ is the complex numbers, $\mathbb{H}$ the quaternions, and $M_3(\mathbb{C})$ the $3 \times 3$ complex matrices. This algebra encodes the gauge group $\mathrm{SU}(3)_c \times \mathrm{SU}(2)_L \times \mathrm{U}(1)_Y$.

\textbf{Hilbert space:}
\begin{equation}
  \mathcal{H}_F = \mathbb{C}^{96},
  \label{eq:4-63}
\end{equation}
encoding 96 fermionic degrees of freedom per generation (left/right $\times$ particle/antiparticle $\times$ color/weak $\times$ 3 generations).

\textbf{Dirac operator:}
\begin{equation}
  D_F = Y,
  \label{eq:4-64}
\end{equation}
the Yukawa coupling matrix, encoding all fermion masses and mixing angles.

\textbf{Inner fluctuations} of the Dirac operator produce the gauge and Higgs fields:
\begin{equation}
  D \to D + A + JAJ^{-1},
  \label{eq:4-65}
\end{equation}
where $A = \sum a_i [D, b_i]$ generates the gauge connection (for $a_i, b_i \in \mathcal{A}_F$). The Higgs doublet arises as the ``discrete gauge field''---the inner fluctuation in the $F$ direction---and its quartic potential is determined by the spectral action~\cite{ch4:CC1997}.


\subsection{The Three-Scalar Architecture}
\label{sec:4-three-scalar}

The Meridian framework contains exactly three scalar fields, each with a distinct NCG origin:

\begin{table*}[htbp]
\centering
\caption{Complete scalar census with NCG identification.}
\label{tab:4-3}
\small
\begin{tabular}{@{}lllcc@{}}
\toprule
Scalar & NCG origin & Location & Prop.? & Role \\
\midrule
Cuscuton $\phi$ & Conformal fluct.\ & 5D bulk & No & Dark energy \\
Radion $T$ & Geometric modulus~\cite{ch4:GoldbergerWise1999} & 4D eff.\ & Yes ($m\!\sim\!\mathrm{TeV}$) & Hierarchy \\
Higgs $H$ & Inner fluct.\ of $D_F$ & IR brane & Yes ($125\;\mathrm{GeV}$) & EWSB \\
\bottomrule
\end{tabular}
\end{table*}

\begin{proposition}[Minimality]\label{prop:4-3}
The three-scalar architecture is the minimal scalar content consistent with the following requirements: (i)~stabilized extra dimension (requires bulk scalar), (ii)~electroweak symmetry breaking (requires Higgs), (iii)~moduli stabilization (radion exists automatically).
\end{proposition}

The cuscuton~\cite{ch4:Afshordi2007} is identified as the conformal fluctuation of the 5D metric---the scalar partner of the graviton in the 5D gravity multiplet. The NCG spectral triple predicts its non-minimal coupling $\xi = 1/6$ (conformal coupling). Crucially, $\xi = 1/6$ is a \emph{consequence} of the spectral action, not a free parameter (see Section~\ref{sec:4-xi-derivation} below). The parameter $\zeta_0$ is free, determined by brane UV physics:
\begin{equation}
  \frac{\xi\phi_0^2}{M_5^3} = \zeta_0\,,
  \label{eq:4-66}
\end{equation}
which for $\xi = 1/6$ gives $\phi_0 = \sqrt{6\zeta_0}\, M_5^{3/2}$. For the junction-condition benchmark $\zeta_0 = 0.001$:
\begin{equation}
  \phi_0 = \sqrt{6 \times 0.001}\, M_5^{3/2} = 0.077\, M_5^{3/2}\,.
  \label{eq:4-67}
\end{equation}
A \textbf{sub-Planckian field value}---natural and consistent with the effective field theory framework. For the CMB-constraint benchmark $\zeta_0 = 0.037$, $\phi_0 = 0.47\, M_5^{3/2}$, also sub-Planckian.


\subsection{The Conformal Coupling \texorpdfstring{$\xi = 1/6$}{xi = 1/6}: Derivation, Not Assumption}
\label{sec:4-xi-derivation}

The results of this section rest on Theorem~\ref{thm:4-axiom-preservation}, which establishes that the junction coupling preserves the NCG axioms required for the brane spectral triple from which $\xi = 1/6$ is derived.

The non-minimal coupling $\xi = 1/6$ is sometimes regarded as a parameter choice. We show here that it is in fact determined by the spectral action principle through three complementary perspectives, all arriving at the same value. These are not independent proofs but different views of a single geometric fact: the scalar is a metric fluctuation, and conformal coupling follows uniquely.

\paragraph{Derivation 1: Seeley--DeWitt $a_2$ coefficient.}
The $a_2$ coefficient for $D^2 = -\nabla^2 + E$ on a $d$-dimensional manifold is~\cite{ch4:Gilkey1984}:
\begin{equation}
  a_2(D^2) = (4\pi)^{-d/2} \int \sqrt{g}\;\left[\tfrac{1}{6}\, R\, \mathrm{tr}(\mathbf{1}) + \mathrm{tr}(E)\right] d^d x.
  \label{eq:4-68}
\end{equation}
For a scalar $\Phi$ with non-minimal coupling $\xi R \Phi^2$, the endomorphism is $E = -\xi R$, so the scalar-gravity coupling in the spectral action reads
\begin{equation}
  S_{\mathrm{gravity}} \supset f_2\, \Lambda^{d-2}\, (4\pi)^{-d/2} \int \sqrt{g}\;\bigl(\tfrac{1}{6} - \xi\bigr)\, R\, \Phi^2\; d^d x.
  \label{eq:4-69}
\end{equation}
In the NCG spectral triple, the endomorphism $E$ is \emph{fixed} by the Dirac operator structure, not freely chosen. The scalar degree of freedom is the conformal fluctuation of the metric: the Weyl factor $\sigma$ in $g \to e^{2\sigma} g$, for which $D_\sigma = e^{-\sigma} D\, e^{-\sigma}$. The conformal anomaly in the heat kernel satisfies
\begin{equation}
  a_2(D_\sigma^2) - a_2(D^2) = (4\pi)^{-d/2} \int \sqrt{g}\;\bigl[0 \cdot R\, \sigma^2\bigr]\; d^d x.
  \label{eq:4-70}
\end{equation}
The vanishing of the $R\, \sigma^2$ term \emph{is} conformal coupling: $\xi = 1/6$. This is not a choice---it is a consequence of the scalar being a metric fluctuation rather than an independent field.

\paragraph{Derivation 2: Radion as metric fluctuation in RS.}
In the Randall--Sundrum background, the radion is the fluctuation of the extra-dimensional metric component: $g_{55} \to 1 + 2\varphi(x)$. Under KK reduction, this becomes a 4D scalar, but one whose coupling to $R_4$ is \emph{determined} by the dimensional reduction, not chosen:
\begin{equation}
  S_{4D} \supset \int \sqrt{g_4}\;\bigl[M_{\mathrm{Pl}}^2/2 - \xi\, \varphi^2\bigr]\, R_4\; d^4x,
  \label{eq:4-71}
\end{equation}
with
\begin{equation}
  \xi = \frac{1}{6}\,\bigl[1 + \mathcal{O}(\zeta_0)\bigr].
  \label{eq:4-72}
\end{equation}
The $\mathcal{O}(\zeta_0)$ corrections arise from the backreaction of the scalar on the warp factor. For the RS orbifold with $ky_c = 35$, the backreaction coefficient is $c_{\mathrm{back}} = 2(1 - e^{-2ky_c})/(ky_c) \approx 0.057$, giving $\delta\xi = (1/6) \cdot c_{\mathrm{back}} \cdot \zeta_0 \approx 9.5 \times 10^{-6}$ (using the junction-condition benchmark $\zeta_0 = 9.64 \times 10^{-4}$)---a fractional correction of $0.006\%$. The effect on $w_0$ is $\delta w_0 \approx 1.6 \times 10^{-8}$, which is negligible compared to $\sigma_{w_0} = 0.002$. This derivation gives the same result as Derivation~1 because the radion \emph{is} the conformal fluctuation of the 5D metric projected to 4D.

\paragraph{Derivation 3: Weyl invariance of the spectral action.}
The $a_2$ term of the spectral action $\mathrm{Tr}[f(D^2/\Lambda^2)]$ is Weyl-invariant if and only if
\begin{equation}
  \xi = \frac{d-2}{4(d-1)}.
  \label{eq:4-73}
\end{equation}
For $d = 4$: $\xi = 2/12 = 1/6$. For $d = 5$: $\xi = 3/16$. The 4D effective theory arises from KK reduction, and the 4D scalar inherits the 4D conformal coupling $\xi_{4D} = 1/6$, not the 5D value, because the scalar's kinetic term is canonically normalized in 4D after integrating over the extra dimension.

\paragraph{Sensitivity analysis.}
From the D2 analysis, $\partial w_0 / \partial \xi \big|_{\xi = 1/6} = -0.0017$. The deviation $\delta\xi$ required to shift $w_0$ by its $1\sigma$ uncertainty ($\sigma_{w_0} = 0.002$) is $\delta\xi \approx 1.15$, corresponding to $690\%$ of $\xi = 1/6$. The prediction is therefore robust: even order-unity errors in $\xi$ would not significantly alter $w_0$.

\paragraph{Summary.}
Three complementary perspectives yield the same result. Derivation~1 (the Seeley--DeWitt $a_2$ coefficient) and Derivation~3 (Weyl invariance) are mathematically equivalent: Weyl invariance of $a_2$ \emph{is} the statement that $\xi = 1/6$ annihilates the $R\sigma^2$ term. Derivation~2 (radion as metric fluctuation) gives the same result because the radion is the conformal fluctuation projected to 4D. The convergence is therefore a consistency check from three angles on a single structural fact, not three independent proofs. Nevertheless, the result elevates $\xi = 1/6$ from a parameter choice to a structural prediction of the spectral action principle. The quoted uncertainty $\sigma_{w_0} = 0.002$ does not need to be enlarged for $\xi$ uncertainty.


\subsection{The Gauge-Gravity Separation}
\label{sec:4-gauge-gravity}

The layered architecture produces a clean separation between the gravitational and gauge sectors:

\begin{enumerate}
  \item \textbf{Gravity (Layer~1):} Einstein--Hilbert $+$ Gauss--Bonnet $+$ cosmological constant $+$ cuscuton kinetic term $+$ radion stabilization. All from the 5D bulk spectral action.

  \item \textbf{Gauge (Layer~2):} $\mathrm{SU}(3) \times \mathrm{SU}(2) \times \mathrm{U}(1)$ gauge fields $+$ Higgs $+$ Yukawa couplings. All from the brane spectral action on $M_4 \times F$.

  \item \textbf{Coupling:} Israel junction conditions~\cite{ch4:Shiromizu2000} $+$ boundary Seeley--DeWitt terms (brane tensions, DGP gravity, CS terms). For Gauss--Bonnet gravity, the junction conditions are modified following Davis~\cite{ch4:Davis2003}.
\end{enumerate}

This separation is not imposed by hand---it is \emph{enforced} by the NCG framework through the odd-dimension obstruction (Proposition~\ref{prop:4-1}). The 5D bulk \emph{cannot} carry gauge fields in the NCG sense; they are naturally confined to the 4D brane.



%% ============================================================
\section{The Octonionic Extension: Three Generations from Geometry}
\label{sec:4-octonionic}

The Chamseddine--Connes--Marcolli (CCM) construction of Sections~\ref{sec:4-brane-triple}--\ref{sec:4-gauge-gravity} takes the number of fermion generations $N_g$ as a free parameter: the finite algebra $\mathcal{A}_F = \mathbb{C} \oplus \mathbb{H} \oplus M_3(\mathbb{C})$ encodes the gauge group $G_{\mathrm{SM}} = \mathrm{SU}(3)_c \times \mathrm{SU}(2)_L \times \mathrm{U}(1)_Y$ but says nothing about $N_g$. One simply declares $\mathcal{H}_F = \mathbb{C}^{32 N_g}$ and sets $N_g = 3$ by hand. In this section we show that extending the finite algebra from the direct sum $\mathbb{C} \oplus \mathbb{H} \oplus M_3(\mathbb{C})$ to the complexified \emph{Dixon algebra}
\begin{equation}
  \mathcal{T}_{\mathbb{C}} = \mathbb{C}_{\mathbb{C}} \otimes_{\mathbb{C}} \mathbb{H}_{\mathbb{C}} \otimes_{\mathbb{C}} \mathbb{O}_{\mathbb{C}}
  \label{eq:4-74}
\end{equation}
derives $N_g = 3$ as an algebraic rigidity result, while preserving the gauge group, the 4D effective spectral action ratios $(C^2, E_4, R^2) = (-18, +11, 0)$, the conformal coupling $\xi = 1/6$, and all bulk physics. (These are the 4D effective ratios from the dimensionally reduced $a_4$ coefficient; the 5D bulk $a_3$ ratios are $(-18, +11, -90)$ before the $R^2 = 0$ identity applies; see Section~\ref{sec:4-ncg-as}, Theorem~\ref{thm:4-critical}.)

The construction rests on four independent algebraic theorems (Hurwitz, Albert, triality uniqueness, and the Fano plane structure), each forcing $N_g \leq 3$ through a different mechanism. The resulting spectral triple $(\mathcal{T}_{\mathbb{C}}, \mathcal{H}_{\mathrm{oct}}, D_{\mathrm{oct}}, J_{\mathrm{oct}}, \gamma_{\mathrm{oct}})$ satisfies all seven NCG axioms---with the first-order condition holding in the Boyle--Farnsworth modified sense for alternative algebras~\cite{ch4:BoyleFarnsworth2020}.%
\footnote{The modified first-order condition used here~\cite{ch4:BoyleFarnsworth2020,ch4:ChamseddineConnesSuijlekom2013} is not universally accepted; see~\cite{ch4:Devastato2018} for an alternative approach that preserves the original first-order condition by enlarging the Hilbert space. The three-generation structure derived in this section depends on the modified condition; the Gauss--Bonnet coupling, radion physics, and all results outside this section do not.} We verify orientability by two independent proofs (envelope surjectivity and explicit Hochschild cycle) and establish Poincar\'{e} duality through the associative envelope $A_{\mathrm{env}}(\mathbb{O}) = M_8(\mathbb{R})$.

The division of labour between the algebraic (octonionic) and geometric (5D warp) sectors is clean: the octonions determine the \emph{structure}---$N_g = 3$, the gauge group, the $S_3$ generation symmetry---while the 5D warp factor determines the \emph{values}---the mass hierarchy and the CKM/PMNS mixing angles.


\subsection{The Dixon Algebra and the Three-Generation Problem}
\label{sec:4-dixon}

\paragraph{The problem.}
Within the CCM framework, the orbifold $S^1/\mathbb{Z}_2$ cannot fix $N_g$. The Atiyah--Patodi--Singer index theorem on $M_4 \times [0, y_c]$ with orbifold boundary conditions yields
\begin{equation}
  \mathrm{Index}(D_5) = \int_{M_4 \times I} \hat{A}(R)\, \mathrm{ch}(F) - \tfrac{1}{2}\bigl[\eta(D_{\partial}) + \dim\ker D_{\partial}\bigr].
  \label{eq:4-75}
\end{equation}
For flat $M_4$ and the topologically trivial interval $I = [0, y_c]$ (which satisfies $H^n(I, \mathbb{Z}) = 0$ for all $n \geq 1$), both the bulk integral and the boundary $\eta$-invariant vanish. There is no topological charge in the extra dimension to source a generation number. The path to fixing $N_g = 3$ therefore requires extending the \emph{finite algebra} $\mathcal{A}_F$ to one whose representation theory forces three generations.

\paragraph{The complexified Dixon algebra.}
The Dixon algebra is the tensor product of the three normed division algebras over $\mathbb{C}$~\cite{ch4:Furey2025,ch4:Dixon1994}:
\begin{equation}
  \mathcal{T}_{\mathbb{C}} = \mathbb{C}_{\mathbb{C}} \otimes_{\mathbb{C}} \mathbb{H}_{\mathbb{C}} \otimes_{\mathbb{C}} \mathbb{O}_{\mathbb{C}},
  \label{eq:4-76}
\end{equation}
where $\mathbb{C}_{\mathbb{C}} = \mathbb{C}$, $\mathbb{H}_{\mathbb{C}} = \mathbb{C} \otimes_{\mathbb{R}} \mathbb{H} \cong M_2(\mathbb{C})$, and $\mathbb{O}_{\mathbb{C}} = \mathbb{C} \otimes_{\mathbb{R}} \mathbb{O}$ is the 8-dimensional complex alternative algebra (non-associative). The complex dimensions are:
\begin{equation}
  \dim_{\mathbb{C}} \mathcal{T}_{\mathbb{C}} = 1 \times 4 \times 8 = 32.
  \label{eq:4-77}
\end{equation}
This is precisely the number of fermionic states (16~particles $+$ 16~antiparticles) in one Standard Model generation. The algebra is alternative (every subalgebra generated by two elements is associative) but not associative, owing to the octonionic factor.

\paragraph{The $\mathbb{Z}_2^5$-grading.}
The Dixon algebra $\mathcal{T} = \mathbb{C} \otimes \mathbb{H} \otimes \mathbb{O}$ carries five independent $\mathbb{Z}_2$ gradings~\cite{ch4:Furey2025}: complex conjugation on $\mathbb{C}$, the main involution on $\mathbb{H}$, octonionic conjugation on $\mathbb{O}$, and two mixed gradings from the tensor product structure. The group $\mathbb{Z}_2^5$ has $2^5 = 32$ elements, and after complexification each grading sector is one-dimensional over $\mathbb{C}$, exhausting $\mathcal{T}_{\mathbb{C}}$. The 32 grading sectors are in one-to-one correspondence with the 32 fermionic states of one SM generation: the quantum numbers (colour, weak isospin, hypercharge, chirality, particle/antiparticle) are determined by the five grading parities.

\paragraph{The KO-dimension chain.}
The Meridian geometry $M_4 \times S^1/\mathbb{Z}_2 \times F$ contributes three Clifford algebra factors:
\begin{equation}
  \mathrm{Cl}(4) \otimes \mathrm{Cl}(1) \otimes \mathrm{Cl}(6) = \mathrm{Cl}(11),
  \label{eq:4-78}
\end{equation}
with the KO-dimension sum $4 + 1 + 6 = 11$. The $\mathbb{Z}_2$ orbifold projection selects the even subalgebra:
\begin{equation}
  \mathrm{Cl}^+(11) \cong \mathrm{Cl}(10).
  \label{eq:4-79}
\end{equation}
After complexification, $\mathrm{Cl}(10)_{\mathbb{C}} = M_{32}(\mathbb{C})$, which has a \emph{unique} irreducible representation of dimension~32---matching $\dim_{\mathbb{C}} \mathcal{T}_{\mathbb{C}}$ exactly. The $\mathbb{Z}_2$ orbifold of the RS geometry naturally produces the $\mathrm{Cl}(10)$ structure within which the octonionic construction lives.


\subsection{Four Proofs of \texorpdfstring{$N_g = 3$}{Ng = 3}}
\label{sec:4-ng-proofs}

The number of fermion generations is determined by the number of independent complex structures on the octonions. A \emph{complex structure} on $\mathbb{O}$ is a linear map $J\colon \mathbb{O} \to \mathbb{O}$ (viewed as $\mathbb{R}^8$) satisfying $J^2 = -\mathrm{Id}$ and compatible with the norm: $|J(x)| = |x|$. Right multiplication by any unit imaginary octonion $e \in \mathrm{Im}(\mathbb{O})$ defines such a map: $J_e(x) = x \cdot e$, with $J_e^2 = -\mathrm{Id}$ following from alternativity and $e^2 = -1$.

While $\mathbb{O}$ has seven imaginary units $e_1, \ldots, e_7$, only three of the corresponding complex structures are independent (up to $G_2 = \mathrm{Aut}(\mathbb{O})$ equivalence after fixing the $\mathrm{SU}(3) = \mathrm{Stab}_{G_2}(e_1)$ subgroup). Each complex structure $J_i$ defines a different $\mathbb{Z}_2^5$-grading of $\mathcal{T}_{\mathbb{C}}$, and each grading yields one generation of SM fermions. The three gradings are related by the $S_3$ subgroup of $\mathrm{Aut}(\mathbb{O})$ that permutes the Fano lines through a given point.

\begin{theorem}[$N_g \leq 3$; Equality from Maximality Postulate]\label{thm:4-3}
The octonions $\mathbb{O}$ possess at most three independent complex structures ($N_g \leq 3$), established by four independent algebraic results. The upper bound is proven. The lower bound $N_g \geq 3$ is not a theorem but a \emph{maximality postulate}: it follows from requiring that the spectral triple exhaust the full algebraic content of the Dixon algebra (representational completeness, see below). Without this postulate, $N_g = 1$ or $N_g = 2$ would be algebraically consistent but representationally incomplete.

\begin{enumerate}
  \item \textbf{Hurwitz's theorem.} $\mathbb{R}$, $\mathbb{C}$, $\mathbb{H}$, and $\mathbb{O}$ are the only normed division algebras over $\mathbb{R}$~\cite{ch4:Hurwitz1898}. The quaternions $\mathbb{H}$ are the largest \emph{associative} normed division algebra, with $\dim_{\mathbb{R}} \mathbb{H} = 4$. The three independent complex structures on $\mathbb{O}$ arise from any quaternionic subalgebra $\mathbb{H} \hookrightarrow \mathbb{O}$: the three imaginary units of $\mathbb{H}$ define three independent $J_i$, and no fourth can be added without either lying in the $G_2$-orbit of an existing one or violating normed-division-algebra constraints.

  \item \textbf{Albert's theorem.} The exceptional Jordan algebra $J_n(\mathbb{O})$ (Hermitian $n \times n$ matrices over $\mathbb{O}$ with symmetrized product) exists only for $n \leq 3$~\cite{ch4:Albert1934,ch4:TodorovDV2018}. For $n = 3$, $J_3(\mathbb{O})$ has rank~3 (every element has exactly three eigenvalues), and $\mathrm{Aut}(J_3(\mathbb{O})) = F_4$. The maximal rank~3 bounds the number of generations that can be accommodated within an octonionic Jordan structure.

  \item \textbf{Triality uniqueness.} Among all simple Lie groups, $\mathrm{Spin}(8) = D_4$ is the \emph{unique} group possessing an outer automorphism of order~3~\cite{ch4:Baez2002}. This triality permutes three inequivalent 8-dimensional representations $\mathbf{8}_v \to \mathbf{8}_s \to \mathbf{8}_c \to \mathbf{8}_v$ and provides the $S_3$ permutation symmetry between generations. Since $\mathrm{Cl}(8) \subset \mathrm{Cl}(10)$ is naturally embedded in the Meridian Clifford chain, the triality structure is intrinsic to the framework.

  \item \textbf{Fano plane structure.} The multiplication table of $\mathrm{Im}(\mathbb{O})$ is encoded in the Fano plane (7~points, 7~lines), where each point lies on exactly 3~lines~\cite{ch4:Baez2002}. The three lines through a given point $e_i$ determine the three complex-conjugate pairs in the decomposition $\mathbb{O} \cong \mathbb{C}^4$ under $J_{e_i}$. There are exactly three independent such decompositions, corresponding to three imaginary units forming an associative triple (a single Fano line).
\end{enumerate}
\end{theorem}

\begin{proof}[Proof outline]
(1)~is classical~\cite{ch4:Hurwitz1898}. (2)~Albert showed that $J_n(\mathbb{O})$ fails to satisfy the Jordan identity for $n \geq 4$ because the non-associativity of~$\mathbb{O}$ makes $n \times n$ Hermitian matrix products ill-defined~\cite{ch4:Albert1934}. (3)~follows from inspection of the Dynkin diagrams: $D_4$ is the unique diagram with $\mathbb{Z}_3$ symmetry (and full outer automorphism group $S_3$). (4)~is verified by direct enumeration: the seven imaginary units modulo the $G_2$ action on $S^6 \subset \mathrm{Im}(\mathbb{O})$ yield $\mathrm{Stab}_{G_2}(e_1) = \mathrm{SU}(3)$, under which $\mathrm{Im}(\mathbb{O}) = \langle e_1 \rangle \oplus V_3 \oplus \bar{V}_3$, and the three independent complex structures are right multiplication by the three imaginary units of any quaternionic subalgebra $\mathbb{H} \hookrightarrow \mathbb{O}$.
\end{proof}

\paragraph{Lower bound: $N_g \geq 3$ from representational completeness.}
The four results above establish $N_g \leq 3$. The lower bound---why not $N_g = 1$ or $N_g = 2$---requires a separate argument. The full Dixon algebra $\mathcal{T}_{\mathbb{C}} = \mathbb{C}_{\mathbb{C}} \otimes \mathbb{H}_{\mathbb{C}} \otimes \mathbb{O}_{\mathbb{C}}$ requires all three independent complex structures on $\mathbb{O}$ to generate the complete 48-state representation (16~states per generation $\times$ 3 generations $= 48$). Using fewer than three complex structures produces a proper subalgebra of the finite spectral triple: with only one complex structure $J_1$, the resulting representation is $\mathcal{H}_1 \cong \mathbb{C}^{32}$ (one generation). With two complex structures, the $S_3$ triality symmetry of $D_4$ (which permutes $\mathbf{8}_v$, $\mathbf{8}_s$, $\mathbf{8}_c$) is broken to $\mathbb{Z}_2$, and the missing third representation means the spectral triple cannot reproduce the full $G_2$-invariant structure of the octonionic multiplication table (the Fano plane requires three lines through each point, not two). We emphasize that one or two generations would still be anomaly-free and internally consistent---the Standard Model is anomaly-free per generation. The lower bound is therefore \emph{not} a no-go theorem against fewer generations; it is a maximality postulate requiring that the spectral triple exhaust the full algebraic content of the Dixon algebra $\mathcal{T}_{\mathbb{C}}$. Fewer than three complex structures give a consistent but \emph{incomplete} spectral triple that does not utilize the octonions' full structure.

\begin{remark}
The three generations come from the octonionic \emph{algebra} (three complex structures on~$\mathbb{O}$), not from the Spin(10) spinor decomposition. Under $\mathrm{Spin}(10) \supset \mathrm{Spin}(8) \times \mathrm{U}(1)$, the Dirac spinor branches as $\mathbf{32} \to (\mathbf{8}_s \oplus \mathbf{8}_c)_{+1/2} \oplus (\mathbf{8}_c \oplus \mathbf{8}_s)_{-1/2}$, yielding only \emph{two} sectors. The $\mathbf{8}_v$ is absent from the spinor decomposition; it enters through the octonionic \emph{multiplication} (the algebra), not through the \emph{representation} (the spinor). Triality provides the $S_3$ symmetry that permutes the three generations---it does not generate them via the spinor branching rule.
\end{remark}


\subsection{The Dirac Operator \texorpdfstring{$D_{\mathrm{oct}}$}{D\_oct}}
\label{sec:4-doct}

The octonionic Hilbert space is $\mathcal{H}_{\mathrm{oct}} = \mathcal{H}_1 \oplus \mathcal{H}_2 \oplus \mathcal{H}_3$, where each $\mathcal{H}_i \cong \mathbb{C}^{32}$ carries the left regular representation of $\mathcal{T}_{\mathbb{C}}$ under the $i$-th $\mathbb{Z}_2^5$-grading (defined by complex structure $J_i$). The total dimension is $\dim_{\mathbb{C}} \mathcal{H}_{\mathrm{oct}} = 96$.

The finite Dirac operator $D_{\mathrm{oct}}$ is a $96 \times 96$ Hermitian matrix with block structure
\begin{equation}
  D_{\mathrm{oct}} = \begin{pmatrix}
    D_{\mathrm{gen}} & \tfrac{1}{2}\, D_{\mathrm{gen}} & \tfrac{1}{2}\, D_{\mathrm{gen}} \\[4pt]
    \tfrac{1}{2}\, D_{\mathrm{gen}} & D_{\mathrm{gen}} & \tfrac{1}{2}\, D_{\mathrm{gen}} \\[4pt]
    \tfrac{1}{2}\, D_{\mathrm{gen}} & \tfrac{1}{2}\, D_{\mathrm{gen}} & D_{\mathrm{gen}}
  \end{pmatrix},
  \label{eq:4-80}
\end{equation}
where $D_{\mathrm{gen}}$ is the standard CCM one-generation Dirac operator (a $32 \times 32$ matrix encoding the Yukawa couplings $Y_\nu$, $Y_e$, $Y_u$, $Y_d$ and the Majorana mass $M_R$ for a single generation~\cite{ch4:CCM2007,ch4:vanSuijlekom2024}). The intra-generation blocks $D_{ii} = D_{\mathrm{gen}}$ are identical for all three generations, guaranteed by the $S_3$ symmetry of the three complex structures. The inter-generation blocks $D_{ij} = \tfrac{1}{2}\, D_{\mathrm{gen}}$ ($i \neq j$) are determined by the \emph{overlap matrix} $M_{\mathrm{oct}}$.

\paragraph{The overlap matrix.}
For each complex structure $J_i$, the complexified octonions decompose as $\mathbb{O}_{\mathbb{C}} = V_i^+ \oplus V_i^-$ (eigenspaces of $J_i$ with eigenvalues $\pm i$), where $\dim_{\mathbb{C}} V_i^{\pm} = 4$. The projector onto $V_i^+$ is $P_i = \tfrac{1}{2}(\mathrm{Id} - i\, J_i)$. The mixing matrix is
\begin{equation}
  (M_{\mathrm{oct}})_{ij} = \frac{1}{\dim V^+}\, \mathrm{Tr}(P_i P_j) = \frac{1}{4}\, \mathrm{Tr}(P_i P_j).
  \label{eq:4-81}
\end{equation}

\begin{theorem}[Democratic Inter-Generation Mixing]\label{thm:4-4}
The overlap matrix is the democratic matrix
\begin{equation}
  M_{\mathrm{oct}} = \frac{1}{2}(I_3 + J_3) = \begin{pmatrix}
    1 & 1/2 & 1/2 \\ 1/2 & 1 & 1/2 \\ 1/2 & 1/2 & 1
  \end{pmatrix},
  \label{eq:4-82}
\end{equation}
where $J_3$ is the $3 \times 3$ all-ones matrix. Its eigenvalues are
\begin{equation}
  \lambda_1 = \lambda_2 = \tfrac{1}{2} \quad (\text{doubly degenerate}), \qquad \lambda_3 = 2,
  \label{eq:4-83}
\end{equation}
with eigenvectors spanning the two-dimensional irreducible and the trivial representations of~$S_3$, respectively.
\end{theorem}

\begin{proof}
The three complex structures $J_1, J_2, J_4$ (right multiplication by $e_1, e_2, e_4$) are related by the $S_3$ subgroup of $\mathrm{Aut}(\mathbb{O}) = G_2$ permuting the three Fano lines through any given point. This $S_3$ acts transitively on $\{J_1, J_2, J_4\}$, forcing: (i)~$\mathrm{Tr}(P_i P_j)$ to be the same for all pairs $i \neq j$; (ii)~$\mathrm{Tr}(P_i P_i)$ to be the same for all~$i$. The diagonal entries are $1$ by normalization ($P_i^2 = P_i$, $\mathrm{Tr}(P_i) = 4$, so $\mathrm{Tr}(P_i P_i)/4 = 1$). The off-diagonal entries are $1/2$: the singular values of the overlap matrices $U_{ij} = (V_i^+)^\dagger V_j^+$ are uniformly $1/\sqrt{2}$ for all three pairs, meaning the $+i$~eigenspaces of different complex structures overlap at exactly $45^\circ$---maximally non-aligned subject to being 4-dimensional subspaces of $\mathbb{C}^8$. The eigenvalue computation is elementary: $M_{\mathrm{oct}} = \tfrac{1}{2} I_3 + \tfrac{1}{2} J_3$, and $J_3$ has eigenvalues $\{0, 0, 3\}$.
\end{proof}

\begin{remark}
The democratic form of $M_{\mathrm{oct}}$ is forced by the $S_3$ symmetry---it has zero free parameters. The eigenvalue ratio $1\!:\!1\!:\!4$ is not the observed fermion mass hierarchy (which has ratios $\sim\!10^{-5}\!:\!10^{-2}\!:\!1$ for up-type quarks). The observed hierarchy arises from the 5D warp factor (Section~\ref{sec:4-yukawa-assessment}).
\end{remark}

\paragraph{Properties of $D_{\mathrm{oct}}$.}
The operator $D_{\mathrm{oct}}$ is Hermitian, has rank~96 (full rank, no zero modes), and its spectrum consists of 48~positive and 48~negative eigenvalues. The real structure compatibility $\|D_{\mathrm{oct}} - J_{\mathrm{oct}}\, D_{\mathrm{oct}}^*\, J_{\mathrm{oct}}^{-1}\| = 0$ holds exactly for KO-dimension~6. The grading condition $\{D_{\mathrm{oct}}, \gamma_{\mathrm{oct}}\} = 0$ is satisfied in the proper CCM basis where the chirality decomposition ($\gamma_{\mathrm{oct}} = \mathrm{diag}(+I_{48}, -I_{48})$) aligns with the particle--antiparticle decomposition: since $D_{\mathrm{gen}}$ is off-diagonal in the chirality basis (the Yukawa matrix connects $L$ to $R$), the inter-generation mixing $M_{\mathrm{oct}} \otimes D_{\mathrm{gen}}$ inherits this off-diagonal structure.


\subsection{Poincar\'{e} Duality via the Associative Envelope}
\label{sec:4-poincare}

The Poincar\'{e} duality axiom requires non-degeneracy of the intersection form on $K_*(\mathcal{A}_F) \times K_*(\mathcal{A}_F) \to \mathbb{Z}$. For the associative CCM algebra $\mathcal{A}_F = \mathbb{C} \oplus \mathbb{H} \oplus M_3(\mathbb{C})$, this follows from the standard K-theory of direct sums of matrix algebras: $K_0(\mathcal{A}_F) = \mathbb{Z}^3$ and the intersection form is a non-degenerate $3 \times 3$ integer matrix~\cite{ch4:CCM2007}. For the non-associative Dixon algebra $\mathcal{T}_{\mathbb{C}}$, standard K-theory does not directly apply: matrix algebras $M_n(\mathbb{O}_{\mathbb{C}})$ are ill-defined for $n > 1$ because non-associativity prevents consistent matrix multiplication.

The resolution passes through the \emph{associative envelope} of the octonions.

\begin{theorem}[Associative Envelope and Poincar\'{e} Duality]\label{thm:4-5}
\leavevmode
\begin{enumerate}
  \item[\emph{(a)}] The associative envelope of the octonions is $A_{\mathrm{env}}(\mathbb{O}) \cong M_8(\mathbb{R})$.
  \item[\emph{(b)}] The associative envelope of the full Dixon algebra is $A_{\mathrm{env}}(\mathcal{T}_{\mathbb{C}}) \cong M_{16}(\mathbb{C})$.
  \item[\emph{(c)}] $K_0(A_{\mathrm{env}}(\mathcal{T}_{\mathbb{C}})) = \mathbb{Z}$ and $K_1(A_{\mathrm{env}}(\mathcal{T}_{\mathbb{C}})) = 0$ (Morita equivalence $M_{16}(\mathbb{C}) \sim \mathbb{C}$).
  \item[\emph{(d)}] The representation K-theory $K_0(\pi(\mathcal{A}_{\mathrm{oct}}) \text{ on } \mathcal{H}_{\mathrm{oct}})$ is isomorphic to~$\mathbb{Z}^3$, matching the CCM result.
  \item[\emph{(e)}] The intersection form is non-degenerate, and Poincar\'{e} duality holds.
\end{enumerate}
\end{theorem}

\begin{proof}
For~(a): the envelope is generated by the left and right multiplication operators $\{L_{e_i}, R_{e_i} : i = 0, \ldots, 7\}$, which are elements of $\mathrm{End}_{\mathbb{R}}(\mathbb{O}) = M_8(\mathbb{R})$. The linear span of $\{L_i, R_i\}$ alone has dimension~15 ($\mathrm{Id}$ plus 7~left and 7~right operators). Including pairwise products $\{L_i L_j, L_i R_j, R_i L_j, R_i R_j\}$ increases the dimension to $64 = \dim M_8(\mathbb{R})$, so the envelope is all of $M_8(\mathbb{R})$.

For~(b): the complexified factors satisfy $(\mathbb{C}_{\mathbb{C}})_{\mathrm{env}} = \mathbb{C}$ and $(\mathbb{H}_{\mathbb{C}})_{\mathrm{env}} = M_2(\mathbb{C})$ (both already associative), and $A_{\mathrm{env}}(\mathbb{O}_{\mathbb{C}}) = M_8(\mathbb{C})$ (complexification of part~(a)). Tensoring gives $\mathbb{C} \otimes M_2(\mathbb{C}) \otimes M_8(\mathbb{C}) = M_{16}(\mathbb{C})$.

For~(c): $M_{16}(\mathbb{C})$ is Morita equivalent to $\mathbb{C}$, giving $K_0 = \mathbb{Z}$ and $K_1 = 0$.

For~(d): the gauge group extraction (Section~\ref{sec:4-gauge-oct}) gives $G_{\mathrm{SM}} = \mathrm{SU}(3) \times \mathrm{SU}(2) \times \mathrm{U}(1)$, unchanged from CCM. The representation of $\mathcal{A}_{\mathrm{oct}}$ on $\mathcal{H}_{\mathrm{oct}} = \mathbb{C}^{96}$ decomposes into the same irreducible pieces as in CCM (three generations of $16 + 16$ states with identical gauge quantum numbers), so the K-theoretic data is isomorphic.

Part~(e) follows from~(d) and the explicit CCM intersection form computation~\cite{ch4:CCM2007}.
\end{proof}

\begin{remark}
The octonionic extension is a ``soft deformation'' in the K-theoretic sense. The non-associativity introduces corrections proportional to the octonionic associator $[a, b, c] = (ab)c - a(bc)$, but these corrections: (i)~preserve representation dimensions ($\dim \mathcal{H}_{\mathrm{oct}} = 96$); (ii)~preserve the gauge group; (iii)~vanish on the associative subalgebra generating the gauge group; and (iv)~do not change the homotopy type of the algebra. The K-groups are therefore expected to be stable under the non-associative deformation. A rigorous proof of homotopy stability for non-associative deformations of spectral triples does not yet exist in the literature; we treat part~(d) of Theorem~\ref{thm:4-5} as strongly supported by the representation-level matching but conditional on the validity of this stability assumption. See Open Problem~6.
\end{remark}


\subsection{Orientability and the \texorpdfstring{$1/6$}{1/6} Coefficient}
\label{sec:4-orientability-oct}

The orientability axiom for a spectral triple $(\mathcal{A}, \mathcal{H}, D, J, \gamma)$ requires the existence of a Hochschild $n$-cycle $c \in Z_n(\mathcal{A}, \mathcal{A} \otimes \mathcal{A}^{\mathrm{op}})$ with $\pi(c) = \gamma$, where $\gamma$ is the grading and $n$ is the metric dimension~\cite{ch4:Connes1994}. For the finite spectral triple (metric dimension~0), the Hochschild cycle is a 0-cycle---an element $c \in \mathcal{A} \otimes \mathcal{A}^{\mathrm{op}}$ satisfying the trivial condition $d_0(c) = 0$ (since $d_0$ maps to the trivial degree~$-1$ module). The axiom thus reduces to a \emph{surjectivity condition}: does $\gamma$ lie in the image of the bimodule representation $\pi\colon \mathcal{A} \otimes \mathcal{A}^{\mathrm{op}} \to \mathrm{End}(\mathcal{H})$?

For non-associative algebras, three complications arise: (i)~the standard Hochschild differential satisfies $d^2 \neq 0$ at degrees $\geq 2$; (ii)~the nucleus $N(\mathbb{O}) = \{n \in \mathbb{O} : [n, a, b] = 0\ \forall\, a, b\} = \mathbb{R} \cdot e_0$ is too small to contain the grading; (iii)~the bimodule structure involves parenthesisation. However, at degree~0 none of these obstructions apply: $d_0 = 0$ regardless of associativity, and the surjectivity of the bimodule map can be established directly.

\begin{theorem}[Orientability of the Octonionic Spectral Triple]\label{thm:4-6}
The orientability axiom holds for $(\mathcal{T}_{\mathbb{C}}, \mathcal{H}_{\mathrm{oct}}, D_{\mathrm{oct}}, J_{\mathrm{oct}}, \gamma_{\mathrm{oct}})$, established by two independent proofs.

\textbf{Proof~A (Associative envelope surjectivity).}
From Theorem~\ref{thm:4-5}(a), $A_{\mathrm{env}}(\mathbb{O}) = M_8(\mathbb{R})$. The bimodule map $\mu\colon \mathbb{O} \otimes_{\mathbb{R}} \mathbb{O}^{\mathrm{op}} \to \mathrm{End}_{\mathbb{R}}(\mathbb{O})$, defined by $\mu(a \otimes b^{\mathrm{op}})(x) = (a \cdot x) \cdot b$, has image $\mathrm{span}\{L_{e_i} R_{e_j} : i, j = 0, \ldots, 7\}$. This span has rank~$64 = \dim M_8(\mathbb{R})$, so $\mu$ is surjective. In particular, the grading $\gamma_{\mathbb{O}} = \mathrm{diag}(+1, -1, \ldots, -1)$ (octonionic conjugation: $\gamma_{\mathbb{O}}(x) = x^*$) lies in $\mathrm{Im}(\mu)$. For the full algebra $A_{\mathrm{env}}(\mathcal{T}_{\mathbb{C}}) = M_{16}(\mathbb{C})$, surjectivity of the bimodule map follows from the double commutant theorem (Burnside's theorem for simple algebras). The three-generation grading $\gamma_{\mathrm{oct}} = I_3 \otimes \gamma_{1\mathrm{gen}}$ lies in the image.

\textbf{Proof~B (Explicit Hochschild cycle).}
For a normed division algebra $K$ of real dimension~$n$, the universal trace formula~\cite{ch4:Baez2002} gives
\begin{equation}
  \sum_{i=0}^{n-1} (e_i \cdot x) \cdot e_i = (2 - n)\, x^*,
  \label{eq:4-84}
\end{equation}
where $\{e_0, \ldots, e_{n-1}\}$ is an orthonormal basis and $x^*$ is the conjugation. For $\mathbb{O}$ ($n = 8$):
\begin{equation}
  \sum_{i=0}^{7} (e_i \cdot x) \cdot e_i = -6\, x^*.
  \label{eq:4-85}
\end{equation}
Define the orientation cycle
\begin{equation}
  c_{\mathbb{O}} = -\frac{1}{6}\, \sum_{i=0}^{7} e_i \otimes e_i^{\mathrm{op}} \;\in\; \mathbb{O} \otimes_{\mathbb{R}} \mathbb{O}^{\mathrm{op}}.
  \label{eq:4-86}
\end{equation}
Then
\begin{equation}
  \pi(c_{\mathbb{O}})(x) = -\frac{1}{6}\, \sum_{i=0}^{7} (e_i \cdot x) \cdot e_i = -\frac{1}{6}(-6\, x^*) = x^* = \gamma_{\mathbb{O}}(x).
  \label{eq:4-87}
\end{equation}
The element $c_{\mathbb{O}}$ is a Hochschild 0-cycle trivially ($d_0 = 0$), and the full orientation cycle is $c_{\mathrm{oct}} = I_3 \otimes c_{\mathbb{C}} \otimes c_{\mathbb{H}} \otimes c_{\mathbb{O}}$, where $c_{\mathbb{C}}$ and $c_{\mathbb{H}}$ are the standard orientation cycles for the associative factors.
\end{theorem}

\begin{remark}
The coefficient $-1/6$ in~\eqref{eq:4-86} arises from $1/(n - 2) = 1/(8 - 2) = 1/6$, where $n = 8$ is the octonionic dimension. This is numerically equal to the conformal coupling $\xi = 1/6$ derived in Section~\ref{sec:4-xi-derivation}, which arises from $\xi = (d - 2)/[4(d - 1)] = 2/12 = 1/6$ for spacetime dimension $d = 4$. Both are consequences of dimensionality (8~for octonions, 4~for spacetime) through different mechanisms. Whether this coincidence has deeper significance through the Meridian bulk--brane correspondence---where the 8-dimensional octonionic structure on the brane is related to 4-dimensional spacetime through the RS warped geometry---is an open question noted here as an observation.
\end{remark}

\paragraph{Basis independence.}
The orientation cycle $c_{\mathbb{O}}$ is basis-independent: for any orthonormal basis $\{f_0, \ldots, f_7\}$ of $\mathbb{O}$ with $f_0 = e_0$, the trace operator $T(x) = \sum_i f_i\, x\, f_i$ is invariant under $\mathrm{Aut}(\mathbb{O}) = G_2$, which acts transitively on orthonormal bases of $\mathrm{Im}(\mathbb{O})$.


\subsection{Complete Axiom Verification}
\label{sec:4-axiom-table}

With orientability established, the octonionic spectral triple satisfies all seven NCG axioms:

\begin{table*}[htbp]
\centering
\caption{NCG axiom verification for the octonionic spectral triple $(\mathcal{T}_{\mathbb{C}}, \mathcal{H}_{\mathrm{oct}}, D_{\mathrm{oct}}, J_{\mathrm{oct}}, \gamma_{\mathrm{oct}})$.}
\label{tab:4-4}
\small
\begin{tabular}{@{}clll@{}}
\toprule
\# & Axiom & Status & Key technique \\
\midrule
1 & Compact resolvent & Holds & Finite-dimensional (trivial) \\
2 & First-order condition & Modified$^\dagger$ & Boyle--Farnsworth~\cite{ch4:BoyleFarnsworth2020}: holds up to associator \\
3 & Orientability & Holds & Trace formula $+$ envelope surjectivity (Theorem~\ref{thm:4-6}) \\
4 & Poincar\'{e} duality & Holds & Via $A_{\mathrm{env}} = M_8(\mathbb{R})$, $K_0 = \mathbb{Z}$ (Theorem~\ref{thm:4-5}) \\
5 & Reality & Holds & Conjugations on $\mathbb{C}$, $\mathbb{H}$, $\mathbb{O}$ are anti-involutions \\
6 & Finiteness & Holds & Finite-dimensional (trivial) \\
7 & Regularity & Holds & Finite-dimensional (trivial) \\
\bottomrule
\multicolumn{4}{@{}l}{\footnotesize $^\dagger$\,Modified in the Boyle--Farnsworth sense~\cite{ch4:BoyleFarnsworth2020}; see discussion below.}\\
\end{tabular}
\end{table*}


\paragraph{The modified first-order condition: implications and caveats.}
The first-order condition $[[D, a], Jb^* J^{-1}] = 0$ is one of the seven axioms defining a real spectral triple in the Connes formalism. For the octonionic extension, this condition does \emph{not} hold in the standard sense on the full algebra $\mathcal{T}_{\mathbb{C}}$. The violation is controlled by the octonionic associator:
\begin{equation}
  [[D_{\mathrm{oct}}, a], J_{\mathrm{oct}} b^* J_{\mathrm{oct}}^{-1}] = \Delta(a, b),
  \label{eq:4-foc-violation}
\end{equation}
where $\Delta(a, b)$ is totally antisymmetric and vanishes whenever both $a$ and $b$ lie in the associative subalgebra $\mathbb{C} \oplus \mathbb{H} \oplus M_3(\mathbb{C})$. Since this subalgebra generates the Standard Model gauge group $G_{\mathrm{SM}}$, the physical gauge sector satisfies the first-order condition exactly. The violation is confined to the octonionic directions that do not participate in gauge interactions.

Boyle and Farnsworth~\cite{ch4:BoyleFarnsworth2020} proposed a relaxation of the first-order condition to accommodate non-associative algebras, requiring only that the commutator vanish on the associative envelope. Our construction satisfies this modified condition. However, this is a substantive mathematical caveat, not a minor technicality: the Boyle--Farnsworth modification is itself a subject of active debate in the NCG community. Whether the modified axiom preserves the full reconstruction theorems of Connes' programme---in particular, whether the spectral characterization of manifolds extends to the non-associative setting---remains an open question. The $N_g = 3$ derivation (Theorem~\ref{thm:4-3}) is rigorous \emph{within} the Boyle--Farnsworth framework; its status in the standard Connes axiomatics is conditional on the resolution of this mathematical question.

\subsection{The Yukawa Structure: An Honest Assessment}
\label{sec:4-yukawa-assessment}

The octonionic associator $[a, b, c] = (ab)c - a(bc)$ is totally antisymmetric (a 3-form on $\mathrm{Im}(\mathbb{O})$), vanishes on the 7~Fano triples (associative sub-triples), and is non-zero on exactly 28 of the 35 independent triples of imaginary units~\cite{ch4:Baez2002}. The 35-parameter space $\Lambda^3(\mathbb{R}^7)$ decomposes under $G_2 = \mathrm{Aut}(\mathbb{O})$ as
\begin{equation}
  \Lambda^3(\mathbb{R}^7) = \mathbf{1} \oplus \mathbf{7} \oplus \mathbf{27},
  \label{eq:4-88}
\end{equation}
where the $\mathbf{1}$ is the $G_2$-invariant associative calibration $\varphi$ (non-zero on Fano triples, zero elsewhere) and the associator itself lies in the $\mathbf{7} \oplus \mathbf{27}$ complement.

A critical point: \emph{the octonionic associator has zero free parameters}. It is completely determined by the Fano plane. The 35 independent components are all fixed---28 non-zero values of magnitude~$2$ with determined signs, and 7~zeros on the Fano lines. Therefore the association between ``35~parameters'' and the 20~physical SM mixing parameters (4~CKM, 4~PMNS, 12~fermion masses) is a counting coincidence, not a physical mechanism.

\paragraph{The $S_3$ obstruction.}
\label{sec:4-s3-catastrophe}
The three complex structures $J_1, J_2, J_4$ are related by an exact $S_3$ symmetry, forcing $M_{\mathrm{oct}}$ to be democratic (Theorem~\ref{thm:4-4}). The up-type and down-type quarks see the \emph{same} overlap matrix, so the CKM matrix at the algebraic level is
\begin{equation}
  V_{\mathrm{CKM}}^{\mathrm{alg}} = U_u^\dagger\, U_d = U_M^\dagger\, U_M = I_3,
  \label{eq:4-89}
\end{equation}
where $U_M$ diagonalises $M_{\mathrm{oct}}$. This predicts zero CKM mixing---obviously wrong.

\paragraph{Bulk--brane complementarity.}
The resolution is the 5D warp factor. The effective Yukawa coupling between generations $i$ and $j$ for fermion type~$f$ is
\begin{equation}
  Y_{ij}^{f,\,\mathrm{eff}} = (M_{\mathrm{oct}})_{ij}\, y_f \int_0^{y_c} f_i^f(y)\, f_j^f(y)\, h(y)\, e^{-4A(y)}\, dy,
  \label{eq:4-90}
\end{equation}
where $f_i^f(y)$ is the extra-dimensional profile of the $i$-th generation fermion of type~$f$, $h(y)$ is the Higgs profile, and $A(y) = -ky$ is the warp factor. The profiles depend on the bulk mass parameters~$c_i^f$, which are different for up-type and down-type quarks. The exponential $e^{-c\, k y_c}$ with $ky_c \approx 37$ converts $\mathcal{O}(1)$ differences in~$c$ to enormous mass hierarchies ($e^{-0.5 \times 37} \sim 10^{-8}$, comparable to $m_u/m_t$), and the difference between $Y_u^{\mathrm{eff}}$ and $Y_d^{\mathrm{eff}}$ generates non-trivial CKM mixing.

\begin{table*}[htbp]
\centering
\caption{Division of labour between the algebraic and geometric sectors.}
\label{tab:4-5}
\small
\begin{tabular}{@{}lll@{}}
\toprule
Structure & Determines & Free parameters \\
\midrule
Octonionic (brane) & $N_g = 3$, gauge group, $S_3$ symmetry, Fano topology & 0 \\
5D warp (bulk) & Mass hierarchy, CKM/PMNS mixing, flavour violation & $\sim$20 (bulk masses $c_i^f$) \\
\bottomrule
\end{tabular}
\end{table*}

\begin{table*}[htbp]
\centering
\caption{Cumulative parameter budget at each stage of the derivation chain. ``Total free'' counts all parameters not fixed by the stated axioms.}
\label{tab:4-param-budget}
\small
\begin{tabular}{@{}llcc@{}}
\toprule
Stage & New inputs & New free & Total free \\
\midrule
\multicolumn{4}{@{}l}{\emph{Core cosmological framework (Chapters~1, 3, 5)}} \\
A1--A2: 5D orbifold + scalar & RS geometry, $P(X,\Phi)$ & 0 & 0 \\
A3--A6: sequestering, NCG, $\xi$, tadpole & Spectral action on orbifold & 0 & 0 \\
Self-tuning + cuscuton uniqueness & $w_0(\zeta_0)$ functional form & 0 & 0 \\
GB correction ($\epsilon_1$ from spectral action) & $C_\mathrm{KK}$ & $0^\dag$ & $0^\dag$ \\
GW stabilisation parameter $\varepsilon$ & $\zeta_0$, specific $w_0$ value & 1 & 1 \\[3pt]
\multicolumn{4}{@{}l}{\emph{NCG spectral geometry (Chapter~4, cosmological sector)}} \\
Spectral action on warped orbifold & GB, CS, DGP terms & 0 & 1 \\
$C_\mathrm{GB} = 2/3$, $\hat{\alpha}$ evaluation & $\epsilon_1 = 0.010 \pm 0.002$ & 0 & 1 \\[3pt]
\multicolumn{4}{@{}l}{\emph{Particle physics extensions (Chapter~4, speculative)}} \\
Octonionic Dixon algebra & $N_g = 3$ (upper), $G_\mathrm{SM}$ & 0 & 1 \\
Maximality postulate$^\ddag$ & $N_g = 3$ (lower) & 0 & 1 \\
Quark sector bulk masses $c_i^{u,d}$ & Mass hierarchy, CKM & 9 & 10 \\
Charged lepton bulk masses & $m_e, m_\mu, m_\tau$ & 3 & 13 \\
Yukawa couplings $y_u, y_d$ & Overall scale & 2 & 15 \\
Majorana sector ($a, b, \Delta c_\nu$) & $m_{\nu_i}$, PMNS & 3 & 18 \\
\bottomrule
\multicolumn{4}{@{}l}{\footnotesize $^\dag$\,$\epsilon_1$ has a 20\% irreducible cutoff-function ambiguity (not a free parameter in the usual sense).} \\
\multicolumn{4}{@{}l}{\footnotesize $^\ddag$\,Aesthetic assumption, not a mathematical necessity (see text).}
\end{tabular}
\end{table*}

This is not a weakness but the correct architectural division: the octonions fix what \emph{can} exist (three generations with $S_3$ symmetry and identical gauge quantum numbers), while the 5D geometry fixes what \emph{does} exist in our universe (the specific mass ratios and mixing angles).

\paragraph{Quantitative first-order condition.}
The preceding qualitative picture can be made precise. The NCG first-order condition $[[D_F, \pi(a)], \pi^\circ(b)] = 0$ for the full $8 \times 8$ octonionic mass matrix $M$ becomes
\begin{equation}
  [[M, L_{e_i}], R_{e_j}] = 0 \quad \text{for all } i,j \in \{1,\ldots,7\},
  \label{eq:4-fo-full}
\end{equation}
where $L_{e_i}$, $R_{e_j}$ are left and right multiplication by imaginary octonion units. Direct computation shows that this system eliminates \textbf{63 of the 64} entries of~$M$: the SVD of the constraint matrix yields 63 active singular values and a one-dimensional null space. The single surviving degree of freedom is proportional to the $8 \times 8$ identity---it is $S_3$-invariant. Non-associativity constrains the \emph{shape} of the mass matrix almost completely, but the one free parameter it leaves is democratic across generations.

\paragraph{Basin determination.}
The practical consequence is quantified by a global fit of the $S_3$-doublet-degenerate model ($c_{a1} = c_{a2}$ for each sector) to all 9~charged fermion masses and 6~CKM elements. The result is $\chi^2 = 1554$ for $N_{\mathrm{dof}} = 2$, with $V_{us} = V_{ub} = 0$ identically---the $S_3$ doublet degeneracy forces the first two generations to be indistinguishable, making nontrivial quark mixing impossible. An SVD null-space analysis of the Jacobian reveals 5~unconstrained parameter directions, all traceable to the $S_3$ doublet structure.

\paragraph{The $J_3(\mathbb{O})$ resolution.}
The exceptional Jordan algebra $J_3(\mathbb{O})$ of $3 \times 3$ Hermitian octonionic matrices provides the natural completion. Its Peirce decomposition with respect to three orthogonal primitive idempotents $p_1, p_2, p_3$ gives:
\begin{equation}
  J_3(\mathbb{O}) = \bigoplus_{i} J_{ii} \;\oplus\; \bigoplus_{i < j} J_{ij}, \qquad
  \dim J_{ii} = 1, \quad \dim J_{ij} = 8,
  \label{eq:4-peirce}
\end{equation}
where the off-diagonal Peirce spaces $J_{ij}$ carry the octonionic inter-generation coupling (reproducing $M_{\mathrm{oct}}$) and the diagonal spaces $J_{ii}$ carry generation-specific real parameters unconstrained by~$S_3$. Every element of $J_3(\mathbb{O})$ is $F_4$-conjugate to a unique diagonal element $\mathrm{diag}(\lambda_1, \lambda_2, \lambda_3)$ with $\lambda_1 \leq \lambda_2 \leq \lambda_3$---the three Jordan eigenvalues are the complete $F_4$-orbit invariants.

Fitting the full model with all three $c_i$ independent per sector (18~parameters, 15~observables) yields $\chi^2 = 0.82$, reproducing all 9~masses to within 6\% and all 6~CKM elements to within 1\%. The $S_3$ degeneracy catastrophe is resolved: generation~1--2 splittings range from $|\Delta c| = 0.006$ (down quarks) to $0.19$ (charged leptons). Of the 7~null directions in the Jacobian, 3~are trivially from parameter excess and 4~from reparameterisation degeneracies---none are structural impossibilities.

The division of Table~\ref{tab:4-5} is thus refined:
\begin{itemize}
  \item \textbf{Octonionic first-order condition} ($63/64$): fixes the mass matrix template and $M_{\mathrm{oct}}$ structure (topology).
  \item \textbf{Jordan diagonal $(\lambda_1, \lambda_2, \lambda_3)$}: parameterises the generation-specific bulk masses (metric).
  \item \textbf{RS warp factor}: converts $\mathcal{O}(1)$ parameter differences to the observed mass hierarchy (dynamics).
\end{itemize}
The algebraic structure is fully determined; the specific values are empirical. This is the basin map of the fermion sector.

\begin{remark}[Negative results that sharpen the boundary]
\label{rem:4-negative-results}
Two alternative approaches to resolving the $S_3$ degeneracy have been tested and excluded:
\begin{enumerate}
\item \textbf{$G_2$ cross-sector constraints.} The automorphism group $G_2 = \mathrm{Aut}(\mathbb{O})$ acts on the Fano plane and therefore constrains the octonionic overlap matrix $M_\mathrm{oct}$. However, $G_2$ acts identically on all three complex structures---it does not provide cross-sector constraints that could differentiate quark from lepton masses. The $G_2$-invariant content is exhausted by the democratic structure already present in $M_\mathrm{oct}$.
\item \textbf{$\mathrm{SU}(5)$ GUT embedding.} Constraining all fermion profiles within each generation to a common $\mathrm{SU}(5)$ multiplet ($c_{10} = c_{\bar{5}}$) reduces the parameter count but yields $\chi^2/\mathrm{dof} = 518$---a catastrophic failure. The GUT relations are incompatible with the RS warp-factor mechanism, which requires different bulk masses for different SM representations. This is a genuine prediction: the Meridian framework excludes minimal $\mathrm{SU}(5)$ unification at the level of the fermion profiles.
\end{enumerate}
These negative results are as informative as the positive ones: they map the boundary between what the algebraic structure determines (topology) and what only the geometry can provide (metric). The octonionic NCG is a \emph{selection principle}, not a mass formula.
\end{remark}


\subsection{The Gauge Group from Octonions}
\label{sec:4-gauge-oct}

The gauge group is extracted from the inner automorphisms of the Dixon algebra modulo its centre. The automorphism group of $\mathbb{O}$ is $G_2$~\cite{ch4:Baez2002}, but the choice of a preferred complex structure (equivalently, a preferred imaginary unit $e_1$) breaks $G_2$ to the stabiliser:
\begin{equation}
  \mathrm{Stab}_{G_2}(e_1) = \mathrm{SU}(3).
  \label{eq:4-91}
\end{equation}
This is the colour group. The breaking $G_2 \to \mathrm{SU}(3)$ is automatic once the generation structure is defined---it is not an additional assumption. Combined with the automorphisms of the other Dixon factors:
\begin{equation}
  G_{\mathrm{SM}} = \mathrm{Stab}_{G_2}(e_1) \times \mathrm{Aut}(\mathbb{H}) \times \mathrm{Aut}(\mathbb{C}) = \mathrm{SU}(3) \times \mathrm{SU}(2) \times \mathrm{U}(1),
  \label{eq:4-92}
\end{equation}
matching the CCM gauge group exactly~\cite{ch4:Krasnov2025,ch4:Furey2025}.

\paragraph{Spectral action preservation.}
The gravitational spectral action coefficients depend only on the 5D Dirac operator (the bulk geometry), not on the finite space~$F$. Since the octonionic extension modifies only the brane algebra while $\dim \mathcal{H}_{\mathrm{oct}} = 96 = \dim \mathcal{H}_F^{\mathrm{CCM}}$:
\begin{equation}
  (C^2,\; E_4,\; R^2) = (-18,\; +11,\; 0) \qquad \text{(unchanged)}.
  \label{eq:4-93}
\end{equation}
The $R^2 = 0$ identity (the structural relation $15/12 - 8/12 - 7/12 = 0$) is a property of the 5D Dirac operator and holds for \emph{any} finite space. The conformal coupling $\xi = 1/6$ (Section~\ref{sec:4-xi-derivation}), the self-tuning mechanism, and the sound speed $c_s \in [12c,\, 15c]$ are all bulk-determined and therefore preserved.

\paragraph{The matter spectral action.}
The matter sector $S_{\mathrm{matter}} = \Tr_{\mathcal{H}_F}[f(D_F^2/\Lambda^2)]$ has its leading coefficient
\begin{equation}
  a_0^{\mathrm{matter}} = \Tr(\mathbf{1}_{\mathcal{H}_F}) = 96,
  \label{eq:4-94}
\end{equation}
which is identical in both the CCM ($N_g = 3$, assumed) and octonionic ($N_g = 3$, derived) constructions. The higher Seeley--DeWitt coefficients $a_2^{\mathrm{matter}}$ and $a_4^{\mathrm{matter}}$ depend on the specific form of $D_F$ (the Yukawa matrix). If the inter-generation mixing from $M_{\mathrm{oct}}$ differs from the CCM Yukawa matrix, these coefficients will receive $\mathcal{O}(\tfrac{1}{2})$ corrections from the off-diagonal blocks. The gravitational predictions of Meridian---which depend only on the bulk spectral action and the leading matter trace---are unaffected.


\subsection{Anomaly Cancellation from Octonionic \texorpdfstring{$\mathrm{Spin}(10)$}{Spin(10)}}
\label{sec:4-anomaly}

Each generation of the octonionic spectral triple fills exactly one $\mathbf{16}$-dimensional spinor representation of $\mathrm{Spin}(10)$, containing 16~Weyl fermions: the 15~of the SM ($q_L, u_R, d_R, \ell_L, e_R$) plus $\nu_R$.  The anomaly cancellation of the $\mathbf{16}$ is a representation-theoretic consequence, but we verify all six independent 4D conditions explicitly using exact (fraction) arithmetic---not floating-point approximations.

\begin{table*}[htbp]
\centering
\caption{Anomaly cancellation for one SM generation with $\nu_R$.  All conditions verified with exact fraction arithmetic.}
\label{tab:4-anomaly}
\small
\begin{tabular}{@{}lll@{}}
\toprule
Anomaly & Condition & Result \\
\midrule
$\mathrm{SU}(3)^3$ & Sum of cubic Casimirs & $0$ \\
$\mathrm{SU}(3)^2 \times \mathrm{U}(1)_Y$ & Mixed gauge--hypercharge & $0$ \\
$\mathrm{SU}(2)^2 \times \mathrm{U}(1)_Y$ & Mixed gauge--hypercharge & $0$ \\
$\mathrm{U}(1)_Y^3$ & Cubic hypercharge & $0$ \\
$\mathrm{U}(1)_Y \times \mathrm{grav}^2$ & Mixed gravitational & $0$ \\
Witten $\mathrm{SU}(2)$ global & Number of doublets mod~$2$ & $0$ (12 doublets, even) \\
\bottomrule
\end{tabular}
\end{table*}

\begin{proposition}[Anomaly Freedom Requires $\nu_R$]\label{prop:4-nuR-anomaly}
The right-handed neutrino $\nu_R$ is necessary for anomaly cancellation.  Without it, each generation has $n_L = 8$ and $n_R = 7$, giving:
\begin{enumerate}
  \item[(i)] $\mathrm{SU}(2)^2 \times \mathrm{grav}$ anomaly $= 2$ (two unpaired doublets);
  \item[(ii)] gravitational anomaly $n_L - n_R = 1 \neq 0$.
\end{enumerate}
With $\nu_R$, both cancel: $n_L = n_R = 8$ per generation.  The right-handed neutrino is not a phenomenological addition---it is demanded by 5D consistency.
\end{proposition}

\paragraph{Explicit verification of $\mathrm{U}(1)_Y^3$.}
The cubic hypercharge anomaly is the most non-trivial condition.  With hypercharge assignments in the convention $Q_{\mathrm{em}} = T_3 + Y$, the anomaly is $\mathcal{A}[\mathrm{U}(1)_Y^3] = \sum_{\mathrm{left}} n_c\, Y^3 - \sum_{\mathrm{right}} n_c\, Y^3$, where the sums run over left-handed and right-handed Weyl fermions respectively (each weak doublet component counted separately).  For one generation with $\nu_R$:
\begin{equation}
\begin{aligned}
  \sum_L Y^3 &= 3 \cdot 2 \cdot \bigl(\tfrac{1}{6}\bigr)^3 + 1 \cdot 2 \cdot \bigl({-}\tfrac{1}{2}\bigr)^3 \notag\\
  &= \tfrac{1}{36} - \tfrac{1}{4}
  = -\tfrac{2}{9}\,, \\[4pt]
  \sum_R Y^3 &= 3\bigl(\tfrac{2}{3}\bigr)^3 + 3\bigl({-}\tfrac{1}{3}\bigr)^3 + (-1)^3 + 0^3 \notag\\
  &= \tfrac{8}{9} - \tfrac{1}{9} - 1 + 0
  = -\tfrac{2}{9}\,, \\[4pt]
  \mathcal{A}[\mathrm{U}(1)_Y^3] &= -\tfrac{2}{9} - \bigl(-\tfrac{2}{9}\bigr) = 0\,.
\end{aligned}
  \label{eq:4-anomaly-explicit}
\end{equation}
The cancellation $\sum_L Y^3 = \sum_R Y^3 = -2/9$ is exact, verified with fraction arithmetic.  Without $\nu_R$ ($Y_{\nu_R} = 0$), $\sum_R$ is unchanged and the condition still holds; however, the gravitational anomaly $n_L - n_R = 1 \neq 0$ then fails (Proposition~\ref{prop:4-nuR-anomaly}).

\begin{remark}
Each generation fills exactly one $\mathbf{16}$-spinor of $\mathrm{Spin}(10)$.  The anomaly cancellation therefore follows from the representation theory of $\mathrm{Spin}(10)$, which itself follows from the octonionic Clifford structure $\mathrm{Cl}(10)_{\mathbb{C}} = M_{32}(\mathbb{C})$ (Section~\ref{sec:4-dixon}).  The cancellation is geometric at its root.
\end{remark}


\subsection{Chern--Simons Inflow Consistency}
\label{sec:4-cs-inflow}

The 5D orbifold requires consistency of the Chern--Simons inflow mechanism: the bulk anomaly polynomial must factorise so that brane-localised anomalies are cancelled by inflow from the bulk.  We verify all seven independent checks:

\begin{table*}[htbp]
\centering
\caption{Chern--Simons inflow consistency on the warped $S^1/\mathbb{Z}_2$ orbifold.  All seven checks pass.  Note: Checks 1--2 and 7 are consistency conditions (trivially satisfied when $I_6 = 0$); the non-trivial verifications are Checks 3, 5, and 6.}
\label{tab:4-cs-inflow}
\small
\begin{tabular}{@{}cll@{}}
\toprule
\# & Check & Result \\
\midrule
1 & Anomaly polynomial $I_6$ & $= 0$ (no bulk anomaly) \\
2 & Green--Schwarz factorisation & Trivial ($I_6 = 0$) \\
3 & Rep-by-rep inflow matching & All reps, all 3 generations \\
4 & Warp correction consistency & Topological invariance holds \\
5 & Parity anomaly (integer CS levels) & $k_{\mathrm{SU}(3)} = 6$, $k_{\mathrm{SU}(2)} = 6$, $k_{\mathrm{U}(1)} = 10$ \\
6 & $\mathbb{Z}_2$ orbifold consistency & All 6 brane anomalies vanish \\
7 & Octonionic extension & No extra CS structure introduced \\
\bottomrule
\end{tabular}
\end{table*}

The CS inflow mechanism on the warped orbifold is the 5D manifestation of the 4D anomaly cancellation (Table~\ref{tab:4-anomaly}).  The deeper reason both work is the $\mathrm{Spin}(10)$ embedding from octonionic Clifford structure: it automatically ensures that the fermion content is anomaly-free in any dimension where it can be consistently embedded.  The integer CS levels ($k_{\mathrm{SU}(3)} = 6$, $k_{\mathrm{SU}(2)} = 6$, $k_{\mathrm{U}(1)} = 10$) confirm that the $\mathbb{Z}_2$ orbifold projection is fully consistent with the gauge structure---no additional discrete anomaly matching conditions arise from the orbifold.


%% ============================================================
\section{Fermion Mass Hierarchy from Warping}
\label{sec:4-fermion-mass}

The octonionic democratic matrix $M_{\mathrm{oct}}$ (Section~\ref{sec:4-octonionic}) produces degenerate inter-generation couplings with eigenvalues $\{1/2,\, 1/2,\, 2\}$. The observed mass hierarchy---spanning more than five orders of magnitude from $m_e = 0.511\;\mathrm{MeV}$ to $m_t = 172.69\;\mathrm{GeV}$---must arise from a different mechanism. In the Meridian framework, this mechanism is the Gherghetta--Pomarol (GP) localization~\cite{ch4:GherghettaPomarol2000} of fermion zero modes in the warped extra dimension, embedded within the NCG spectral triple constructed in Section~\ref{sec:4-spectral-triple}.

\subsection{Zero-Mode Profiles and the Profile Overlap Function}
\label{sec:4-zero-mode-profiles}

Each bulk fermion $\Psi_i$ with 5D mass parameter $c_i$ (in units of the AdS curvature $k$) possesses a chiral zero mode whose profile in the extra dimension is determined by the orbifold $\mathbb{Z}_2$ projection (Section~\ref{sec:4-z2-spinors}):
\begin{equation}
  f_L^{(0)}(y) = N_L\, e^{(2-c_i)ky}, \qquad
  f_R^{(0)}(y) = N_R\, e^{(2+c_i)ky},
  \label{eq:4-95}
\end{equation}
with normalization determined by the warped measure $\int_0^{y_c} e^{-4ky}\, |f|^2\, dy = 1$. The effective 4D Yukawa coupling for the $i$-th fermion, obtained by evaluating the brane-localized Higgs overlap, is
\begin{equation}
  Y_i^{\mathrm{eff}} = Y_5 \cdot [g(c_i)]^2, \qquad
  m_i = Y_5\, [g(c_i)]^2 \cdot \frac{v}{\sqrt{2}}\,,
  \label{eq:4-96}
\end{equation}
where $Y_5$ is the universal 5D Yukawa coupling and $g(c)$ is the dimensionless profile overlap function:
\begin{equation}
  g(c) = \sqrt{\frac{1 - 2c}{e^{(1-2c)ky_c} - 1}}\; e^{(1/2 - c)ky_c}\,.
  \label{eq:4-97}
\end{equation}
For $|c - 1/2| \gg 1/(ky_c)$, this reduces to the exponential
\begin{equation}
  g(c) \simeq \sqrt{2c - 1}\; e^{(1/2 - c)ky_c}\,,
  \label{eq:4-98}
\end{equation}
demonstrating that a shift $\Delta c$ in the bulk mass parameter produces a mass ratio $m_1/m_2 \sim e^{2\Delta c\, ky_c}$. With $ky_c = 37$, a spread of $\Delta c \sim 0.65$ in $\mathcal{O}(1)$ parameters suffices to generate the full $3.4 \times 10^5$ hierarchy of charged fermion masses.

\begin{figure*}[htbp]
\centering
\includegraphics[width=0.85\textwidth]{figures/fig5_fermion_profiles.pdf}
\caption{Fermion zero-mode probability density $|\psi(y)|^2$ in the warped extra dimension,
for selected bulk mass parameters $c_i$. The top quark ($c = 0.004$, red) extends nearly flat
across the dimension with large overlap with the IR-localized Higgs, producing an $\mathcal{O}(1)$
Yukawa coupling. The electron ($c = 0.656$, dark blue) is sharply UV-peaked with negligible
Higgs overlap, producing a tiny effective Yukawa. The 20~orders-of-magnitude separation in
probability density at the IR brane generates the observed fermion mass hierarchy
($m_t/m_e = 3.4 \times 10^5$) from $\mathcal{O}(1)$ bulk mass parameters
($\Delta c \sim 0.65$).}
\label{fig:5-fermion-profiles}
\end{figure*}

\subsection{Bulk Mass Parameters for All Fermion Species}
\label{sec:4-ci-table}

Setting $Y_5 = 1.00$ (the natural dimensionless value) and requiring $m_t = 172.69\;\mathrm{GeV}$, the bulk mass parameter $c_{\mathrm{top}} = 0.004$ is determined. All remaining $c_i$ follow from mass ratios. The complete table is:

\begin{table*}[htbp]
\centering
\caption{Bulk mass parameters $c_i$ for all Standard Model fermions. All charged fermion parameters lie in $[0.004,\, 0.656]$, a total spread of $0.65$. The 5D Yukawa coupling is $Y_5 = 1.00$ exactly.}
\label{tab:4-ci}
\small
\begin{tabular}{@{}llccccl@{}}
\toprule
Fermion & Sector & Gen & $c_i$ & $g(c_i)$ & $m_{\mathrm{obs}}$ & Localization \\
\midrule
$t$ & up & 3 & 0.004 & $9.96 \times 10^{-1}$ & 172.69\;GeV & IR brane \\
$b$ & down & 3 & 0.503 & $1.55 \times 10^{-1}$ & 4.18\;GeV & flat \\
$\tau$ & lepton & 3 & 0.523 & $1.01 \times 10^{-1}$ & 1.777\;GeV & flat \\
$c$ & up & 2 & 0.530 & $8.54 \times 10^{-2}$ & 1.27\;GeV & flat \\
$\mu$ & lepton & 2 & 0.574 & $2.46 \times 10^{-2}$ & 105.66\;MeV & UV brane \\
$s$ & down & 2 & 0.576 & $2.32 \times 10^{-2}$ & 93.4\;MeV & UV brane \\
$d$ & down & 1 & 0.623 & $5.18 \times 10^{-3}$ & 4.67\;MeV & UV brane \\
$u$ & up & 1 & 0.635 & $3.52 \times 10^{-3}$ & 2.16\;MeV & UV brane \\
$e$ & lepton & 1 & 0.656 & $1.71 \times 10^{-3}$ & 0.511\;MeV & UV brane \\
\bottomrule
\end{tabular}
\end{table*}

The localization pattern is physically transparent: first-generation fermions ($c > 0.6$) are UV-localized, far from the IR-brane Higgs, and therefore light. Third-generation fermions ($c \lesssim 0.5$) are IR-localized or flat, overlapping strongly with the Higgs, and therefore heavy. The top quark with $c_{\mathrm{top}} \approx 0$ has maximal Higgs overlap, giving the naturally $\mathcal{O}(1)$ top Yukawa coupling $Y_t \approx 1$.

\begin{theorem}[Warp-Generated Mass Hierarchy]\label{thm:4-warp-hierarchy}
Let the inter-generation mixing be governed by the democratic matrix $M_{\mathrm{oct}}$ with eigenvalues $\{1/2,\, 1/2,\, 2\}$, and let each fermion species $i$ have a bulk mass parameter $c_i \in \mathbb{R}$ on the RS orbifold with $ky_c = 37$. Then:
\begin{enumerate}
  \item[(i)] The effective 4D Yukawa coupling is $Y_i^{\mathrm{eff}} = Y_5\, [g(c_i)]^2$ with $g(c)$ given by \eqref{eq:4-97}.
  \item[(ii)] All nine charged fermion masses are reproduced by $c_i \in [0.004,\, 0.656]$ with universal $Y_5 = 1.00$, spanning a mass hierarchy of $3.4 \times 10^5$ from $\mathcal{O}(1)$ parameters.
  \item[(iii)] The division of labor is clean: the octonionic algebra determines the generational \emph{structure} $(N_g = 3,\; S_3$ symmetry$)$, while the warp factor determines the mass \emph{values} through exponential sensitivity $m_i/m_j \sim e^{2(c_j - c_i)ky_c}$.
\end{enumerate}
\end{theorem}

\begin{proof}
Part~(i) follows from integrating the 5D Dirac equation on the RS background with orbifold boundary conditions~\eqref{eq:4-14} and evaluating the IR-brane Higgs overlap~\cite{ch4:GherghettaPomarol2000}. Part~(ii) is verified numerically: the profile function~\eqref{eq:4-97} evaluated at each $c_i$ in Table~\ref{tab:4-ci} reproduces all observed masses to machine precision. Part~(iii) follows from the factorization $Y_i^{\mathrm{eff}} = Y_5 \cdot M_{\mathrm{oct}} \cdot [g(c_i)]^2$, where $M_{\mathrm{oct}}$ is fixed by the octonionic algebra (zero free parameters) and $g(c_i)$ depends only on the geometry.
\end{proof}

\begin{remark}\label{rmk:4-ci-not-predicted}
The bulk mass parameters $c_i$ are inputs, not predictions. The GP mechanism establishes that $\mathcal{O}(1)$ parameters \emph{can} reproduce the hierarchy---the same qualitative status as Yukawa couplings in the Standard Model---but does not predict the specific values. Whether additional algebraic structure from the octonionic spectral triple constrains the $c_i$ quantitatively is an open question (see Section~\ref{sec:4-open}).
\end{remark}


%% ============================================================
\section{Quark and Lepton Mixing}
\label{sec:4-ckm}

\subsection{CKM Matrix from Democratic Matrix and Warp Profiles}
\label{sec:4-ckm-mechanism}

The CKM matrix arises from the mismatch between up-type and down-type diagonalizations. The physical Yukawa matrix for fermion type~$f$ is the Hadamard product
\begin{equation}
  Y_f = Y_5\, M_{\mathrm{oct}} \circ (\mathbf{g}_f\, \mathbf{g}_f^T),
  \label{eq:4-99}
\end{equation}
where $\mathbf{g}_f = (g(c_{f_1}),\, g(c_{f_2}),\, g(c_{f_3}))^T$ is the vector of profile overlaps. Diagonalizing $Y_u$ and $Y_d$ via $Y_f = V_f\, \Sigma_f\, V_f^T$ (the matrices are real symmetric in the single-$c$ model), the CKM matrix is
\begin{equation}
  V_{\mathrm{CKM}} = V_u^\dagger\, V_d\,.
  \label{eq:4-100}
\end{equation}

The hierarchical profile ratios $g_1 \ll g_2 \ll g_3$ within each sector cause the diagonalizing rotations to be nearly aligned with the generation basis. The CKM matrix is therefore \emph{automatically} near-diagonal---a structural prediction of the democratic $M_{\mathrm{oct}}$ combined with hierarchical warping.

\subsection{The Gatto--Sartori--Tonin Relation}
\label{sec:4-gst}

The GST relation~\cite{ch4:GST1968} emerges as a genuine consequence of the mechanism. For the $(1,2)$ sector, the off-diagonal rotation angle satisfies
\begin{equation}
  |V_{us}| \sim \frac{g_d}{g_s} = \sqrt{\frac{m_d}{m_s}} = 0.224\,,
  \label{eq:4-101}
\end{equation}
compared to the observed value $|V_{us}|_{\mathrm{obs}} = 0.2265$ (a 1.3\% discrepancy). The derivation proceeds from the perturbative diagonalization of~\eqref{eq:4-99}: the $S_3$-symmetric off-diagonal entries of $M_{\mathrm{oct}}$ contribute a common factor that cancels between numerator and denominator, leaving only the profile ratio. A similar analysis gives $|V_{ub}| \sim \sqrt{m_u/m_t} = 0.0035$, matching the observed $0.0036$ to 2.8\%.

The full set of GST-type relations:
\begin{equation}
\begin{aligned}
  |V_{us}| &\sim \sqrt{m_d/m_s} = 0.224 \quad &&(\mathrm{obs.\ }0.227,\; 1.3\%),\\
  |V_{ub}| &\sim \sqrt{m_u/m_t} = 0.0035 \quad &&(\mathrm{obs.\ }0.0036,\; 2.8\%),\\
  |V_{cb}| &\sim \sqrt{m_s/m_b} = 0.149 \quad &&(\mathrm{obs.\ }0.041,\; \text{factor 3.7}).
\end{aligned}
  \label{eq:4-102}
\end{equation}

The $|V_{cb}|$ estimate fails because the $(2,3)$ sector is insufficiently hierarchical for the perturbative expansion. The full left--right treatment resolves this, as we now show.

\subsection{Full Left--Right Fermion Sector}
\label{sec:4-lr-fermions}

The standard RS flavour structure~\cite{ch4:GherghettaPomarol2000,ch4:Casagrande2008,ch4:AgashePerezSoni2005,ch4:Huber2003} assigns separate bulk mass parameters to the $\mathrm{SU}(2)_L$ doublet ($c_{Q_i}$) and singlets ($c_{u_i}$, $c_{d_i}$) of each generation. The effective 4D Yukawa matrices become
\begin{align}
  (Y_u)_{ij} &= Y_5^u\, (M_{\mathrm{oct}})_{ij}\, g(c_{Q_i})\, g(c_{u_j})\,,\label{eq:4-yukawa-LR}\\
  (Y_d)_{ij} &= Y_5^d\, (M_{\mathrm{oct}})_{ij}\, g(c_{Q_i})\, g(c_{d_j})\,.\notag
\end{align}
where $g(c)$ is the profile overlap function~\eqref{eq:4-97} and $Y_5^{u,d}$ are the 5D Yukawa couplings for each sector. The doublet profiles $g(c_{Q_i})$ are shared between up-type and down-type Yukawas---this is the structural origin of CKM mixing in the left--right treatment.

\paragraph{Parameter counting.} The quark sector has 9 bulk masses ($c_{Q_{1,2,3}}$, $c_{u_{1,2,3}}$, $c_{d_{1,2,3}}$) plus 2 Yukawa couplings ($Y_5^u$, $Y_5^d$) $= 11$ parameters for 9 observables (6 quark masses $+$ 3 CKM mixing angles). This is accommodation, not prediction---the same parametric status as the Standard Model (which has 6 Yukawa couplings $+$ 4 CKM parameters $= 10$ free parameters for quarks). The geometric origin of the hierarchy (exponential sensitivity of $g(c)$ to $\mathcal{O}(1)$ parameters) is the structural advance.

\paragraph{Resolution of the $|V_{cb}|$ discrepancy.} Numerical optimisation (differential evolution, $10^5$ function evaluations) yields a best-fit with:
\begin{equation}
\renewcommand{\arraystretch}{1.2}
\begin{array}{@{}lll@{}}
  c_{Q_1} = 0.557, & c_{u_1} = 0.661, & c_{d_1} = 0.495, \\
  c_{Q_2} = 0.646, & c_{u_2} = 0.415, & c_{d_2} = 0.465, \\
  c_{Q_3} = 0.247, & c_{u_3} = 0.200, & c_{d_3} = 0.567,
\end{array}
\label{eq:4-LR-bestfit}
\end{equation}
with $Y_5^u = 1.75$ and $Y_5^d = 0.18$. The resulting CKM elements are:

\begin{table*}[htbp]
\centering
\caption{CKM matrix elements: single-$c$ model, left--right model, and observed values.}
\label{tab:4-ckm-LR}
\small
\begin{tabular}{@{}lcccl@{}}
\toprule
Element & Single-$c$ & Left--Right & Observed & Status \\
\midrule
$|V_{us}|$ & $0.224$ & $0.224$ & $0.224$ & $\checkmark$ \\
$|V_{cb}|$ & $0.149$ & $\mathbf{0.042}$ & $0.042$ & $\checkmark$ \textbf{(fixed)} \\
$|V_{ub}|$ & $0.0035$ & $0.0039$ & $0.0039$ & $\checkmark$ \\
\midrule
$|V_{ud}|$ & --- & $0.975$ & $0.974$ & $\checkmark$ \\
$|V_{cs}|$ & --- & $0.974$ & $0.975$ & $\checkmark$ \\
$|V_{tb}|$ & --- & $0.999$ & $0.999$ & $\checkmark$ \\
\bottomrule
\end{tabular}
\end{table*}

The mechanism is transparent: in the single-$c$ model, $|V_{cb}| \sim \sqrt{m_s/m_b}$ because the doublet and singlet profiles are locked together. In the left--right model, $|V_{cb}|$ is controlled by the doublet profile ratio $g(c_{Q_2})/g(c_{Q_3})$, which is decoupled from the mass ratio $m_s/m_b$ (set by the singlet profiles $g(c_{d_2})/g(c_{d_3})$). The Wolfenstein hierarchy $|V_{us}| : |V_{cb}| : |V_{ub}| \sim \lambda : \lambda^2 : \lambda^3$ with $A = |V_{cb}|/\lambda^2 = 0.84$ is reproduced exactly.

\paragraph{Localization structure.} The best-fit reveals a physically transparent pattern: third-generation doublets ($c_{Q_3} = 0.25$) are IR-localised, giving large top and bottom Yukawas. First- and second-generation doublets ($c_{Q_1} = 0.56$, $c_{Q_2} = 0.65$) are UV-localised, suppressing their Yukawa couplings. The singlet parameters follow the same qualitative pattern as the single-$c$ model (Table~\ref{tab:4-ci}), with the top singlet ($c_{u_3} = 0.20$) most IR-localised and the up singlet ($c_{u_1} = 0.66$) most UV-localised.

\subsection{CP Violation: The Real-Sector Obstruction}
\label{sec:4-cp-real-sector}

All bulk mass parameters $c_i$ are real, $M_{\mathrm{oct}}$ is real, and $g(c)$ is real. Therefore, within the real sector alone:
\begin{equation}
  J_{\mathrm{CP}} = \mathrm{Im}(V_{us}\, V_{cb}\, V_{ub}^*\, V_{cs}^*) = 0\,.
  \label{eq:4-103}
\end{equation}
The observed value $J = (3.08 \pm 0.15) \times 10^{-5}$ requires complex phases in the Yukawa sector.  In Section~\ref{sec:4-cp-mechanism}, we show that these phases enter naturally through the complex brane-localised 5D Yukawa couplings $(Y_5^f)_{ij}$, while the bulk masses $c_i$ and octonionic matrix $M_{\mathrm{oct}}$ remain real.  Two other candidate mechanisms---complex octonionic structure constants and radiative KK loops---are excluded (see Section~\ref{sec:4-cp-mechanism} for details).

\subsection{Quark-Lepton Complementarity}
\label{sec:4-qlc}

The empirical relation
\begin{equation}
  \theta_C + \theta_{12}^{\mathrm{PMNS}} = 13.1^\circ + 33.4^\circ = 46.5^\circ \approx 45^\circ
  \label{eq:4-104}
\end{equation}
has a natural qualitative interpretation in the Meridian framework. The democratic $M_{\mathrm{oct}}$ provides a $45^\circ$ ``maximal mixing'' starting point. The strong quark hierarchy breaks this by $\sim 32^\circ$ (giving $\theta_C = 13^\circ$), while the weaker lepton hierarchy plus seesaw enhancement breaks it by only $\sim 12^\circ$ (giving $\theta_{12} = 33^\circ$). The sum recovers the democratic baseline. This explanation is qualitative; the specific values depend on the bulk mass parameters.


%% ============================================================
\section{Majorana Sector and the Seesaw Mechanism}
\label{sec:4-majorana}

\subsection{$S_3$-Constrained Majorana Mass Matrix}
\label{sec:4-s3-majorana}

The $S_3$ generation symmetry derived from octonionic triality (Section~\ref{sec:4-octonionic}) constrains the $3 \times 3$ Majorana mass matrix $M_R$ for right-handed neutrinos. Since $S_3$ acts by simultaneous permutation of rows and columns, the most general $S_3$-invariant real symmetric matrix has only two free parameters:
\begin{equation}
  M_R = a\, \mathbf{I}_3 + b\, (\mathbf{J}_3 - \mathbf{I}_3) = (a - b)\, \mathbf{I}_3 + b\, \mathbf{J}_3\,,
  \label{eq:4-105}
\end{equation}
where $\mathbf{J}_3$ is the $3 \times 3$ all-ones matrix. This represents a reduction from six to two free parameters---a 67\% reduction enforced by the octonionic algebra.

The eigenvalues and eigenvectors are:
\begin{equation}
\begin{aligned}
  M_d &= a - b \quad &(\text{doublet, degeneracy 2}),\\
  M_s &= a + 2b \quad &(\text{singlet, degeneracy 1}),
\end{aligned}
  \label{eq:4-106}
\end{equation}
with eigenvectors forming the tribimaximal (TBM) basis:
\begin{equation}
  \mathbf{v}_1 = \frac{1}{\sqrt{2}}\begin{pmatrix} 1\\-1\\0 \end{pmatrix}, \quad
  \mathbf{v}_2 = \frac{1}{\sqrt{6}}\begin{pmatrix} 1\\1\\-2 \end{pmatrix}, \quad
  \mathbf{v}_3 = \frac{1}{\sqrt{3}}\begin{pmatrix} 1\\1\\1 \end{pmatrix}.
  \label{eq:4-107}
\end{equation}

\subsection{Tribimaximal Mixing as Leading-Order Prediction}
\label{sec:4-tbm}

\begin{theorem}[Tribimaximal Mixing from $S_3$]\label{thm:4-tbm}
In the $S_3$-symmetric limit where all neutrino bulk mass parameters are equal ($c_{\nu_1} = c_{\nu_2} = c_{\nu_3} = c_0$), the Dirac matrix $m_D = m_{D0}\, M_{\mathrm{oct}}$ and the Majorana matrix $M_R$ share the TBM eigenbasis~\eqref{eq:4-107}. The Type-I seesaw formula $m_\nu = -m_D\, M_R^{-1}\, m_D^T$ then yields:
\begin{equation}
  U_{\mathrm{PMNS}} = U_{\mathrm{TBM}} = \begin{pmatrix}
    \sqrt{2/3} & 1/\sqrt{3} & 0 \\
    -1/\sqrt{6} & 1/\sqrt{3} & -1/\sqrt{2} \\
    -1/\sqrt{6} & 1/\sqrt{3} & 1/\sqrt{2}
  \end{pmatrix},
  \label{eq:4-108}
\end{equation}
predicting
\begin{equation}
  \sin^2\theta_{12} = \tfrac{1}{3} = 0.333, \qquad
  \sin^2\theta_{23} = \tfrac{1}{2} = 0.500, \qquad
  \sin^2\theta_{13} = 0\,.
  \label{eq:4-109}
\end{equation}
\end{theorem}

\begin{proof}
In the $S_3$-symmetric limit, $m_D$ and $M_R$ are both elements of the two-dimensional algebra $\mathrm{span}\{\mathbf{I}_3,\, \mathbf{J}_3\}$, which is closed under multiplication and inversion. Therefore $m_\nu = -m_D\, M_R^{-1}\, m_D^T$ also lies in this algebra, and shares the TBM eigenvectors~\eqref{eq:4-107}. Since the charged lepton mass matrix is likewise proportional to $M_{\mathrm{oct}}$ (same eigenvectors), the PMNS matrix $U = V_e^\dagger\, V_\nu$ is determined by the rotation from the charged lepton mass eigenstate basis to the TBM basis. By explicit computation using the $M_{\mathrm{oct}}$ eigenvalues $\{1/2, 1/2, 2\}$, the rotation matrix is $U_{\mathrm{TBM}}$.
\end{proof}

Comparison with NuFIT~5.2 (normal ordering)~\cite{ch4:NuFIT2024}: $\sin^2\theta_{12} = 0.307 \pm 0.013$ ($2.0\sigma$ from TBM), $\sin^2\theta_{23} = 0.546 \pm 0.021$ ($2.2\sigma$), $\sin^2\theta_{13} = 0.0220 \pm 0.0007$ ($31\sigma$ from zero). The first two angles are approximated well; $\theta_{13} = 0$ is excluded by Daya Bay at $> 30\sigma$~\cite{ch4:DayaBay2012}. $S_3$-breaking from different neutrino bulk mass parameters $c_{\nu_i}$ is mandatory.

\subsection{Neutrino Mass Accommodation}
\label{sec:4-nu-mass}

With $S_3$-broken Majorana parameters (three $M_R$ eigenvalues + three $c_{\nu_i}$, six parameters for six observables), the framework accommodates all neutrino data:
\begin{equation}
\begin{aligned}
  m_1 &= 1.27 \times 10^{-3}\;\mathrm{eV},\\
  m_2 &= 9.69 \times 10^{-3}\;\mathrm{eV},\\
  m_3 &= 4.08 \times 10^{-2}\;\mathrm{eV},\\
  \textstyle\sum m_\nu &= 0.052\;\mathrm{eV} \;\;(\text{below } 0.12\;\mathrm{eV}\text{ bound}),
\end{aligned}
  \label{eq:4-110}
\end{equation}
with normal hierarchy preferred by the $S_3$ singlet being naturally heavier than the doublet. The effective Majorana mass for neutrinoless double-beta decay is
\begin{equation}
  m_{ee} = (3\text{--}5) \times 10^{-3}\;\mathrm{eV} \qquad \text{(tribimaximal limit)}\,,
  \label{eq:4-111}
\end{equation}
which reduces to $m_{ee} \sim 1.5$--$3.7\;\mathrm{meV}$ when realistic $S_3$-breaking ($\theta_{13} \neq 0$) is included (see Section~\ref{sec:4-experiment-forecast}). Both ranges fall below next-generation sensitivity (nEXO: $\sim 5\text{--}18\;\mathrm{meV}$; LEGEND-1000: $\sim 9\text{--}21\;\mathrm{meV}$), consistent with normal ordering. A null result in $0\nu\beta\beta$ searches is a concrete negative prediction of the framework.


\subsection{\texorpdfstring{$S_3$}{S3} Breaking and the Octonionic Mass Matrix}
\label{sec:4-s3-breaking}

The $S_3$ generation symmetry, derived from the three independent complex structures on the octonions (Theorem~\ref{thm:4-3}), governs the inter-generation mixing through $M_{\mathrm{oct}}$ (Theorem~\ref{thm:4-4}).  A critical clarification: \emph{$M_{\mathrm{oct}}$ governs the Dirac mass matrix, not the Majorana mass matrix.}  The Majorana sector is controlled by the Type-I seesaw mechanism with parameters determined by brane physics, not by the octonionic overlap.

The eigenvalues of $M_{\mathrm{oct}}$ are $\{1/2,\, 1/2,\, 2\}$ (equation~\eqref{eq:4-83}), verified under all six $S_3$ permutations.  The doublet degeneracy (two eigenvalues at $1/2$) is the $S_3$ representation-theoretic origin of the near-degeneracy $M_{R_2} \approx M_{R_3}$, which drives ARS leptogenesis (Section~\ref{sec:4-baryogenesis}).  The Fano plane overlap between any two complex structures is $1/2$---a geometric constant determined by the quaternionic subalgebra intersections.

\paragraph{Parameter count: $6 \to 6$.}
The $S_3$ symmetry alone does \emph{not} reduce the number of free parameters in the neutrino sector.  The six parameters (three bulk masses $c_{\nu_i}$ and three Majorana parameters $a$, $b$, $\Delta M$) remain independent.  This is a negative result.  What $S_3$ constrains is \emph{qualitative}, not quantitative:
\begin{itemize}
  \item Mass hierarchy topology: $M_1 \ll M_2 \approx M_3$ (doublet $+$ singlet);
  \item ARS leptogenesis viability (near-degeneracy is structural, not tuned);
  \item Fano plane overlap $= 1/2$ (geometric, from quaternionic subalgebra intersections).
\end{itemize}
The individual bulk masses $c_{\nu_i}$, the $S_3$-breaking splitting $\delta c$ (ARS-viable range: $10^{-9}$ to $10^{-7}$), and the absolute neutrino masses remain undetermined by the algebra.

The $S_3$ symmetry from octonions provides no quantitative predictive advantage over the standard $\nu$MSM for the neutrino sector; its value is structural (explaining \emph{why} the near-degeneracy pattern exists) rather than parametric.

\paragraph{Comparison with Gresnigt.}
Gresnigt's $S_3$ from $\mathrm{Aut}(S_{16})$ and Meridian's $S_3$ from $\mathrm{Aut}(\mathbb{O}) = G_2$ permute the same three objects (three 8-dimensional representations mapped to three generations).  The groups are related: $G_2$ contains the $S_3$ that acts on the complex structures, and since $\mathrm{Aut}(\mathbb{O})$ embeds in $\mathrm{Aut}(S_{16})$, Meridian's $S_3$ is a subgroup of Gresnigt's larger automorphism group.  Same permutation symmetry ($S_3$ on three objects); whether the Fano-plane overlap matrix $M_\mathrm{oct}$ agrees between the two derivations remains an open comparison.


\subsection{Conditional Neutrino Predictions}
\label{sec:4-conditional-nu}

The RS warp-profile mechanism constrains the Dirac mass matrix: $m_D = f_L \cdot Y_5 \cdot f_R$, where $f_L$ and $f_R$ are warp-factor overlaps determined by bulk masses.  This eliminates the Casas--Ibarra freedom present in generic Type-I seesaw.  The number of predictions depends critically on the 5D Yukawa coupling matrix $Y_5$:

\begin{table*}[htbp]
\centering
\caption{Neutrino prediction count as a function of $Y_5$ structure.}
\label{tab:4-nu-predictions}
\small
\begin{tabular}{@{}lcccc@{}}
\toprule
Scenario & $N_{\mathrm{params}}$ & $N_{\mathrm{obs}}$ & Net predictions & Status \\
\midrule
$S_3$ only (baseline) & 6 & 6 & 0 & Negative \\
$Y_5$ fixed by geometry & 4 & 8 & 4 & Conditional \\
$Y_5$ texture (2 params) & 6 & 8 & 2 & Conditional \\
$Y_5$ anarchic (5 params) & 9 & 8 & $-1$ & No prediction \\
\bottomrule
\end{tabular}
\end{table*}

Two structural exclusions hold regardless of $Y_5$ determination:
\begin{itemize}
  \item \textbf{Diagonal $Y_5$ is excluded:} it produces zero mixing angles.
  \item \textbf{Universal $Y_5$ is excluded:} it gives rank-1 $m_D$ (only one non-zero neutrino mass).
\end{itemize}

\textbf{If $Y_5$ is geometric} (i.e., determined by the NCG vacuum selection or AS fixed point; see Section~\ref{sec:4-open}), four quantities become predictions: (i)~the Dirac CP phase $\delta_{\mathrm{CP}}$, (ii)~the lightest neutrino mass $m_1$, (iii)~the effective Majorana mass $m_{ee}$, and (iv)~the Majorana phases.  A CSD(3) texture, if selected, predicts $\delta_{\mathrm{CP}} = -87^\circ$, separated from the NuFIT best-fit ($-163^\circ$) by $5.1\sigma$ at DUNE.  This would be the most discriminating near-term test of the framework's neutrino sector.


\subsection{The $\nu$MSM Embedding}
\label{sec:4-numsm}

The three right-handed neutrinos required by the spectral triple (Section~\ref{sec:4-brane-triple}) embed naturally into the $\nu$MSM framework of Asaka, Blanchet, and Shaposhnikov~\cite{ch4:ABS2005}:
\begin{itemize}
  \item $\nu_{R_1}$ ($M_{R_1} \sim \mathrm{keV}$): dark matter candidate (Section~\ref{sec:4-dm}).
  \item $\nu_{R_2},\, \nu_{R_3}$ ($M_{R_{2,3}} \sim \mathrm{GeV}$): generate light neutrino masses via seesaw; their CP-violating decays produce the lepton asymmetry for baryogenesis and Shi--Fuller production of $\nu_{R_1}$.
\end{itemize}
This mass pattern requires maximal $S_3$-breaking in $M_R$, achieved through $\mathcal{O}(1)$ differences in the neutrino bulk mass parameters: $c_{\nu_1} \sim 1.17$ (UV-localized, giving $M_{R_1} \sim 7\;\mathrm{keV}$) versus $c_{\nu_{2,3}} \sim 1.0$ (moderately UV-localized, giving $M_{R_{2,3}} \sim \mathrm{GeV}$). The GP mechanism converts the $\mathcal{O}(0.17)$ difference in $c_\nu$ into the required six orders of magnitude in $M_R$.

\subsection{Baryogenesis via Leptogenesis}
\label{sec:4-baryogenesis}

The CP-violating decays and oscillations of $\nu_{R_2}$ and $\nu_{R_3}$ generate the baryon asymmetry via the ARS mechanism~\cite{ch4:ARS1998,ch4:AsakaShaposhnikov2005}. Three structural features of the Meridian framework align with this scenario:

\paragraph{1. Automatic near-degeneracy.} In the standard $\nu$MSM, the mass degeneracy $M_{R_2} \approx M_{R_3}$ required for the ARS mechanism must be imposed by hand. In the Meridian framework, it is a structural consequence of the $S_3$ doublet: both $\nu_{R_2}$ and $\nu_{R_3}$ carry mass $M_d = a - b$ at leading order. The splitting $\Delta M = M_{R_3} - M_{R_2}$ is generated by $S_3$-breaking corrections to the bulk mass parameters, controlled by $\Delta c_\nu = c_{\nu_3} - c_{\nu_2}$.

\paragraph{2. O(1) CP violation.} The CP-violating invariant relevant for leptogenesis is
\begin{equation}
  I_{\mathrm{CP}} = \frac{\mathrm{Im}\bigl[(F^\dagger F)_{23}^2\bigr]}{(F^\dagger F)_{22}\, (F^\dagger F)_{33}}\,,
  \label{eq:4-ICP}
\end{equation}
where $F = m_D / v$ is the Yukawa coupling matrix. Using the Casas--Ibarra parametrization with the Meridian neutrino masses~\eqref{eq:4-110} and a complex angle $\omega_{23} = \pi/4 + 0.5\, i$ (corresponding to 13\% Yukawa amplification), we find
\begin{equation}
  I_{\mathrm{CP}} = -0.61\,.
  \label{eq:4-ICP-value}
\end{equation}
This is naturally $\mathcal{O}(1)$ --- the CP violation from the brane Yukawa phases (Section~\ref{sec:4-cp-mechanism}) is fully sufficient for leptogenesis. No additional CP sources are required.

\paragraph{3. Baryon asymmetry magnitude.} In the weak-washout regime ($K_{\mathrm{prod}} \lesssim 1$), the ARS asymmetry is
\begin{equation}
  \eta_B \sim \frac{28}{79}\, \frac{135\, \zeta(3)}{4\pi^4 g_*}\, K_{\mathrm{prod}}^2\, I_{\mathrm{CP}}\, \sin\!\bigl(\phi_{\mathrm{sph}}\bigr)\,,
  \label{eq:4-etaB-ARS}
\end{equation}
where $K_{\mathrm{prod}} = \Gamma_{\mathrm{prod}}/H|_{T_{\mathrm{sph}}}$ and $\phi_{\mathrm{sph}} = \Delta M^2 M_{\mathrm{Pl}} / (4 \times 3.15 \times 1.66\, \sqrt{g_*}\, T_{\mathrm{sph}}^3)$ is the oscillation phase at sphaleron freeze-out. For the Meridian parameters, the envelope of $|\eta_B|$ reaches $\sim 7 \times 10^{-4}$, \emph{overshooting} the observed $\eta_B^{\mathrm{obs}} = 6.14 \times 10^{-10}$ by six orders of magnitude. The observed value is achieved when $\sin(\phi_{\mathrm{sph}}) \sim 10^{-6}$, which occurs at multiple values of $\Delta M$ in the range $0.01$--$1\;\mathrm{keV}$.

\begin{proposition}[Baryogenesis Viability]\label{prop:4-BAU}
The Meridian $\nu$MSM embedding produces $\eta_B = (6.14 \pm 0.02) \times 10^{-10}$ for $\Delta M / M \in [10^{-8},\, 10^{-6}]$, corresponding to $\Delta c_\nu \in [10^{-10},\, 10^{-8}]$ via the GP mechanism. The $S_3$ doublet provides the required near-degeneracy automatically; the precise $\Delta M$ constrains the $S_3$-breaking scale from UV physics.
\end{proposition}

This is an accommodation, not a prediction: $\Delta M$ and $\mathrm{Im}(\omega_{23})$ are adjusted to match $\eta_B^{\mathrm{obs}}$. However, the framework reduces the standard $\nu$MSM parameter count from 11 (three $M_R$ eigenvalues plus eight Yukawa parameters) to 3 ($a$, $b$, and $\Delta c_\nu$), and provides a structural explanation for the most unnatural feature of the $\nu$MSM --- the near-degeneracy of $M_{R_2}$ and $M_{R_3}$.


%% ============================================================
\section{Radiative Stability of the Democratic Matrix}
\label{sec:4-radiative-stability}

A necessary consistency check is whether the democratic structure of $M_{\mathrm{oct}}$ survives renormalization group evolution from the Kaluza--Klein scale $M_{\mathrm{KK}} \sim 3\;\mathrm{TeV}$ to the electroweak scale $m_Z = 91.2\;\mathrm{GeV}$.

\begin{proposition}[$S_3$ Preservation under SM RGE]\label{prop:4-s3-rge}
If $Y_f(M_{\mathrm{KK}}) = y_f\, M_{\mathrm{oct}}$ for $f \in \{u, d, e\}$, then $Y_f(\mu) = y_f(\mu)\, M_{\mathrm{oct}}(\mu)$ where $M_{\mathrm{oct}}(\mu)$ remains $S_3$-symmetric at all scales $\mu < M_{\mathrm{KK}}$ under the 1-loop Standard Model RGE.
\end{proposition}

\begin{proof}
The SM Yukawa RGE has the structure $16\pi^2\, dY_f/dt = Y_f(\alpha_f\, Y_f^\dagger Y_f + \beta_f\, Y_{f'}^\dagger Y_{f'}) + c_f\, Y_f$, where $c_f = T - G_f$ is a scalar (trace plus gauge factors). For $Y_f = y_f\, M_{\mathrm{oct}}$:
\begin{equation}
  Y_f^\dagger Y_f = y_f^2\, M_{\mathrm{oct}}^2, \qquad
  Y_f\, Y_f^\dagger Y_f = y_f^3\, M_{\mathrm{oct}}^3\,.
  \label{eq:4-112}
\end{equation}
By the Cayley--Hamilton theorem, $M_{\mathrm{oct}}$ has minimal polynomial of degree~2 (two distinct eigenvalues: $1/2$ with multiplicity~2 and $2$ with multiplicity~1), so $M_{\mathrm{oct}}^3$ lies in the span of $\{\mathbf{I}_3,\, M_{\mathrm{oct}}\}$. The right-hand side of the RGE is therefore a linear combination of $M_{\mathrm{oct}}$ and $\mathbf{I}_3$, both of which are $S_3$-invariant. The flow remains in the $S_3$-symmetric subspace at all scales.
\end{proof}

Numerical integration confirms this: the off-diagonal spread of $Y_f$ is zero to machine precision ($10^{-16}$) at all intermediate scales. The eigenvalue ratio $\lambda_{\min}/\lambda_{\max}$ shifts by at most 2.3\% (up-type quarks) from its initial value of $1/4$, while the overall Yukawa magnitude changes by 17--23\% (quarks) and 2.4\% (leptons).

\begin{corollary}\label{cor:4-no-ckm-from-rge}
No CKM or PMNS mixing is generated by RGE evolution from a democratic starting point. All physical mixing must originate from the bulk mass parameters $c_i$.
\end{corollary}

The argument extends to all perturbative orders: any polynomial in $M_{\mathrm{oct}}$ lies in $\mathrm{span}\{\mathbf{I}_3, M_{\mathrm{oct}}\}$, so higher-loop Yukawa insertions $Y(Y^\dagger Y)^n$ cannot escape the $S_3$-symmetric subspace. The democratic $M_{\mathrm{oct}}$ is a \emph{quasi-fixed point} of the SM RGE: its $S_3$ symmetry and eigenvector structure are exactly preserved, with eigenvalue ratios drifting by $< 3\%$.


%% ============================================================
\section{Dark Matter}
\label{sec:4-dm}

\subsection{Systematic Candidate Survey}
\label{sec:4-dm-survey}

The Meridian framework contains a finite set of potential dark matter candidates. We evaluate each systematically:

\begin{table*}[htbp]
\centering
\caption{Dark matter candidate survey in the Meridian framework.}
\label{tab:4-dm}
\small
\begin{tabular}{@{}llll@{}}
\toprule
Candidate & Stabilizing symmetry? & Lifetime & Status \\
\midrule
Lightest KK particle & No (KK parity broken by warping) & $\tau \ll t_{\mathrm{univ}}$ & \textbf{Excluded} \\
Radion & No (trace anomaly coupling) & $\tau \ll t_{\mathrm{univ}}$ & \textbf{Excluded} \\
Cuscuton excitation & No (zero mode = DE; KK unstable) & --- & \textbf{Excluded} \\
Sterile neutrino $\nu_{R_1}$ & Approximate (small mixing) & $\tau \gg t_{\mathrm{univ}}$ & \textbf{Viable} \\
\bottomrule
\end{tabular}
\end{table*}

\paragraph{LKP exclusion.} In the Randall--Sundrum orbifold, KK parity is broken: the warp factor $A(y) = -ky$ distinguishes the UV brane ($y = 0$) from the IR brane ($y = y_c$), so the transformation $y \to y_c - y$ is \emph{not} a symmetry of the metric. All KK excitations decay promptly. Even if hypothetically stable, the KK graviton's spin-independent cross-section $\sigma_{\mathrm{SI}} \sim m_N^4/(\pi\, \Lambda_{\mathrm{IR}}^4) \sim 8 \times 10^{-44}\;\mathrm{cm}^2$ exceeds the LZ bound~\cite{ch4:LZ2024} by two orders of magnitude.

\paragraph{Radion exclusion.} The radion couples universally through the trace anomaly $\mathcal{L}_{\mathrm{int}} = -(r/\Lambda_r)\, T^\mu{}_\mu$. The dominant decay $r \to gg$ gives $\Gamma \sim \alpha_s^2\, m_r^3\, b_3^2/(32\pi^3\, \Lambda_r^2)$, with lifetime $\tau \ll t_{\mathrm{univ}}$ for all $m_r > 30\;\mathrm{meV}$. No discrete symmetry forbids this decay.

\paragraph{Cuscuton exclusion.} The cuscuton has $Q_s = P_X + 2X\, P_{XX} = 0$ at leading order---zero propagating degrees of freedom. The $\epsilon_1\, X$ Gauss--Bonnet correction introduces a propagating mode, but its zero mode \emph{is} the dark energy field ($m \sim H_0$) and its KK excitations are unstable (broken KK parity).

\subsection{Sterile Neutrino: The Unique Viable Candidate}
\label{sec:4-dm-sterile}

\begin{theorem}[Unique Dark Matter Candidate]\label{thm:4-dm}
Within the particle content of the Meridian framework (spectral triple on the RS orbifold with cuscuton self-tuning), the lightest right-handed neutrino $\nu_{R_1}$ is the unique viable cold/warm dark matter candidate. The three exclusions above are structural: broken KK parity eliminates all KK modes, the trace anomaly eliminates the radion, and the cuscuton constraint eliminates the bulk scalar.
\end{theorem}

\begin{proof}
The spectral triple requires $\nu_R$ as part of the 16-dimensional representation per generation (Section~\ref{sec:4-brane-triple}). Being a gauge singlet, $\nu_{R_1}$ interacts only through Yukawa mixing. For $c_{\nu_1} \gtrsim 1.19$ (an $\mathcal{O}(1)$ bulk mass parameter), the GP mechanism gives $M_{R_1} \sim 7\;\mathrm{keV}$ and active-sterile mixing $\sin^2 2\theta \lesssim 2 \times 10^{-11}$, consistent with the XRISM non-detection of the 3.5~keV line~\cite{ch4:XRISM2024}. The lifetime $\tau_{\nu_{R_1}} > 10^{28}\;\mathrm{s} \gg t_{\mathrm{univ}}$ ensures DM stability. The Shi--Fuller resonant production mechanism~\cite{ch4:ShiFuller1999}, driven by lepton asymmetry from the heavier $\nu_{R_{2,3}}$ decays (Section~\ref{sec:4-baryogenesis}), yields $\Omega_s h^2 \simeq 0.12$ for $\sin^2 2\theta \gtrsim 10^{-13}$. The exclusions of the other three candidates are structural consequences of the RS warping (KK parity), the trace anomaly (radion), and the cuscuton constraint (bulk scalar), independent of parameter choices.
\end{proof}

The sterile neutrino DM prediction is testable via X-ray spectroscopy: the decay $\nu_{R_1} \to \gamma + \nu_{\mathrm{active}}$ produces a monochromatic line at $E = m_s/2 = 3.5\;\mathrm{keV}$.

\subsection{Detection Strategy and Experimental Constraints}
\label{sec:4-dm-detection}

\paragraph{X-ray line.} The radiative decay rate is~\cite{ch4:PalWolfenstein1982}
\begin{multline}
  \Gamma(\nu_s \!\to\! \nu_a + \gamma) = \frac{9\, \alpha\, G_F^2}{1024\, \pi^4}\, \sin^2 2\theta\, m_s^5
  \\
  \simeq 1.5 \!\times\! 10^{-28}\, \biggl(\frac{\sin^2 2\theta}{10^{-11}}\biggr) \!\biggl(\frac{m_s}{7\;\mathrm{keV}}\biggr)^{\!5}\mathrm{s}^{-1},
  \label{eq:4-Gamma-rad}
\end{multline}
producing a monochromatic line at $E = m_s/2 = 3.5\;\mathrm{keV}$ from dark matter halos.

\paragraph{Current constraints.} The XRISM Resolve spectrometer set $\sin^2 2\theta < 2.4 \times 10^{-11}$ at $m_s = 7.1\;\mathrm{keV}$ (99.7\% CL) from a 230~ks observation of the Perseus cluster~\cite{ch4:XRISM2024}. The Meridian framework accommodates this constraint with $c_{\nu_1} \gtrsim 1.19$, a 1.3\% shift from the na\"ive benchmark $c_{\nu_1} = 1.17$. The mixing angle depends exponentially on $c_{\nu_1}$ through the GP overlap factor $g(c) \propto \exp[-(c - \tfrac{1}{2})\, k y_c]$, so even this small shift reduces $\sin^2 2\theta$ from $7 \times 10^{-11}$ to below $2 \times 10^{-11}$.

\paragraph{Viable parameter space.} The Shi--Fuller production mechanism requires $\sin^2 2\theta \gtrsim 10^{-13}$ to achieve the correct relic abundance $\Omega_s h^2 \simeq 0.12$. Combined with the XRISM upper bound, the viable range is
\begin{equation}
  10^{-13} \lesssim \sin^2 2\theta \lesssim 2.4 \times 10^{-11}\,,
  \label{eq:4-viable-mixing}
\end{equation}
corresponding to $c_{\nu_1} \in [1.19,\, 1.25]$.

\paragraph{Future experiments.} The Athena X-ray Observatory (projected launch $\sim$2035) will achieve sensitivity $\sin^2 2\theta \sim 3 \times 10^{-13}$ with a deep Perseus observation, covering \emph{the entire viable Meridian parameter space}~\eqref{eq:4-viable-mixing}. Non-detection by Athena would exclude $m_s = 7\;\mathrm{keV}$ sterile neutrino DM within the Shi--Fuller framework.

\paragraph{Structure formation.} Lyman-$\alpha$ forest constraints require $m_s \gtrsim 3.3\;\mathrm{keV}$ for thermal warm DM (relaxed to $\sim 2\;\mathrm{keV}$ for resonant Shi--Fuller production)~\cite{ch4:Viel2013}. At $m_s = 7\;\mathrm{keV}$, the free-streaming length $\lambda_{\mathrm{fs}} \approx 0.05\;\mathrm{Mpc}$ is well below the smallest observed structures. The MW satellite count constraint $m_s \gtrsim 6.5\;\mathrm{keV}$ (95\% CL)~\cite{ch4:Nadler2021} is marginally satisfied.


%% ============================================================
\section{NCG--Asymptotic Safety Bridge}
\label{sec:4-ncg-as}

\subsection{$R^2 = 0$ from Conformal Structure}
\label{sec:4-r2-zero}

The $a_4$ Seeley--DeWitt coefficient for the Dirac operator $D^2 = -\nabla^2 + R/4$ (Lichnerowicz formula) in the $(C^2,\, \mathcal{E}_4,\, R^2)$ basis takes the form~\cite{ch4:Vassilevich2003,ch4:Gilkey1984}:
\begin{equation}
  a_4 \propto d_S\, \bigl[-18\, C^2 + 11\, \mathcal{E}_4 + 0 \cdot R^2\bigr]\,,
  \label{eq:4-113}
\end{equation}
where $d_S$ is the spinor dimension and $C^2 = R_{MNPQ}^2 - 2\, R_{MN}^2 + \tfrac{1}{3}\, R^2$ is the Weyl tensor squared. The $R^2$ coefficient vanishes due to the algebraic identity
\begin{equation}
  \tfrac{5}{4} - \tfrac{2}{3} - \tfrac{7}{12} = \tfrac{15 - 8 - 7}{12} = 0\,,
  \label{eq:4-114}
\end{equation}
which holds for all spinor dimensions $d_S$ and all spacetime dimensions~$d$. This is a \emph{structural} property of the Dirac operator, arising from the conformal covariance of the Dirac equation: the trace anomaly for a Dirac fermion contains $C^2$ and $\mathcal{E}_4$ but no $R^2$ term, because $R^2$ would break conformal structure.

\subsection{Warped Orbifold Preserves $R^2 = 0$}
\label{sec:4-r2-warped}

On the 5D RS orbifold, two additional contributions could in principle modify this result:

\paragraph{Bulk contribution.} The 5D curvature invariants decompose into 4D curvature plus warp terms. The quadratic terms $5\, \hat{R}^2 - 8\, \hat{R}_{MN}^2 - 7\, \hat{R}_{MNPQ}^2$ acquire a factor $e^{4ky}$ from the 5D-to-4D curvature conversion, which is exactly cancelled by the inverse factor $e^{-4ky}$ in the 5D volume element $\sqrt{g_5} = e^{-4ky}\, \sqrt{\hat{g}}$. The warp factor drops out, and the effective 4D coefficients retain the ratio $(-18,\, +11,\, 0)$.

\paragraph{Boundary contribution.} On the orbifold, the 5D spinor has 2~components with Neumann BCs and 2~with Dirichlet BCs. The boundary $a_4$ terms are proportional to the Robin parameter $\chi$: $+1$ for Neumann, $-1$ for Dirichlet. The sum $(+1)(2) + (-1)(2) = 0$ ensures exact cancellation of boundary curvature-squared corrections.

\subsection{The Spectral Action on the Critical Surface}
\label{sec:4-critical-surface}

\begin{theorem}[Spectral Action Critical Surface]\label{thm:4-critical}
The spectral action on the RS orbifold sits exactly on the codimension-1 critical surface of the Reuter fixed point of asymptotic safety. Specifically:
\begin{enumerate}
  \item[(i)] The spectral action direction in coupling space is $(\omega,\, \sigma) \propto (-18,\, 0)$, where $\omega$ is the $C^2$ coupling and $\sigma$ the $R^2$ coupling.
  \item[(ii)] The Reuter fixed point has $\sigma$ as the UV-repulsive direction (critical exponent $\theta_\sigma = -2.11$).
  \item[(iii)] The spectral action has zero projection onto the UV-repulsive direction: $\sigma = 0$ exactly.
  \item[(iv)] One-loop graviton corrections generate $\sigma > 0$~\cite{ch4:AvramidiBarvinsky1985,ch4:FradkinTseytlin1982}, pushing the theory \emph{into} the basin of attraction.
\end{enumerate}
\end{theorem}

\begin{proof}
Part~(i) follows from~\eqref{eq:4-113}. Part~(ii) is established in the AS literature~\cite{ch4:Benedetti2010,ch4:Codello2009}. Part~(iii) is the content of~\eqref{eq:4-114}. Part~(iv) is verified explicitly in Proposition~\ref{prop:4-sigma1}: the one-loop $R^2$ counterterm from the graviton (Lichnerowicz) and ghost (vector Laplacian) sectors yields $\sigma_1 = +0.403 > 0$, consistent with the Avramidi--Barvinsky result~\cite{ch4:AvramidiBarvinsky1985}.
\end{proof}

The physical interpretation is that $\sigma = 0$ means the scalaron ($R^2$ scalar mode) is \emph{non-propagating} at tree level---its mass $m_{\mathrm{scalaron}}^2 \sim 1/(6\sigma)$ is infinite. This is consistent with the Meridian framework, where the dynamical scalar degree of freedom is the bulk cuscuton/radion, not the scalaron. The scalaron acquires a finite mass at one loop, as computed explicitly in the next subsection.

\subsection{One-Loop Verification: Explicit Computation of \texorpdfstring{$\sigma_1$}{sigma1}}
\label{sec:4-oneloop-r2}

We now verify part~(iv) of Theorem~\ref{thm:4-critical} by explicit computation. The one-loop effective action from quantum gravitational fluctuations around a fixed background receives contributions from two sectors:

\paragraph{Graviton sector.} The Lichnerowicz operator $\Delta_L$ acts on symmetric 2-tensors ($N = d(d{+}1)/2 = 10$ components in $d = 4$). Its endomorphism involves the Riemann tensor (cf.~Barvinsky--Vilkovisky 1985): $E^{\mu\nu}{}_{\rho\sigma} = -2\, R^{(\mu}{}_{(\rho}\,{}^{\nu)}{}_{\sigma)} + R^{(\mu}{}_{(\rho}\,\delta^{\nu)}{}_{\sigma)} + \tfrac{1}{2}\,R\,\delta^{(\mu}{}_{(\rho}\,\delta^{\nu)}{}_{\sigma)}$.

\paragraph{Ghost sector.} The vector Laplacian ($N = d = 4$) has endomorphism $E_{\mu\nu} = R_{\mu\nu}$ and connection curvature $\Omega_{\mu\nu}{}^{\rho\sigma} = R^{\rho}{}_{\mu\nu}{}^{\sigma}$.

The Seeley--DeWitt $a_4$ coefficient for an operator $-\nabla^2 + E$ is~\cite{ch4:Vassilevich2003}:
\begin{equation}
  b_4 = \frac{1}{360}\bigl[N(5R^2 - 2R_{\mu\nu}^2 + 2R_{\mu\nu\rho\sigma}^2) + 60R\,\mathrm{tr}\, E + 180\,\mathrm{tr}\, E^2 + 30\,\mathrm{tr}\,\Omega^2\bigr]\,.
  \label{eq:4-b4-general}
\end{equation}
Evaluating on three backgrounds with linearly independent curvature invariants ($S^4$, $S^2 \times S^2$ with two curvature scale ratios) and solving the resulting $3 \times 3$ system yields:
\begin{align}
  b_4^{\mathrm{ghost}} &= \frac{1}{360}\bigl[80\, R^2 + 172\, R_{\mu\nu}^2 + 38\, R_{\mu\nu\rho\sigma}^2\bigr]\,,
  \label{eq:4-b4-ghost} \\
  b_4^{\mathrm{grav}} &= \frac{1}{360}\bigl[-130\, R^2 + 1780\, R_{\mu\nu}^2 + 380\, R_{\mu\nu\rho\sigma}^2\bigr]\,.
  \label{eq:4-b4-grav}
\end{align}
The ghost coefficients $(80,\, 172,\, 38)$ match the known analytical result~\cite{ch4:AvramidiBarvinsky1985}, cross-validating the computation. The one-loop divergence $\Gamma_1^{\mathrm{div}} = \tfrac{1}{2}\, b_4^{\mathrm{grav}} - b_4^{\mathrm{ghost}}$ in the $(C^2,\, \mathcal{E}_4,\, R^2)$ basis is:
\begin{equation}
  \Gamma_1^{\mathrm{div}} = \frac{1}{16\pi^2 \varepsilon}\!\int\! \sqrt{g}\;
  \bigl[\underbrace{1.842}_{\sigma_{C^2}} C^2
  - \underbrace{1.419}_{\sigma_{E_4}} \mathcal{E}_4
  + \underbrace{0.403}_{\sigma_1} R^2\bigr].
  \label{eq:4-oneloop-r2}
\end{equation}

\begin{proposition}[One-Loop $R^2$ Coefficient]\label{prop:4-sigma1}
The one-loop $R^2$ coefficient is $\sigma_1 = +0.403 > 0$. This confirms Theorem~\ref{thm:4-critical}(iv): quantum gravitational fluctuations generate a \emph{positive} $R^2$ coupling, pushing the theory into the basin of attraction of the Reuter fixed point.
\end{proposition}

Three features of this result merit emphasis:

\begin{enumerate}
  \item \textbf{Perturbative smallness.} The ratio $\sigma_1 / |\sigma_{C^2}^{\mathrm{tree}}| = 0.403 / 18 \approx 2.2\%$, confirming perturbative control.
  \item \textbf{Mass independence.} The $R^2$ coefficient in~$b_4$ depends only on the curvature of the endomorphism~$E$ and connection~$\Omega$, not on the mass. A mass shift $E \to E + m^2 \cdot \mathbf{1}$ modifies only the $R$ and constant terms in~$b_4$, leaving the curvature-squared structure invariant. Consequently, \emph{each KK mode contributes the same $\sigma_1$}, and the total KK tower sum $N_{\mathrm{KK}} \cdot \sigma_1$ (regulated by the 5D cutoff) does not change the sign.
  \item \textbf{Structural origin of $R^2 = 0$ breaking.} The tree-level identity $\alpha + \beta/3 + \gamma/3 = 0$ is specific to the Dirac operator. The ghost ($\alpha + \beta/3 + \gamma/3 = 0.417$) and graviton ($1.639$) operators both violate it, as they are not conformally structured. The breaking is expected and irreducible.
\end{enumerate}


\subsection{Quantitative RG Flow Analysis}
\label{sec:4-rg-flow}

We now strengthen Theorem~\ref{thm:4-critical} with a quantitative analysis of the RG flow from the Reuter fixed point to the spectral action cutoff scale~$\Lambda \sim 10^{17}\;\mathrm{GeV}$.

\paragraph{Eigenvector alignment.}
The stability matrix $M_{ij} = \partial \beta_i / \partial g_j |_{\mathrm{FP}}$ of the four-derivative truncation~\cite{ch4:Benedetti2010} has eigenvalues $\lambda_\alpha = -\theta_\alpha$ corresponding to the four critical exponents. The right eigenvector~$V_4$ associated with the UV-irrelevant direction ($\theta_4 = -2.11$) determines the normal to the critical surface at the fixed point. Perturbative estimates of the off-diagonal stability matrix elements $\partial \beta_\sigma / \partial g \sim 0.34$ (with non-perturbative enhancement $\sim 4\times$; see Codello--Percacci--Rahmede~\cite{ch4:Codello2009}, Eq.~4.18), combined with the eigenvalue gap $|\lambda_4 - \lambda_{\mathrm{EH}}| \approx 4.6$, give the mixing of~$V_4$ into the Einstein--Hilbert sector:
\begin{equation}
  |V_4^{(g)}| \sim 0.07, \quad |V_4^{(\lambda)}| \sim 0.04, \quad |V_4^{(\omega)}| \sim 0, \quad |V_4^{(\sigma)}| \approx 1.
  \label{eq:4-v4-alignment}
\end{equation}
After normalisation, $|V_4^{(\sigma)}|^2 \approx 0.993$: the UV-irrelevant eigenvector is 99\%~aligned with the~$\sigma$ direction. Consequently, $\sigma = 0$ is effectively the critical surface condition.

\paragraph{Projection analysis.}
The deviation $\delta g_i = g_i^{\mathrm{SA}} - g_i^*$ from the NGFP to the spectral action couplings at $k = \Lambda$ is dominated by the~$C^2$ direction ($|\delta\omega| \gg |\delta g|, |\delta\lambda|, |\delta\sigma|$). The projection onto the UV-irrelevant eigenvector satisfies
\begin{equation}
  \frac{|W_4 \cdot \delta g|}{|\delta g|} \approx 0.006\,,
  \label{eq:4-projection}
\end{equation}
i.e., the spectral action is within 0.6\% of the critical surface in coupling space.

\paragraph{$C^2$ trajectory consistency.}
The~$C^2$ coupling is strongly UV-relevant ($\theta_3 = 8.40$). On the critical surface, the linearised flow amplifies deviations by $\exp(-\theta_3\, t_{\mathrm{SA}}) \sim 10^{17}$, where $t_{\mathrm{SA}} = \ln(\Lambda / M_{\mathrm{Pl}}) \approx -4.8$. The spectral action value $\omega_{\mathrm{SA}} \approx -10.9$ (for the Standard Model NCG, $d_S = 96$) requires an initial condition $C_3 \sim 10^{-17}$ at the Planck scale---exponentially small, as expected for a relevant direction that is only slightly perturbed in the UV. This is a generic feature of asymptotic safety, not a fine-tuning: the large amplification is precisely the mechanism by which the NGFP generates predictive low-energy physics from UV-universal initial conditions.

\paragraph{$R^2$ on the critical surface.}
Since $C_4 = 0$ on the critical surface, the~$R^2$ coupling receives no contribution from the UV-irrelevant direction along the entire RG trajectory. It evolves only through the small mixing terms from $V_{1,2,3}$, remaining close to~$\sigma^*$ throughout the flow. The spectral action value $\sigma_{\mathrm{SA}} = 0$ is therefore consistent with the trajectory: $\sigma$ has evolved from~$\sigma^* = -0.010$ at the NGFP to $\sigma \approx 0$ at $k = \Lambda$ through the small mixing contributions alone.

\begin{remark}
The linearised analysis is strictly valid only near the NGFP. The spectral action scale $\Lambda \sim 10^{17}\;\mathrm{GeV}$ is two orders of magnitude below $M_{\mathrm{Pl}}$, so the couplings have departed from the fixed-point values (particularly $g(\Lambda) \sim 10^{-5} \ll g^* \approx 0.86$). However, the structural content---$\sigma = 0$ gives zero projection onto the UV-irrelevant direction---is exact and independent of the linearisation. A full nonlinear trajectory computation (requiring numerical integration of the FRG flow equations from the NGFP through the crossover regime) would quantify the subleading corrections but cannot invalidate the tree-level placement.
\end{remark}


%% ============================================================
\section{Inflation from Modulus Geometry}
\label{sec:4-inflation}

\subsection{$R^2 = 0$ Kills Starobinsky Inflation}
\label{sec:4-starobinsky-killed}

The Starobinsky action $S = (M_{\mathrm{Pl}}^2/2)\, R + \alpha_R\, R^2$ requires $\alpha_R \neq 0$. Theorem~\ref{thm:4-critical} establishes $\alpha_R = 0$ as a structural identity of the spectral action. \emph{Starobinsky inflation is therefore excluded in the Meridian framework.}

\textbf{Observational caveat.} The modulus inflation mechanism (Section~\ref{sec:4-modulus-inflation}) produces an $\alpha = 1$ attractor potential formally identical to the Starobinsky potential~\eqref{eq:4-117}. The inflationary observables ($n_s$, $r$) are therefore indistinguishable from Starobinsky $R^2$ inflation for at least the next decade: LiteBIRD + Simons Observatory can confirm $r \approx 0.004$ at $\sim 7\sigma$ but cannot distinguish the two mechanisms. The discriminant lies in the reheating sector---modulus decay to $WW/ZZ$ with BR $> 85\%$ via the trace anomaly coupling vs.\ democratic Starobinsky decay---which affects $N_*$ by $\sim 1$--$2$ e-folds and $\Delta r \sim 2 \times 10^{-5}$, below near-term sensitivity. The $R^2 = 0$ identity is thus a theoretical prediction whose observational consequences are indirect on decade timescales.

The conformal plateau mechanism of Higgs inflation~\cite{ch4:BezrukovShaposhnikov2008} also fails: for $V_J = (\lambda/4)\, r^4$ with $\xi = 1/6$, the Einstein-frame potential $V_E = V_J/\Omega^4$ \emph{diverges} as $r \to \sqrt{6}\, M_{\mathrm{Pl}}$ rather than flattening. The plateau requires $\xi \gg 1$; at $\xi = 1/6$ the conformal stretching is too weak.

\subsection{Modulus Inflation via K\"ahler Geometry}
\label{sec:4-modulus-inflation}

The resolution comes from the RS modulus $T = e^{ky_c}$, which parameterizes the size of the extra dimension. At energies above the Goldberger--Wise stabilization scale, $T$ is an approximately flat direction with a K\"ahler potential
\begin{equation}
  K = -3\, M_{\mathrm{Pl}}^2\, \ln(T + T^*)\,,
  \label{eq:4-115}
\end{equation}
defining the field-space metric $G_{TT^*} = 3\, M_{\mathrm{Pl}}^2/(T + T^*)^2$ on the coset $\mathrm{SL}(2,\mathbb{R})/\mathrm{U}(1)$, a space of constant negative curvature.

\emph{Note:} The K\"ahler structure~\eqref{eq:4-115} is an additional assumption beyond A1--A6. It follows naturally from $\mathcal{N} = 1$ supergravity descriptions of the RS modulus sector~\cite{ch4:Luty1999}, but the Meridian framework is not supersymmetric. We adopt~\eqref{eq:4-115} as the simplest field-space metric consistent with the $\mathrm{SL}(2,\mathbb{R})$ symmetry of the modulus; alternative K\"ahler potentials (e.g., with higher-order corrections) would modify the inflationary predictions quantitatively but not the qualitative $\alpha$-attractor structure, which depends only on the sign of $\mathcal{R}_K$.
\begin{equation}
  \mathcal{R}_K = -\frac{2}{3}\,.
  \label{eq:4-116}
\end{equation}
The canonical field $\sigma = \sqrt{3/2}\, M_{\mathrm{Pl}}\, \ln(T/T_0)$ converts a potential of the form $V(T) = V_0\, [1 - c/T + \cdots]$ into
\begin{equation}
  V(\sigma) = V_0\, \Bigl[1 - e^{-\sqrt{2/3}\;\sigma/M_{\mathrm{Pl}}}\Bigr]^2,
  \label{eq:4-117}
\end{equation}
which is formally identical to the Starobinsky potential. In the Kallosh--Linde $\alpha$-attractor classification~\cite{ch4:KalloshLinde2013}, this corresponds to
\begin{equation}
  \alpha_K = -\frac{3}{\mathcal{R}_K} = 1\,.
  \label{eq:4-118}
\end{equation}

\subsection{The Seventh $1/6$: Inflationary Consistency}
\label{sec:4-seventh-one-sixth}

The non-minimal coupling $\alpha$-attractor parameter is $\alpha_\xi = 1/(6\xi)$. At $\xi = 1/6$:
\begin{equation}
  \alpha_\xi = \frac{1}{6 \cdot 1/6} = 1 = \alpha_K\,.
  \label{eq:4-119}
\end{equation}
This agreement holds \emph{only} at $\xi = 1/6$---for any other value of $\xi$, the K\"ahler and NMC $\alpha$-parameters disagree. We note that $\alpha_K = 1$ depends on the assumed K\"ahler potential~\eqref{eq:4-115} (an additional assumption; see the caveat in Section~\ref{sec:4-modulus-inflation}), while $\alpha_\xi = 1$ follows from $\xi = 1/6$ (derived within the framework). The coincidence is therefore a consistency check between an assumption and a derivation---it gains its force from the fact that $\alpha_K = 1$ was not engineered to match $\alpha_\xi$, but follows from the standard RS modulus K\"ahler geometry. This is one of four such checks that complement the three independent derivations established in Section~\ref{sec:4-xi-derivation}:
\begin{enumerate}
  \item Seeley--DeWitt $a_2$ coefficient (algebraic),
  \item Radion as metric fluctuation (geometric),
  \item Weyl invariance of spectral action (functional),
  \item Self-tuning exact to 15 significant figures (physical),
  \item $\det(\mathbf{Z}) = 1$ / no ghost (kinetic stability),
  \item AS anomalous dimension $\eta_m = 0$ (RG protection),
  \item $\alpha_K = \alpha_\xi = 1$ (inflationary consistency).
\end{enumerate}

\subsection{Predictions}
\label{sec:4-inflation-predictions}

The slow-roll parameters for the $\alpha = 1$ attractor potential~\eqref{eq:4-117} are:
\begin{equation}
  n_s = 1 - \frac{2}{N_*} + \mathcal{O}(N_*^{-2}), \qquad
  r = \frac{12}{N_*^2}\,,
  \label{eq:4-120}
\end{equation}
where $N_*$ is the number of $e$-folds before the end of inflation.

\begin{theorem}[Inflationary Observables]\label{thm:4-inflation}
The modulus inflation mechanism in the Meridian framework predicts:
\begin{equation}
  n_s = 0.965 \pm 0.003, \qquad r = 0.004 \pm 0.001\,,
  \label{eq:4-121}
\end{equation}
for $N_* = 50$--$60$ (determined by the reheating temperature $T_{\mathrm{reh}} \sim 10^{8}$--$10^{10}\;\mathrm{GeV}$). The tensor-to-scalar ratio $r \sim 0.004$ is detectable by LiteBIRD $(\sigma(r) \sim 10^{-3})$ at $\sim 4\sigma$ significance, and definitively testable by the combination of LiteBIRD with the Simons Observatory $(\sigma(r) \sim 5 \times 10^{-4})$ at $\sim 7\sigma$; see Section~\ref{sec:4-bmode-forecast}.
\end{theorem}

\begin{table*}[htbp]
\centering
\caption{Inflationary predictions for selected $e$-fold numbers.}
\label{tab:4-inflation}
\small
\begin{tabular}{@{}ccccl@{}}
\toprule
$N_*$ & $n_s$ & $r$ & $\phi_*/M_{\mathrm{Pl}}$ & Planck tension \\
\midrule
50 & 0.9582 & 0.0049 & 5.14 & $1.6\sigma$ \\
55 & 0.9621 & 0.0041 & 5.26 & $0.7\sigma$ \\
57 & 0.9635 & 0.0038 & 5.30 & $0.3\sigma$ \\
60 & 0.9654 & 0.0034 & 5.37 & $0.1\sigma$ \\
\bottomrule
\end{tabular}
\end{table*}

The predictions are observationally identical to Starobinsky $R^2$ inflation---the distinguishing features lie in the reheating sector (the modulus decays through the trace anomaly coupling to $WW/ZZ$, dominant BR $> 85\%$, rather than the democratic Starobinsky decay) and in the absence of the scalaron degree of freedom.

\begin{figure*}[htbp]
\centering
\includegraphics[width=0.85\textwidth]{figures/fig4_inflation_nsr.pdf}
\caption{Meridian inflationary predictions in the $n_s$--$r$ plane. The $\alpha = 1$
K\"ahler modulus attractor (blue line, $N_* = 50$--$60$) predicts $n_s = 0.965 \pm 0.003$
and $r = 0.004 \pm 0.001$, consistent with Planck~2018 at $0.3\sigma$. The QQG framework
(red shaded, $r \geq 0.01$) provides a sharp observational discriminant: Simons Observatory
(${\sim}2027$, $\sigma(r) \sim 0.003$) and CMB-S4 (${\sim}2029$, $\sigma(r) \sim 0.001$)
will cleanly separate the two predictions. LiteBIRD sensitivity shown as dotted line.}
\label{fig:4-inflation}
\end{figure*}

\subsection{Primordial B-Mode Forecast and Experimental Tests}
\label{sec:4-bmode-forecast}

The tensor-to-scalar ratio $r = 0.004 \pm 0.001$ lies an order of magnitude below the current best limit $r < 0.036$ (95\% CL, BICEP/Keck BK18~\cite{ch4:BK18}). Testing this prediction requires next-generation B-mode polarization experiments.

\paragraph{LiteBIRD.} The JAXA-led LiteBIRD satellite~\cite{ch4:LiteBIRD2023}, scheduled for launch in JFY~2032, will observe the full sky from Sun--Earth $L_2$ for three years using 4{,}508 TES bolometers across 15 frequency bands (34--448~GHz). The design requirement is $\delta r < 10^{-3}$, with statistical sensitivity $\sigma_{\mathrm{stat}}(r) = 6 \times 10^{-4}$ and a systematic budget $\sigma_{\mathrm{syst}}(r) = 5.7 \times 10^{-4}$. For Meridian's $r = 0.004$:
\begin{equation}
  \mathrm{SNR}_{\mathrm{LiteBIRD}} = \frac{r}{\sigma(r)} \approx \frac{0.004}{10^{-3}} = 4.0\,,
  \label{eq:4-SNR-LB}
\end{equation}
constituting strong evidence but falling short of the conventional $5\sigma$ discovery threshold.

\paragraph{Simons Observatory.} The ground-based Simons Observatory~\cite{ch4:SO2019}, currently observing from the Atacama Desert, will reach $\sigma(r) \approx 1.2 \times 10^{-3}$ with six small-aperture telescopes, giving SNR $\approx 3$ independently.

\paragraph{Combined detection.} Since LiteBIRD (full-sky, large scales) and the Simons Observatory (deep patch, intermediate scales) observe largely independent sky regions, their Fisher matrices add:
\begin{equation}
  \sigma(r)_{\mathrm{LB+SO}} = \Bigl(\sigma_{\mathrm{LB}}^{-2} + \sigma_{\mathrm{SO}}^{-2}\Bigr)^{-1/2} \approx 5.4 \times 10^{-4}\,,
  \label{eq:4-sigma-combined}
\end{equation}
yielding $\mathrm{SNR} \approx 7$ for $r = 0.004$---a definitive detection.

\begin{table*}[htbp]
\centering
\caption{Detection forecast for the Meridian prediction $r = 0.004$.}
\label{tab:4-bmode-forecast}
\small
\begin{tabular}{@{}llccc@{}}
\toprule
Experiment & Timeline & $\sigma(r)$ & SNR & Status \\
\midrule
BICEP/Keck BK18 & Current & $9 \times 10^{-3}$ & 0.4 & Operating \\
Simons Observatory (6 SATs) & 2028 & $1.2 \times 10^{-3}$ & 3.1 & Observing \\
LiteBIRD & 2033--36 & $10^{-3}$ & 4.0 & Approved \\
LiteBIRD + SO & $\sim$2037 & $5.4 \times 10^{-4}$ & 6.8 & Combined \\
CMB-S4 & --- & $5 \times 10^{-4}$ & 7.4 & Cancelled (2025) \\
\bottomrule
\end{tabular}
\end{table*}

\paragraph{CMB-S4 cancellation.} The DOE/NSF CMB-S4 project~\cite{ch4:CMBS4} was cancelled in July 2025 due to infrastructure limitations at the South Pole site. Its projected $\sigma(r) = 5 \times 10^{-4}$ would have provided the most sensitive single-experiment test of $r = 0.004$. The cancellation makes LiteBIRD the world's definitive B-mode experiment for the coming two decades.

\paragraph{Falsifiability.} The prediction is sharply falsifiable:
\begin{itemize}
  \item \emph{Detection} of $r \approx 0.004$ confirms the $\alpha = 1$ universality class shared with Starobinsky $R^2$ inflation. The reheating sector (modulus decay to $WW/ZZ$ with BR $> 85\%$ vs.\ democratic Starobinsky decay) and the radion (Section~\ref{sec:4-collider}) provide the discriminants.
  \item \emph{Non-detection} ($r < 0.002$ at 95\% CL) excludes $\alpha = 1$ attractors at $> 4\sigma$, requiring $\alpha < 0.5$ and ruling out the K\"ahler modulus mechanism.
  \item Combined LiteBIRD + SO makes non-detection of $r = 0.004$ an exclusion at $\sim 7\sigma$---definitive either way.
\end{itemize}

\subsection{Reheating and $N_*$ Determination}
\label{sec:4-reheating}

The modulus oscillates about the Goldberger--Wise minimum after inflation ends, decaying to Standard Model particles through the trace anomaly coupling $\mathcal{L}_{\mathrm{int}} = (\sigma/\Lambda_r)\, T^\mu{}_\mu$.  The dominant channels for $m_\sigma > 2m_W$ are $WW + ZZ > 85\%$, with total width
\begin{equation}
  \Gamma_{\mathrm{total}} \sim \frac{m_\sigma^3}{32\pi\, \Lambda_r^2}\,,
  \label{eq:4-Gamma-reh}
\end{equation}
yielding a reheating temperature
\begin{equation}
  T_{\mathrm{reh}} = \biggl(\frac{90}{\pi^2\, g_*}\biggr)^{1/4}\!\sqrt{\Gamma\, M_{\mathrm{Pl}}}
  \;\sim\; 5 \times 10^7\text{--}3 \times 10^9\;\mathrm{GeV}
  \label{eq:4-Treh}
\end{equation}
for the natural modulus mass range $m_\sigma \in [200,\, 2000]\;\mathrm{GeV}$.  This determines the $e$-fold number:
\begin{equation}
  N_* \approx 53\text{--}56 \qquad \Longrightarrow \qquad
  n_s = 0.9634 \pm 0.0009, \quad r = 0.0040 \pm 0.0002\,,
  \label{eq:4-Nstar-reh}
\end{equation}
tightening the generic prediction ($N_* \in [50,\, 60]$) by a factor of four.  The reheating temperature is comfortably above the BBN bound ($T_{\mathrm{reh}} \gg 1\;\mathrm{MeV}$) and above the GeV-scale sterile neutrino mass required for ARS leptogenesis (Section~\ref{sec:4-baryogenesis}).

The trace-anomaly dominance of the modulus decay ($WW/ZZ > 85\%$) is structurally distinct from the democratic Starobinsky decay (proportional to $m^2$, with $f\bar{f} \sim 70\%$).  While this difference shifts $N_*$ by only $\sim 1$--$2$ and $\Delta r \sim 2 \times 10^{-5}$ (below LiteBIRD sensitivity), it provides a theoretical discriminant between the two inflation mechanisms.


\subsection{Gravitational Wave Signatures}
\label{sec:4-gw-signatures}

If the RS radion triggers a first-order confinement--deconfinement phase transition (Hawking--Page type)---as argued in~\cite{ch4:Creminelli2002,ch4:Nardini2007} but not computed within this framework---the transition temperature $T_c \sim k\, e^{-ky_c}/8 \approx 190\;\mathrm{GeV}$ would generate a stochastic gravitational wave background.

The dominant contribution is from sound waves in the plasma~\cite{ch4:Hindmarsh2017}, with peak frequency and amplitude
\begin{align}
  f_{\mathrm{peak}} &\approx 0.3\;\mathrm{mHz}\;\biggl(\frac{\beta/H}{10}\biggr)\biggl(\frac{T_c}{200\;\mathrm{GeV}}\biggr),\label{eq:4-gw-pt}\\
  \Omega_{\mathrm{GW}}\, h^2 &\sim 10^{-7}\;\biggl(\frac{\alpha}{1}\biggr)^2\biggl(\frac{10}{\beta/H}\biggr),\notag
\end{align}
where $\alpha$ is the ratio of latent heat to radiation energy and $\beta/H$ is the nucleation rate.  For the RS phase transition, $\alpha \sim \mathcal{O}(1)$ and $\beta/H \sim \mathcal{O}(10$--$100)$ (from the Goldberger--Wise stabilisation potential).

\paragraph{Detection prospects.} The peak frequency $\sim 0.3\;\mathrm{mHz}$ lies squarely in the LISA sensitivity band~\cite{ch4:LISA2017}.  For $\alpha \geq 0.1$, the signal exceeds LISA's power-law sensitivity ($\Omega_{\mathrm{min}} \sim 10^{-13}$) by several orders of magnitude:
\begin{equation}
  \frac{\Omega_{\mathrm{signal}}(f_{\mathrm{LISA}})}{\Omega_{\mathrm{noise}}(f_{\mathrm{LISA}})} \gg 1 \qquad (\text{for } \alpha \gtrsim 0.1,\; \beta/H \lesssim 100)\,.
  \label{eq:4-gw-snr}
\end{equation}
BBO and DECIGO would detect the signal in the deci-Hz band.  The inflationary gravitational wave background ($\Omega_{\mathrm{GW}}^{\mathrm{inf}} \sim r\, A_s\, \Omega_{\mathrm{rad}}/24 \sim 3 \times 10^{-17}$ from $r = 0.004$) is far below all planned detector sensitivities and is accessible only through CMB B-mode polarization (Section~\ref{sec:4-bmode-forecast}).

This constitutes the framework's third independent detection channel: B-mode polarization (inflation), LISA (phase transition), and collider resonances (radion and KK tower)---probing three distinct energy scales with three distinct experimental techniques.


%% ============================================================
\section{Collider Signatures}
\label{sec:4-collider}

\subsection{Higgs--Radion Mixing at $\xi = 1/6$}
\label{sec:4-higgs-radion}

The Giudice--Rattazzi--Wells (GRW) formalism~\cite{ch4:GRW2001} parameterizes the Higgs--radion kinetic and mass mixing through the dimensionless ratio
\begin{equation}
  \gamma = \frac{v_{\mathrm{EW}}}{\Lambda_r}\,, \qquad
  \Lambda_r = \sqrt{6}\, M_{\mathrm{Pl}}\, e^{-ky_c}\,,
  \label{eq:4-122}
\end{equation}
where $\Lambda_r$ is the radion coupling scale. The physical mixing angle is~\cite{ch4:GRW2001}
\begin{equation}
  \tan 2\theta = \frac{12\xi\, \gamma\, m_h^2}{m_r^2(1 + 6\xi\gamma^2) - m_h^2}\,.
  \label{eq:4-123}
\end{equation}
For $\xi = 0$ (the AS prediction for generic scalars): $\theta = 0$, no mixing. For $\xi = 1/6$ (Meridian): $\theta > 0$, proportional to $\gamma$. This is the core observable distinction.

\subsection{The Conformal $\det(\mathbf{Z}) = 1$ Property}
\label{sec:4-det-z}

At $\xi = 1/6$, the kinetic mixing matrix determinant satisfies
\begin{equation}
  \det(\mathbf{Z}) = 1 + 6\xi\gamma^2(1 - 6\xi) = 1\,,
  \label{eq:4-124}
\end{equation}
\emph{exactly}, for all $\gamma$. This means: (i)~no ghost modes ($\det(\mathbf{Z}) > 0$ always), (ii)~no eigenvalue pathology, (iii)~the mixing is purely mass-mixing, not kinetic-mixing. This is a structural consequence of conformal symmetry with no analogue at $\xi = 0$.

\subsection{Coupling Modifiers and the $d > c$ Diagnostic}
\label{sec:4-coupling-modifiers}

The Higgs-like mass eigenstate $h_1 = \cos\theta\, h - \sin\theta\, r$ has modified couplings~\cite{ch4:GRW2001}:
\begin{align}
  d &= \cos\theta + \gamma\, \sin\theta \quad (\text{coupling to } VV),
  \label{eq:4-125}\\
  c &= \cos\theta \quad (\text{coupling to } f\bar{f})\,.
  \notag
\end{align}
The distinctive prediction is $d > c$: the radion admixture \emph{enhances} the $VV$ coupling (through the trace anomaly) relative to fermion couplings. This pattern---$\kappa_V > 1$, $\kappa_f < 1$---is the diagnostic signature of the conformal non-minimal coupling.

For standard RS1 parameters ($k \sim M_{\mathrm{Pl}}$, $ky_c = 35$, $\Lambda_r = 3.76\;\mathrm{TeV}$):
\begin{equation}
  \gamma = 0.065, \qquad \theta = 0.78^\circ \;\;(m_r = 300\;\mathrm{GeV}), \qquad
  \kappa_V - 1 = 8.0 \times 10^{-4}\,.
  \label{eq:4-126}
\end{equation}
This deviation is below the sensitivity of any planned collider (HL-LHC: $\sigma(\kappa_V) \sim 0.017$; FCC-ee: $\sim 0.002$; muon collider: $\sim 0.001$). The primary experimental path is therefore not precision Higgs measurements but \emph{direct radion discovery} as a diboson resonance, followed by coupling-ratio measurements that determine $\xi$.

\paragraph{Predicted radion mass.} The quantum-corrected radion mass from the NCG spectral action is $m_r \approx 120\;\mathrm{GeV}$ (see Chapter~\ref{ch:foundation}, Section~\ref{sec:1-radion}), dominated by one-loop curvature-squared contributions ($99.7\%$ of the mass). The proximity to $m_h = 125.1\;\mathrm{GeV}$ enhances the Higgs--radion mixing: at $m_r = 120\;\mathrm{GeV}$, the mixing angle from Eq.~\eqref{eq:4-123} increases to
\begin{equation}
  \theta(m_r = 120\;\mathrm{GeV}) \approx 18.5^\circ\,, \qquad
  \kappa_V - 1 \approx 0.021\,,
  \label{eq:4-126b}
\end{equation}
which enters the sensitivity range of HL-LHC ($\sigma(\kappa_V) \sim 0.017$). The near-degeneracy $m_r \approx m_h$ is a distinctive prediction of the framework, as the NCG quantum correction and the Higgs mechanism both operate at the electroweak scale through the common origin in the spectral action. The illustrative examples in Table~\ref{tab:4-collider} use $m_r = 300\;\mathrm{GeV}$ for comparison with the existing RS literature; the $m_r \approx 120\;\mathrm{GeV}$ prediction improves the collider prospects considerably.

\subsection{Sensitivity Assessment}
\label{sec:4-sensitivity}

\begin{table*}[htbp]
\centering
\caption{Sensitivity of collider facilities to the $\xi = 1/6$ prediction, for $m_r = 300\;\mathrm{GeV}$ at standard RS1 parameters.}
\label{tab:4-collider}
\small
\begin{tabular}{@{}lccl@{}}
\toprule
Collider & $\sigma(\kappa_V)$ & Signal/$\sigma$ & Can distinguish $\xi = 1/6$ from $0$? \\
\midrule
LHC Run~2 & 0.050 & 0.02 & No \\
HL-LHC & 0.017 & 0.05 & No \\
FCC-hh & 0.004 & 0.20 & No \\
FCC-ee & 0.002 & 0.40 & No \\
Muon Collider & 0.001 & 0.80 & Marginal \\
\bottomrule
\end{tabular}
\end{table*}

If $\Lambda_r < 500\;\mathrm{GeV}$ (possible in modified RS setups with $k < M_{\mathrm{Pl}}$), the deviation $\kappa_V - 1$ reaches $\sim 1\%$, detectable at HL-LHC. The strongest theoretical discriminator remains the self-tuning mechanism: any $\xi \neq 1/6$ destroys the cosmological constant cancellation to 120 orders of magnitude (Section~\ref{sec:4-xi-derivation}).

\begin{remark}\label{rmk:4-falsification}
The $\xi = 1/6$ prediction is falsifiable in principle at any collider precision. If a radion-like resonance is discovered and its coupling ratio $d/c = 1 + \gamma\tan\theta$ is measured to deviate from the $\xi = 1/6$ prediction at $> 3\sigma$, the Meridian framework is falsified. Until radion discovery, the strongest constraints come from cosmological probes of the dark energy equation of state (DESI) and the internal consistency of the spectral action.
\end{remark}

\subsection{Radion Discovery Strategy}
\label{sec:4-radion-discovery}

Direct radion discovery as a narrow diboson resonance is the primary experimental path to testing the $\xi = 1/6$ prediction.

\paragraph{Decay pattern.} The radion decays predominantly through the QCD and electroweak trace anomalies.  For $m_r > 2m_W$, the dominant channels are $WW + ZZ > 85\%$, with the branching ratio pattern:
\begin{align}
  \mathrm{BR}(r \to WW) &\approx 0.63, \notag\\
  \mathrm{BR}(r \to ZZ) &\approx 0.28, \notag\\
  \mathrm{BR}(r \to gg) &\approx 0.07\,,
  \label{eq:4-radion-BR}
\end{align}
at $m_r = 300\;\mathrm{GeV}$.  Above the $t\bar{t}$ threshold, the top channel opens ($\mathrm{BR} \sim 20$--$30\%$) but $WW + ZZ$ remains dominant.  The total width $\Gamma/m_r < 5 \times 10^{-4}$---the radion is an extremely narrow resonance at all masses.

\paragraph{Discovery reach.} At standard RS1 parameters ($\Lambda_r = 3.76\;\mathrm{TeV}$), the production cross-section $\sigma(gg \to r) \propto 1/\Lambda_r^2$ is suppressed by a factor $(v_{\mathrm{EW}}/\Lambda_r)^2 \sim 4 \times 10^{-3}$ relative to a SM-like heavy scalar.  Current ATLAS/CMS diboson resonance searches are insensitive at this coupling scale.

\begin{table*}[htbp]
\centering
\caption{Radion discovery reach at current and future colliders ($\Lambda_r = 3.76\;\mathrm{TeV}$).}
\label{tab:4-radion-reach}
\small
\begin{tabular}{@{}llcl@{}}
\toprule
Facility & Luminosity & $m_r$ reach ($5\sigma$) & Channel \\
\midrule
LHC Run~2 & 139~fb$^{-1}$ & --- & Unconstrained \\
HL-LHC & 3~ab$^{-1}$ & $\lesssim 300\;\mathrm{GeV}$ (if $\Lambda_r < 1\;\mathrm{TeV}$) & $WW \to \ell\nu jj$ \\
FCC-hh & 30~ab$^{-1}$ & $\sim 1\;\mathrm{TeV}$ & $WW/ZZ$ \\
FCC-ee & 150~ab$^{-1}$ & --- (precision $\kappa_V$ to $0.2\%$) & $ZH$ \\
Muon Collider & 10~ab$^{-1}$ at $10\;\mathrm{TeV}$ & $\sim 5\;\mathrm{TeV}$ & $s$-channel \\
\bottomrule
\end{tabular}
\end{table*}

FCC-hh provides the first realistic discovery opportunity at standard RS1 parameters.  A muon collider at $\sqrt{s} = 10\;\mathrm{TeV}$ could produce the radion in the $s$-channel and measure its coupling ratios with sufficient precision to test $d > c$.

\paragraph{Discrimination from heavy Higgs.} Four observables distinguish the radion from a heavy Higgs boson:
\begin{enumerate}
  \item \emph{Coupling ratio} $d/c > 1$: the radion enhances $VV$ couplings (trace anomaly), while a heavy Higgs reduces both $VV$ and $f\bar{f}$ couplings.
  \item \emph{Width}: the radion is narrow ($\Gamma/m < 0.1\%$); MSSM heavy Higgs bosons can be broad ($\Gamma/m \sim \text{few}\%$) at large $\tan\beta$.
  \item \emph{Spin}: spin-0 (scalar), distinguishable from spin-2 KK graviton via angular distributions in $ZZ \to 4\ell$.
  \item \emph{Branching ratios}: $WW + ZZ > 85\%$ for the radion (trace anomaly dominance) vs.\ $b\bar{b}$-dominant for heavy Higgs below the $t\bar{t}$ threshold.
\end{enumerate}

\paragraph{Radion vs.\ KK graviton.} The first KK graviton has mass $m_1 = x_1\, k\, e^{-ky_c} \approx 5.9\;\mathrm{TeV}$ for $k \sim M_{\mathrm{Pl}}$, $ky_c = 35$ (where $x_1 = 3.832$ is the first zero of $J_1$), above current exclusion limits ($m_1 > 4.3\;\mathrm{TeV}$ for $k/M_{\mathrm{Pl}} = 0.1$).  The radion, with mass set by brane coupling strengths, is generically lighter.  It would be the \emph{first} collider signature of the extra dimension.

\subsection{The 120 GeV Radion: HL-LHC Testability}
\label{sec:4-radion-120}

The NCG quantum correction to the radion mass (Chapter~\ref{ch:foundation}, Section~\ref{sec:1-radion}) places the radion at $m_r \approx 120\;\mathrm{GeV}$---below the Higgs mass $m_h = 125.1\;\mathrm{GeV}$ and within 5~GeV of it.  This near-degeneracy, a structural consequence of the spectral action operating at the electroweak scale, transforms the collider phenomenology relative to the $m_r = 300\;\mathrm{GeV}$ benchmark used in the preceding sections.

\paragraph{Enhanced mixing.}  At $m_r = 120\;\mathrm{GeV}$, the denominator of Eq.~\eqref{eq:4-123} approaches zero, enhancing the mixing angle from $\theta = 0.78^\circ$ ($m_r = 300\;\mathrm{GeV}$) to $\theta \approx 18.5^\circ$---a factor of 24 increase.  The corresponding coupling modifier becomes
\begin{equation}
  \kappa_V - 1 \approx 0.021\,,
  \label{eq:4-kappa-120}
\end{equation}
which enters the HL-LHC sensitivity range ($\sigma(\kappa_V) \sim 0.017$) for the first time.

\paragraph{Updated sensitivity.}  Table~\ref{tab:4-collider-120} presents the sensitivity at $m_r = 120\;\mathrm{GeV}$; the improvement over Table~\ref{tab:4-collider} ($m_r = 300\;\mathrm{GeV}$) is dramatic.

\begin{table*}[htbp]
\centering
\caption{Sensitivity to $\kappa_V - 1 = 0.021$ at $m_r = 120\;\mathrm{GeV}$ (standard RS1 parameters).}
\label{tab:4-collider-120}
\small
\begin{tabular}{@{}lccl@{}}
\toprule
Collider & $\sigma(\kappa_V)$ & Signal/$\sigma$ & Status \\
\midrule
HL-LHC & 0.017 & 1.2 & Marginal ($1.2\,\sigma$) \\
FCC-hh & 0.004 & 5.3 & Discovery ($5\,\sigma$) \\
FCC-ee & 0.002 & 10.5 & Decisive \\
Muon Collider & 0.001 & 21 & Decisive \\
\bottomrule
\end{tabular}
\end{table*}

\paragraph{Decay pattern at $m_r = 120\;\mathrm{GeV}$.}  Below the $WW$ threshold ($2m_W \approx 161\;\mathrm{GeV}$), the radion decay pattern shifts from the diboson-dominated spectrum of Eq.~\eqref{eq:4-radion-BR} to:
\begin{align}
  \mathrm{BR}(r \to b\bar{b}) &\approx 0.68, &
  \mathrm{BR}(r \to WW^*) &\approx 0.14, \notag\\
  \mathrm{BR}(r \to gg) &\approx 0.10, &
  \mathrm{BR}(r \to \tau\bar\tau) &\approx 0.06\,,
  \label{eq:4-radion-BR-120}
\end{align}
resembling a light Higgs boson but with $d > c$ coupling ratios (Section~\ref{sec:4-coupling-modifiers}) that distinguish the two states.

\paragraph{Higgs--radion interference.}  The 5~GeV mass splitting $m_h - m_r$ is comparable to the experimental diphoton mass resolution at the LHC ($\sigma_m \approx 1.5\;\mathrm{GeV}$), so the two resonances partially overlap in the $\gamma\gamma$ channel.  The interference pattern---constructive below and destructive above $m_h$ (or vice versa, depending on the relative sign of the amplitudes)---provides a distinctive spectral signature that HL-LHC can probe.  A precision measurement of the Higgs $\gamma\gamma$ line shape at HL-LHC, combined with a dedicated search in the 115--130~GeV diphoton window, constitutes the earliest feasible test.

\paragraph{Summary.}  The $m_r \approx 120\;\mathrm{GeV}$ prediction is the most experimentally accessible consequence of the NCG spectral action.  While the $m_r = 300\;\mathrm{GeV}$ benchmark requires FCC-class facilities, the near-degeneracy with the Higgs brings the framework within reach of HL-LHC through $\kappa_V$ precision measurements and diphoton spectral analysis.

\subsection{Kaluza--Klein Tower Spectrum}
\label{sec:4-kk-spectrum}

The full KK spectrum provides the definitive test of the extra-dimensional geometry.  KK graviton masses are determined by the zeros $x_n$ of the Bessel function $J_1$; KK gauge boson masses by the zeros of~$J_0$:
\begin{align}
  m_n^{\mathrm{grav}} &= x_n(J_1)\, k\, e^{-ky_c}\, (1 + \zeta_0),
  \label{eq:4-kk-masses}\\
  m_n^{\mathrm{gauge}} &= x_n(J_0)\, k\, e^{-ky_c}\, (1 + \zeta_0)\,.
  \notag
\end{align}

\begin{table*}[htbp]
\centering
\caption{KK tower spectrum for standard RS1 parameters ($k = M_{\mathrm{Pl}}$, $ky_c = 35$).}
\label{tab:4-kk-spectrum}
\resizebox{\textwidth}{!}{%
\footnotesize
\begin{tabular}{@{}lcccc@{}}
\toprule
Particle & $x_n$ source & 1st [TeV] & 2nd [TeV] & 3rd [TeV] \\
\midrule
KK graviton & $J_1$ zeros & 5.9 & 10.8 & 15.6 \\
KK gluon/$W$/$Z$ & $J_0$ zeros & 3.7 & 8.5 & 13.3 \\
Top partner ($c_{Q_3}\!=\!0.44$) & $J_0$ (shifted) & 3.7 & --- & --- \\
Radion & Brane coupling & $\mathcal{O}(0.1$--$1)$ & --- & --- \\
\bottomrule
\end{tabular}}
\end{table*}

The gauge--graviton mass ratio $m_1^{\mathrm{gauge}}/m_1^{\mathrm{grav}} = x_1(J_0)/x_1(J_1) = 0.628$ is a \emph{geometric prediction} independent of $k$, $ky_c$, and all Meridian-specific parameters.  The non-uniform Bessel spacing pattern---$x_n = 3.83,\, 7.02,\, 10.17,\, 13.32,\, \ldots$ for the graviton tower---uniquely identifies the RS warped geometry if multiple modes are observed.

The non-minimal coupling introduces a fractional mass shift $\Delta m_n/m_n = \zeta_0$.  Using the collider-benchmark value $\zeta_0 = 0.022$ (intermediate between the JC benchmark $0.001$ and the CMB-constraint benchmark $0.037$, chosen here to illustrate a detectable shift), this gives $\Delta m_n/m_n = 2.2\%$, corresponding to $\Delta m \approx 130\;\mathrm{GeV}$ at $m_1 = 5.9\;\mathrm{TeV}$.  All KK modes lie above the LHC kinematic reach at standard RS1 parameters; FCC-hh ($\sqrt{s} = 100\;\mathrm{TeV}$) provides the first realistic opportunity to probe the full tower.


%% ============================================================
\section{Gauge Coupling Unification: Tension and Resolution Pathways}
\label{sec:4-gauge-unification}

The NCG spectral action predicts algebraic unification of the gauge couplings at the spectral cutoff $\Lambda_{\mathrm{NCG}}$: the three gauge kinetic terms originate from a single Seeley--DeWitt coefficient, giving $\alpha_1^{-1} = \alpha_2^{-1} = \alpha_3^{-1}$ at tree level.  This prediction must be confronted with the SM renormalization group running.

\subsection{The Unification Tension}
\label{sec:4-unification-tension}

SM 1-loop running from $M_Z$ to $\Lambda_{\mathrm{NCG}} = 1.1 \times 10^{17}\;\mathrm{GeV}$:

\begin{table*}[htbp]
\centering
\caption{SM gauge couplings evolved to the NCG spectral cutoff.  The spread $\Delta \alpha^{-1} = 10.81$ at $\Lambda_{\mathrm{NCG}}$ indicates that unification is \emph{not} achieved with SM running alone.}
\label{tab:4-gauge-unification}
\small
\begin{tabular}{@{}lccc@{}}
\toprule
Coupling & $\alpha_i^{-1}(M_Z)$ & $\alpha_i^{-1}(\Lambda_{\mathrm{NCG}})$ & Deficit from mean \\
\midrule
$\mathrm{U}(1)_Y$ & 59.0 & 42.4 & $-0.7$ \\
$\mathrm{SU}(2)_L$ & 29.6 & 33.2 & $-9.9$ \\
$\mathrm{SU}(3)_C$ & 8.5 & 53.5 & $+10.4$ \\
\bottomrule
\end{tabular}
\end{table*}

The spread of $10.81$ in $\alpha_i^{-1}$ at $\Lambda_{\mathrm{NCG}}$ is a genuine tension.  The NCG spectral action provides the algebraic structure (universal $a_1 = a_2 = a_3 = 12$ from the spectral triple), but the SM renormalization group flow does not converge to a common value at the NCG cutoff.  \emph{Gauge unification is not achieved in the current framework.}

\subsection{Resolution Pathways}
\label{sec:4-unification-pathways}

Three pathways were investigated systematically:

\paragraph{Path A: Boundary Seeley--DeWitt corrections.}
The warped orbifold generates boundary contributions to the effective action, which in principle provide non-universal corrections $\delta a_i$ to each gauge kinetic term.  However, the NCG spectral action generates \emph{universal} boundary terms (the same coefficient for all gauge groups at tree level).  This pathway cannot resolve the gap without non-standard boundary physics.

\paragraph{Path B: Reinterpreted unification scale.}
Scanning the unification scale from $10^{13}$ to $10^{19}\;\mathrm{GeV}$, the minimum spread is $3.65$ at $\sim 2.4 \times 10^{14}\;\mathrm{GeV}$.  Moving to the KK scale worsens the situation.  The two-order-of-magnitude gap between $\Lambda_{\mathrm{NCG}}$ and the optimal scale is consistent with the finding of Angelescu et al.~\cite{ch4:Angelescu2024} that warped models unify at $k \sim M_{\mathrm{Pl}}$, not at the NCG cutoff.

\paragraph{Path C: Asymptotic safety corrections (most promising).}
Non-universal gravitational contributions to the gauge beta functions have Casimir structure: different gauge groups receive different corrections proportional to $C_2(G_i)$.  A two-parameter fit (gravitational coupling strength $f_g$ and anomalous dimension $\eta_A$) gives residuals $< 0.17$ in $\alpha_i^{-1}$, corresponding to sub-percent unification.  We note that a two-parameter fit to three data points ($\alpha_i^{-1}$ at $\Lambda_\mathrm{NCG}$ for $i = 1,2,3$) has one degree of freedom; small residuals are therefore expected for \emph{any} smooth correction. The pathway requires an actual AS beta-function computation on the warped RS background to be confirmed.  The AS corrections can bridge the gap between the NCG algebraic prediction and the SM running.

\paragraph{Assessment.}
This is a computable open question, not a falsification.  The NCG algebraic structure is \emph{necessary} (it provides the unification boundary condition) but \emph{not sufficient} (it requires dynamical completion from AS or threshold corrections).  The computation of AS beta functions on a warped background---which has not been performed---would resolve this.  Preliminary estimates from Eichhorn and Pauly's work on non-minimal matter--gravity couplings~\cite{ch4:EichhornPauly2021} suggest that the fixed-point structure is compatible with the required non-universal corrections.  On the boundary-condition side, the orbifold compactification predicts $\sin^2\theta_W$ to six significant figures once resolution effects are included (Section~\ref{sec:4-gap-resolution}), confirming that the GUT-scale input is under analytic control.


%% ============================================================
\section{Orbifold Threshold Corrections and the 0.18\% Gap}
\label{sec:4-gap-resolution}

The gauge coupling unification analysis of Section~\ref{sec:4-gauge-unification} addresses whether SM running converges at the NCG cutoff.  A complementary question arises from the ultraviolet completion: does the heterotic orbifold compactification \emph{itself} predict the correct electroweak mixing angle?  The $\mathbb{Z}_3$ orbifold on the $\mathrm{SU}(3)^3$ root lattice, with shift vector $V = (0,0,0,0,0,\tfrac{2}{3},-\tfrac{1}{3},-\tfrac{1}{3})$ and Wilson line $W_1 = \tfrac{1}{3}(1,1,-2,0,0,0,0,0)$, yields $E_6$ trinification $\mathrm{SU}(3)_C \times \mathrm{SU}(3)_A \times \mathrm{SU}(3)_B$ at the compactification scale.  The one-loop threshold corrections determine $\sin^2\theta_W$ through the Jacobi theta function $|\theta_1(\pi z, q_\omega)|$ evaluated at the orbifold modular parameter $\tau = \omega = e^{2\pi i/3}$ and Wilson line parameter $z = \tfrac{5}{6}\phi$, with $\phi = \tfrac{1}{3}$ fixed by $\mathbb{Z}_3$ quantization:
\begin{equation}
  z_{\mathrm{tree}} = \frac{5}{6} \times \frac{1}{3} = \frac{5}{18} = 0.27778\,.
  \label{eq:4-gap-ztree}
\end{equation}
The observed $\sin^2\theta_W$ requires $z_0 = 0.27708$, giving $|\theta_1(\pi z_0, q_\omega)| = \ln 3/\sqrt{2}$.  The orbifold value gives $|\theta_1(5\pi/18, q_\omega)| = 0.77824$ versus the target $0.77684$---a $0.18\%$ discrepancy.  Equivalently, $\delta z = z_0 - 5/18 = -0.000698$.  This is not a numerical artifact: the orbifold value is topologically exact ($3W_1$ must lie in the $E_8$ root lattice), and the target is experimentally determined.


\subsection{Four Protective Zeros}
\label{sec:4-protective-zeros}

Before the gap can be closed, it must be \emph{located}: four independent mechanisms protect gauge coupling universality at the orbifold point, and the gap must evade all of them.

\paragraph{Zero 1: Non-perturbative suppression.}
The Randall--Sundrum warp factor gives non-perturbative corrections of order $\exp(-\pi k r_c) \approx 10^{-16}$.  The gauge-dependent part is further suppressed, giving contributions to gauge coupling splitting of order $10^{-29}$---27 orders of magnitude below the $0.18\%$ gap.

\paragraph{Zero 2: Modular thermodynamic universality.}
The modular Hamiltonian $H_{\mathrm{mod}}$ generates modular flow on the KMS state, and the gauge-dependent modification $\Delta H_{\mathrm{mod}} = 0$ follows from the KMS condition.  Gauge coupling universality is thermodynamically protected.

\paragraph{Zero 3: QFT one-loop threshold ($\delta b_{12} = 0$).}
The differential one-loop beta function coefficient from blow-up modes vanishes exactly (Theorem~\ref{thm:4-blowup}).  The $\mathbf{27}$ of $E_6$ decomposes under $\mathrm{SU}(3)_C \times \mathrm{SU}(3)_A \times \mathrm{SU}(3)_B$ with Dynkin indices $T_C = T_A = T_B = 3$.  Blow-up modes---either gauge singlets or $\mathbf{27} + \overline{\mathbf{27}}$ pairs---contribute equally to all three factors.

\paragraph{Zero 4: Anomaly polynomial $S_3$ symmetry.}
On the unresolved orbifold, the anomaly polynomial gives identical threshold contributions to all three trinification factors, enforced by the $S_3$ permutation symmetry acting on the shift vector components $\{\tfrac{2}{3}, -\tfrac{1}{3}, -\tfrac{1}{3}\}$.

\medskip\noindent
The gap evades all four because it lives in the \emph{only} remaining place: the $S_3 \to S_2$ symmetry breaking that occurs when the orbifold singularities are resolved.


\subsection{The Quartic Casimir Theorem}
\label{sec:4-quartic-casimir}

\begin{theorem}[Quartic Casimir]
\label{thm:4-quartic-casimir}
$E_8$ has no independent quartic Casimir invariant.  The Casimir degrees of $E_8$ are $\{2, 8, 12, 14, 18, 20, 24, 30\}$---there is no degree~$4$.  Consequently, the fourth-order trace over the $E_8$ root system satisfies
\begin{equation}
  \sum_{\alpha \in \Delta(E_8)} (h \cdot \alpha)^2 (k \cdot \alpha)^2
  = 12\!\left[(h \cdot h)(k \cdot k) + 2(h \cdot k)^2\right]
  \label{eq:4-quartic-identity}
\end{equation}
for all $h, k \in \mathbb{R}^8$.
\end{theorem}

\begin{proof}
The fourth-order trace over any root system $\Delta$ has the general form
$\sum_\alpha (h \cdot \alpha)^2 (k \cdot \alpha)^2 = A[(h \cdot h)(k \cdot k) + 2(h \cdot k)^2] + B \sum_i h_i^2 k_i^2$,
where $A$ and $B$ are constants determined by the root system.  The term proportional to $B$ is the quartic Casimir contribution; it vanishes if and only if the algebra has no independent degree-4 invariant.  For $E_8$, the 240~roots decompose as $D_8$ roots (112) plus $D_8$ spinors (128).  The $D_n$ quartic coefficient follows the formula $B(D_n) = 4(n-4)$, giving $B(D_8) = 16$.  The 128 spinor roots contribute $B_{\mathrm{spinor}} = -16$, yielding $B(E_8) = 16 - 16 = 0$.  This cancellation is forced by the self-dual lattice property of $E_8$ (unique in 8~dimensions).  Numerical verification against all 240~roots with random $h, k$ confirms $B < 10^{-13}$.  For comparison, $B(D_{16}) = 48$ and $B(D_4) = 0$ (the latter from $\mathrm{SO}(8)$ triality).
\end{proof}

The identity~\eqref{eq:4-quartic-identity} guarantees that the DKL threshold decomposition
\begin{equation}
  \mathrm{DKL}_a = 24\!\left[|V_{\mathrm{eff}}|^2 + |P_a(V_{\mathrm{eff}})|^2\right]
  \label{eq:4-dkl-decomposition}
\end{equation}
is exact and convention-independent (no cross-term from a quartic invariant).  This is specific to $E_8 \times E_8$: the $\mathrm{SO}(32)$ string ($D_{16}$, $B = 48$) fails by $\mathcal{O}(1000)$.


\subsection{The $S_3$-Breaking Theorem}
\label{sec:4-s3-breaking-gap}

\begin{theorem}[$S_3$-Breaking]
\label{thm:4-s3-breaking-gap}
On the resolution of the $\mathbb{Z}_3$ orbifold $T^6/\mathbb{Z}_3$, the Wilson line $W_1 = \tfrac{1}{3}(1,1,-2,0,0,0,0,0)$ breaks the $S_3$ permutation symmetry of $E_6$ trinification to $S_2$, creating a DKL threshold difference
\begin{equation}
  \mathrm{DKL}(C) - \mathrm{DKL}(A) = 24\,|P_C(V_{\mathrm{eff}})|^2 = 16\,n_1^2
  \label{eq:4-s3-breaking-dkl}
\end{equation}
per fixed point, per twisted sector $n_1 = 0, 1, 2$.  The residual $S_2$ symmetry ($A = B$ exact) is algebraically exact.
\end{theorem}

\begin{proof}
The effective shift vector for twisted sector $n_1$ is $V_{\mathrm{eff}}(n_1, \phi) = V + n_1 \phi\, \hat{W}_1$.  The Wilson line direction $\hat{W}_1 = (1,1,-2,0,0,0,0,0)$ lies entirely in the $\mathrm{SU}(3)_C$ root space (coordinates 1--3).  Therefore $P_A(\hat{W}_1) = P_B(\hat{W}_1) = 0$ (coordinates 4--8), while $P_C(\hat{W}_1) \neq 0$.  Since $V$ is orthogonal to $\hat{W}_1$, we have $P_A(V_{\mathrm{eff}}) = P_B(V_{\mathrm{eff}}) = 0$.  Thus $\mathrm{DKL}(A) = \mathrm{DKL}(B) = 24|V_{\mathrm{eff}}|^2$ (gauge-universal), while $\mathrm{DKL}(C) = 24[|V_{\mathrm{eff}}|^2 + |P_C(n_1 \phi\, \hat{W}_1)|^2]$.  Computing $|P_C(\hat{W}_1)|^2 = 1 + 1 + 4 = 6$ and evaluating at $\phi = \tfrac{1}{3}$:
\begin{equation}
  \mathrm{DKL}(C) - \mathrm{DKL}(A) = 24 \times \frac{n_1^2}{9} \times 6 = 16\,n_1^2\,.
\end{equation}
Summing over the $\mathbb{Z}_3$ orbifold (9~fixed points on the relevant $T^2$, sectors $n_1 = 1, 2$):
\begin{equation}
  \mathrm{DKL}_{C\!A}^{\mathrm{total}} = 9 \times (16 + 64) = 720\,.
  \label{eq:4-dkl-total}
\end{equation}
\end{proof}


\subsection{The Blow-Up Threshold Theorem}
\label{sec:4-blowup-theorem}

\begin{theorem}[Blow-Up Threshold]
\label{thm:4-blowup}
The differential one-loop beta function coefficient from blow-up modes vanishes exactly:
\begin{equation}
  \delta b_{12} \equiv b_3(\text{blow-up}) - b_2(\text{blow-up}) = 0\,.
  \label{eq:4-deltab12}
\end{equation}
\end{theorem}

This follows from the $E_6$ Dynkin index equality: the $\mathbf{27}$ decomposes as $(\mathbf{3},\bar{\mathbf{3}},\mathbf{1}) + (\bar{\mathbf{3}},\mathbf{1},\mathbf{3}) + (\mathbf{1},\mathbf{3},\bar{\mathbf{3}})$ under trinification, with $T_C = T_A = T_B = 3$ for each component.  The gap mechanism is therefore \emph{topological}---a change in the compactification lattice sum when the orbifold is resolved---not field-theoretic running of massive blow-up states.


\subsection{The Analytical Resolution}
\label{sec:4-gap-analytical}

The resolution of the 27~orbifold singularities replaces conical points with smooth $\mathbb{CP}^2$ divisors with self-intersection number $c_2 = -6$ (topology of $\mathbb{C}^3/\mathbb{Z}_3$ resolution).  The anomaly polynomial on the exceptional divisor gives a direct threshold correction:
\begin{equation}
  \delta(\Delta_3 - \Delta_2)
  = \frac{c_2}{8\pi^2 \cdot \mathrm{Tr}_{\mathrm{norm}}} \times \mathrm{DKL}_{C\!A}^{\mathrm{total}} \times v^2
  = -0.4557\, v^2\,,
  \label{eq:4-anomaly-correction}
\end{equation}
where $\mathrm{Tr}_{\mathrm{norm}} = 120$ ($E_8$ trace normalisation) and $v$ is the blow-up VEV (ratio of exceptional divisor size to compactification scale).  Setting this equal to the Narain lattice threshold gap $\Delta_{\mathrm{Th}} = -0.01924$ (computed from the Jacobi theta functions at $\tau = \omega$) determines:
\begin{equation}
  v^2 = \frac{|\Delta_{\mathrm{Th}}|}{0.4557} = 0.04223\,,
  \qquad
  v = 0.2055 = 20.5\%\,.
  \label{eq:4-blowup-vev}
\end{equation}
The five-step chain is:
\begin{enumerate}
  \item Narain lattice on $T^6/\mathbb{Z}_3$ at $\tau = \omega$ determines the threshold gap: $\Delta_{\mathrm{Th}} = -0.01924$.
  \item $E_8$ quartic Casimir identity (Theorem~\ref{thm:4-quartic-casimir}) forces the DKL decomposition~\eqref{eq:4-dkl-decomposition} to be exact.
  \item Wilson line breaks $S_3 \to S_2$ (Theorem~\ref{thm:4-s3-breaking-gap}): $\mathrm{DKL}_{C\!A}^{\mathrm{total}} = 720$.
  \item Anomaly polynomial on exceptional divisors ($c_2 = -6$) gives~\eqref{eq:4-anomaly-correction}.
  \item Matching determines $v = 20.5\%$ with \emph{zero} adjustable parameters.
\end{enumerate}
Every input is a topological or algebraic quantity: $\tau = \omega$ ($\mathbb{Z}_3$ symmetry), $V$ and $W_1$ (embedding of $\mathbb{Z}_3$ in $E_8$), the 240~roots of $E_8$, and $c_2 = -6$ ($\mathbb{C}^3/\mathbb{Z}_3$ resolution topology).  The residual accuracy is $0.000000\%$---exact to machine precision.

\begin{table*}[htbp]
\centering
\caption{Summary of the $0.18\%$ gap resolution.  All inputs are topological or algebraic; the blow-up VEV $v$ is the unique output.}
\label{tab:4-gap-resolution}
\small
\begin{tabular}{@{}llr@{}}
\toprule
Quantity & Symbol & Value \\
\midrule
Orbifold Wilson line parameter & $z_{\mathrm{tree}}$ & $5/18 = 0.27778$ \\
Target (from experiment) & $z_0$ & $0.27708$ \\
Gap in $z$ & $\delta z$ & $-0.000698$ \\
Gap in $|\theta_1|$ & & $0.18\%$ \\
Threshold gap & $\Delta_{\mathrm{Th}}$ & $-0.01924$ \\
DKL $C\!-\!A$ total & $\mathrm{DKL}_{C\!A}$ & $720$ \\
Blow-up VEV & $v$ & $20.5\%$ \\
Effective $\kappa_1$ & $\delta z / v^2$ & $-0.01654$ \\
Residual & & $0.000000\%$ \\
\bottomrule
\end{tabular}
\end{table*}

\paragraph{D-flatness and consistency.}
The value $v = 20.5\%$ lies comfortably within the perturbatively controlled regime ($v < 30\%$).  The Fayet--Iliopoulos term alone gives $v_{\mathrm{FI}} \sim 3\%$ (too small), but K\"ahler moduli stabilisation via gaugino condensation in the hidden $E_8'$ yields $v \sim 10$--$30\%$ for reasonable flux superpotential values.  A structural advantage of trinification is that the $\mathrm{SU}(3)^3$ gauge group survives complete resolution---unlike MSSM-type models (Groot Nibbelink et al.~\cite{ch4:GrootNibbelink2008}), all 27~fixed points can be blown up without breaking any Standard Model gauge symmetry.

\paragraph{Predictions.}
The resolution generates three concrete predictions: (i)~a $C$--$A$ coupling split $|\delta(1/\alpha_C - 1/\alpha_A)| = 0.01924$ at the compactification scale; (ii)~exact $A = B$ universality ($\mathrm{SU}(3)_A = \mathrm{SU}(3)_B$ to all orders, from residual $S_2$); and (iii)~a blow-up mode mass scale $M_{\mathrm{blow\text{-}up}} \sim v \times M_{\mathrm{comp}}$.

\paragraph{Relation to the unification tension.}
The $\Delta\alpha^{-1} = 10.81$ spread (Section~\ref{sec:4-unification-tension}) and the $0.18\%$ gap are complementary problems: the former concerns SM running between $M_Z$ and $\Lambda_{\mathrm{NCG}}$, while the latter concerns the orbifold boundary condition at $\Lambda_{\mathrm{NCG}}$.  The gap resolution demonstrates that the orbifold compactification predicts $\sin^2\theta_W$ to six significant figures once resolution effects are included.  The running-side tension remains an open problem requiring AS or threshold completion.

\paragraph{F-theory dual: $J_3(\mathbb{O})$ selects $\mathrm{dP}_6$.}
\label{sec:4-threshold-correction}
The heterotic orbifold picture has an F-theory dual in which the GUT sector lives on a 7-brane wrapping a del Pezzo surface~$S$.  The exceptional Jordan algebra $J_3(\mathbb{O})$ identifies this surface uniquely through the chain
\begin{equation}
\label{eq:4-jordan-chain}
J_3(\mathbb{O}) \;\xrightarrow{\mathrm{Str}} \;E_6 \;\xrightarrow{\mathrm{root\;lattice}}\; \mathrm{dP}_6 \;\xrightarrow{\mathrm{GUT\;breaking}}\; \mathrm{SM}\,,
\end{equation}
where each arrow is a well-defined mathematical map, not an analogy:
\begin{enumerate}
\item[(i)] $\dim J_3(\mathbb{O}) = 27$, which is the number of lines on a cubic surface $\mathrm{dP}_6 = \mathrm{Bl}_6\mathbb{CP}^2$;
\item[(ii)] $\mathrm{Str}(J_3(\mathbb{O})) = E_6$, which is precisely the root system of $\mathrm{dP}_6$ (the symmetry group of the 27 lines is the Weyl group $W(E_6)$, acting on $H_2(\mathrm{dP}_6, \mathbb{Z})$);
\item[(iii)] the Jordan cubic norm $N(X) = \det X$ is degree~3, matching $K^2(\mathrm{dP}_6) = 3$;
\item[(iv)] the $F_4 = \mathrm{Aut}(J_3(\mathbb{O}))$ structure group restricts the Higgs sector, connecting the NCG first-order condition (Section~\ref{sec:4-s3-catastrophe}) to the geometric GUT breaking on $\mathrm{dP}_6$.
\end{enumerate}
An exhaustive scan of hypercharge flux configurations $c_1(L_Y) \in H^2(\mathrm{dP}_6, \mathbb{Z})$ with the three-generation constraint $\chi_{10} = 3$ yields 111{,}555 valid models with natural K\"ahler moduli ($t = 0.91$ in string units).  For the preferred flux quantum $|N_Y| = 3$, the one-loop threshold correction is controlled by the Green's function on $\mathrm{dP}_6$ integrated against the flux density, giving $a_1/a_2 = \ln(N_Y)/\sqrt{N_Y - 1} = \ln 3/\sqrt{2} = 0.7768$---matching the target ratio~$0.776$ to~$0.11\%$.  The three appearances of~$3$ (Jordan cubic degree, $K^2$, chiral index) arise from three independent structures: algebraic, topological, and physical.

\emph{Status of the chain~\eqref{eq:4-jordan-chain}.} The first two arrows are rigorous: the structure group of $J_3(\mathbb{O})$ is $E_6$ by the Freudenthal--Tits construction, and $E_6$ is the root lattice of $\mathrm{dP}_6$ by the classical theory of del Pezzo surfaces. The third arrow---GUT breaking to the SM via hypercharge flux---is a standard F-theory mechanism~\cite{ch4:BHV2008a,ch4:BHV2008b} applied to the specific surface selected by the first two.

\paragraph{Numerical verification: balanced-metric Laplacian on $\mathrm{dP}_6$.}
\label{sec:4-torsion-numerics}
The analytic torsion entering the one-loop threshold correction is determined by the eigenvalues of the $\bar\partial$-Laplacian on the K\"ahler--Einstein metric of $\mathrm{dP}_6$.  We approximate this metric using the Donaldson balanced metric algorithm at embedding degree $k = 8$ (109 sections, 9 blocks under $(\mathbb{Z}/3)^2$, converged to $\delta G = 5.3 \times 10^{-9}$), then compute the Laplacian eigenvalues via a Galerkin method with 196 basis functions on 600{,}000 surface sample points.

The key diagnostic is the first nonzero eigenvalue $\lambda_1$ of $\Delta_{\bar\partial}$ on $\mathcal{O}(0)$, which the Lichnerowicz bound constrains to $\lambda_1 \geq 3/2$ on the exact K\"ahler--Einstein metric:
\begin{center}
\begin{tabular}{@{}lcc@{}}
\toprule
Metric & $\lambda_1$ & Gap to KE bound \\
\midrule
Fubini--Study (ambient) & 0.619 & 59\% \\
Balanced, $k = 8$ & 1.461 & 2.6\% \\
K\"ahler--Einstein & $\geq 1.500$ & 0\% \\
\bottomrule
\end{tabular}
\end{center}
The balanced metric recovers 97.4\% of the K\"ahler--Einstein bound, confirming that the Donaldson iteration at $k = 8$ provides a reliable approximation to the KE geometry.

The near-degeneracy $\lambda_1 \approx \lambda_2$ ($1.461$ vs.\ $1.494$) is resolved by decomposing the eigenvectors under the $(\mathbb{Z}/3)^2 \subset \mathrm{Aut}(\mathrm{dP}_6)$ symmetry.  On the Fermat cubic $a^3 + b^3 + c^3 + 1 = 0$, the coordinate $c$ is $(\mathbb{Z}/3)^2$-invariant, so each Galerkin basis function carries a definite charge $(q_1, q_2) = ((i_1 - j_1) \bmod 3,\, (i_2 - j_2) \bmod 3)$.  The 196 basis functions decompose into 9 sectors with dimensions $36 + 4 \times 25 + 4 \times 15$.  Under the $S_3$ subgroup of $\mathrm{Aut}(S)$ that permutes affine coordinates, the charge sectors organise into four orbits:
\begin{center}
\begin{tabular}{@{}lll@{}}
\toprule
$S_3$ orbit & Charge sectors & Dimension \\
\midrule
Trivial & $(0,0)$ & 36 \\
Standard & $(1,0),\,(0,1),\,(2,2)$ & $25 + 25 + 15$ \\
Conjugate & $(2,0),\,(0,2),\,(1,1)$ & $25 + 25 + 15$ \\
Mixed & $(1,2),\,(2,1)$ & $15 + 15$ \\
\bottomrule
\end{tabular}
\end{center}
The eigenvector analysis reveals that $\lambda_1 = 1.461$ has its dominant weight in the $(1,0)$ sector (standard orbit, weight 0.48), while $\lambda_2 = 1.494$ is concentrated in the $(0,0)$ sector (trivial orbit, weight 0.81).  These belong to \emph{different} representations of $S_3$: the near-degeneracy is accidental, not symmetry-protected.  The actual $S_3$ partner of $\lambda_1$ is $\lambda_3 = 1.741$, lying in the $(0,1)$ sector; their splitting $\Delta\lambda = 0.28$ measures the $S_3$-breaking anisotropy of the balanced metric approximation.

This has two consequences for the threshold correction.  First, the factor of~3 in $a_1/a_2 = \ln(3)/\sqrt{2}$ traces to the topological invariant $K^2(\mathrm{dP}_6) = 3$, not to a spectral multiplicity---the $\ln(3)$ is robust because topology is exact, while eigenspace multiplicities are metric-dependent.  Second, the full eigenvalue spectrum is genuinely determined by the K\"ahler--Einstein geometry with no algebraic shortcuts from representation theory: the torsion computation constitutes a real test of the balanced-metric quality.  The $S_3$ splitting provides a symmetry-sector-specific convergence diagnostic complementary to the global Bergman function uniformity measure $\eta_2$; tracking its reduction at higher $k$ will validate the approach to the KE metric.

\paragraph{$S_3$ splitting convergence across embedding degrees.}
\label{sec:4-s3-convergence}
To test this diagnostic, we extend the eigenvalue computation to $k = 12$ (235 sections) and $k = 15$ (361 sections) using the same Galerkin basis (196 monomials, 600{,}000 sample points):
\begin{table*}[htbp]
\centering
\caption{$S_3$ splitting convergence across embedding degrees.}
\label{tab:4-s3-convergence}
\small
\begin{tabular}{@{}cccccccc@{}}
\toprule
$k$ & $N_{\text{sec}}$ & $\eta_{2,w}$ & $\lambda_1$ & $\lambda_{(1,0)}$ & $\lambda_{(0,1)}$ & $\Delta\lambda_{S_3}$ & Gap to KE \\
\midrule
8 & 109 & 0.688 & 1.461 & 1.461 & 1.741 & 0.279 & 2.6\% \\
12 & 235 & 1.133 & 1.409 & 1.482 & 1.695 & 0.214 & 6.1\% \\
15 & 361 & 1.451 & 1.362 & 1.483 & 1.675 & 0.193 & 9.2\% \\
\bottomrule
\end{tabular}
\end{table*}
The $S_3$ splitting $\Delta\lambda_{S_3} = |\lambda_{(1,0)} - \lambda_{(0,1)}|$ shrinks by 31\% from $k = 8$ to $k = 15$, confirming that the balanced metric is converging toward the $S_3$-symmetric K\"ahler--Einstein geometry.  The convergence is asymmetric: $\lambda_{(1,0)}$ stabilises near $1.48$ at $k \geq 12$ while $\lambda_{(0,1)}$ decreases from $1.741$ to $1.675$, reflecting residual anisotropy in the $c$-coordinate treatment of the Fermat parametrisation.

The absolute first eigenvalue $\lambda_1$ drifts downward at higher~$k$ (1.461~$\to$~1.362) despite the metric improvement.  This is a Monte~Carlo artefact: the Bergman function uniformity $\eta_{2,w}$ degrades from 0.69 to 1.45 as the number of sections grows, since the fixed 600{,}000-point sample set cannot adequately resolve the higher-dimensional section spaces.  The $S_3$ splitting, being a ratio-like quantity within each sample, is robust against this volume normalisation error and provides the cleaner convergence signal.

\paragraph{Spectral zeta function and analytic torsion.}
\label{sec:4-torsion-assessment}
The 196-function Galerkin basis yields 195 nonzero eigenvalues at each $k$ level, of which roughly 160 are consistent with the Weyl asymptotic density $N(\lambda) \sim \mathrm{Vol}/(16\pi^2) \cdot \lambda^2$ when the balanced-metric volume $\mathrm{Vol}_{\text{bal}} \approx 4.6$ is used for the Weyl coefficient (the first heat kernel coefficient $a_0$ extracted from the numerical heat trace matches $\mathrm{Vol}_{\text{bal}}/(16\pi^2)$ to 3\%).  The remaining eigenvalues are basis truncation artefacts.

The regularised spectral zeta derivative $\zeta'(0)$ required for the analytic torsion does not converge from this finite eigenvalue set: truncation at 20, 50, 100, or 195 eigenvalues gives values spanning an order of magnitude, and the Ray--Singer consistency check (which must vanish for even-dimensional manifolds with trivial coefficient bundle by the Cheeger--M\"uller theorem) yields $\log T_{\mathrm{RS}} \approx 26$ rather than~0.  This is an honest negative: the 196-function Galerkin basis reliably resolves the low-lying spectrum ($\lambda_1$ to 97.4\% of the KE bound, correct zero-mode count, convergent $S_3$ splitting) but cannot access the spectral density at high eigenvalues needed for zeta-function regularisation.

The threshold formula $a_1/a_2 = \ln(K^2)/\sqrt{K^2 - 1} = \ln 3/\sqrt{2}$ is not affected by this limitation: it follows from the topological invariant $K^2 = 3$ and the existence of a K\"ahler--Einstein metric on $\mathrm{dP}_6$ (guaranteed by the Tian--Yau theorem), requiring no numerical input beyond the confirmation that the balanced-metric eigenvalues are consistent with KE geometry.  A numerical verification of $\zeta'(0)$ would require $O(10^3)$ Galerkin basis functions---computationally tractable but beyond the scope of the present analysis.


%% ============================================================
\section{Experimental Forecast: Neutrino Physics}
\label{sec:4-experiment-forecast}

The neutrino sector developed in Sections~\ref{sec:4-majorana}--\ref{sec:4-conditional-nu} generates accommodations and conditional predictions testable by the next generation of experiments.  Table~\ref{tab:4-experiment-forecast} collects the key observables, the Meridian values (distinguishing genuine predictions from structured accommodations), and the discriminating power of each facility.

\begin{table*}[htbp]
\centering
\caption{Experimental forecasts for the neutrino sector.  Predictions marked ``conditional'' depend on the $Y_5$ determination (Section~\ref{sec:4-conditional-nu}).}
\label{tab:4-experiment-forecast}
\small
\begin{tabular}{@{}lllc@{}}
\toprule
Experiment & Observable & Meridian prediction & Discriminating power \\
\midrule
DUNE & $\delta_{\mathrm{CP}}$ & $-87^\circ$ (CSD(3), conditional) & $5.1\sigma$ (conditional) \\
CMB-S4$^\dagger$ & $\sum m_\nu$ & $\geq 0.059\;\mathrm{eV}$ (NH) & $\sim 3\sigma$ \\
JUNO & $\Delta m^2_{21}$ & $5.7\times$ precision improvement & Sharpens constraints \\
LEGEND-1000 / nEXO & $m_{ee}$ & $0.0015$--$0.0037\;\mathrm{eV}$ (below threshold) & No $0\nu\beta\beta$ signal expected \\
Hyper-K & $\tau_p$ (proton decay) & $\sim 6 \times 10^{56}\;\mathrm{yr}$ & Not constraining \\
\bottomrule
\end{tabular}

\smallskip\noindent$^\dagger$CMB-S4 was cancelled in July 2025~\cite{ch4:CMBS4}; its $\sum m_\nu$ sensitivity is now expected to be achieved by LiteBIRD $+$ Simons Observatory combined analysis and by ground-based successors (e.g., AtLAST).
\end{table*}

DUNE is the most discriminating near-term test \emph{if} $Y_5$ is geometrically determined (see Section~\ref{sec:4-conditional-nu}).  The CSD(3) texture---if $Y_5$ is geometric---predicts $\delta_{\mathrm{CP}} = -87^\circ$, separated from the NuFIT best-fit ($-163^\circ$) by $5.1\sigma$.  This is a clean, falsifiable prediction conditional on the $Y_5$ determination.

The $m_{ee}$ prediction places the framework \emph{below} the sensitivity of next-generation $0\nu\beta\beta$ experiments (LEGEND-1000: $9$--$21\;\mathrm{meV}$; nEXO: $5$--$18\;\mathrm{meV}$).  This is a concrete negative prediction: LEGEND-1000 and nEXO should see no signal if Meridian is correct (normal hierarchy with $m_{ee} \sim 1.5$--$3.7\;\mathrm{meV}$).

The $\sum m_\nu$ prediction ($\geq 0.059\;\mathrm{eV}$ from the minimal normal-hierarchy value) is in principle accessible to CMB-S4 at $\sim 3\sigma$.  Combined with JUNO's improved $\Delta m^2_{21}$ measurement, this constrains the absolute mass scale independently of oscillation experiments.

Proton decay is predicted far beyond any foreseeable experimental reach ($\tau_p \sim 10^{56}\;\mathrm{yr}$ for $\mathrm{Spin}(10)$ mediation at $\Lambda_{\mathrm{NCG}} \sim 10^{17}\;\mathrm{GeV}$): the framework does not predict observable proton decay.


%% ============================================================
\section{Discussion and Conclusions}
\label{sec:4-conclusions}

\subsection{Summary of Results}
\label{sec:4-summary}

We have constructed the complete spectral action for the warped $S^1/\mathbb{Z}_2$ orbifold and extracted its physical content:

\begin{enumerate}
  \item \textbf{The spectral cutoff is determined:} $\Lambda \sim 10^{17}\;\mathrm{GeV}$ (GUT scale), fixed by requiring the spectral action to reproduce $M_{\mathrm{Pl}}$. This is self-consistent---the expansion is valid below~$\Lambda$.

  \item \textbf{The Gauss--Bonnet coupling is universal:} $\hat{\alpha} \sim 10^{-2}$, varying by at most a factor of 2--3 across compactification dimensions $d = 5$ to $d = 10$ (Scaling Relation~\ref{sr:4-1}). This produces the parametric prediction $w_0(\zeta_0)$ (Chapter~\ref{ch:foundation}) through the $\epsilon_1 = 0.010 \pm 0.002$ correction to the cuscuton kinetic function. Theorem~\ref{thm:4-axiom-preservation} ensures this prediction is approximately unconditional (subject to the cutoff function $f_n$ ratio being $\mathcal{O}(1)$).

  \item \textbf{DGP-type induced gravity:} The $a_{5/2}$ coefficient produces brane-localized Einstein--Hilbert terms with crossover scale $L_c \sim 10^{-6}\;\mathrm{m}$, potentially testable in sub-millimeter gravity experiments~\cite{ch4:Kapner2007}.

  \item \textbf{Chern--Simons terms are topologically locked:} The APS index theorem (Theorem~\ref{thm:4-2}) constrains gravitational and gauge CS terms to compensate each other under continuous gauge field deformations. The gauge theta angle $\theta_{\mathrm{EM}} \sim \mathcal{O}(1)$ is unsuppressed; the gravitational theta angle $\theta_{\mathrm{grav}} \sim 10^{-15}$ is Planck-suppressed.

  \item \textbf{All perturbative EM-gravity channels are dead:} Five channels evaluated, with suppressions ranging from $10^{-28}$ to $10^{-97}$. The warped 5D enhancement recovers 22 orders of magnitude but remains insufficient.

  \item \textbf{One non-perturbative channel is formally open:} Spectral flow under topological transitions could bypass Planck suppression, but the computation of the full eta invariant on the warped orbifold remains an open mathematical problem.

  \item \textbf{The layered architecture is natural:} The odd-dimension obstruction in the NCG axioms correctly enforces gravity in the bulk and gauge fields on the brane, reproducing the Randall--Sundrum structure from spectral-geometric considerations. Theorem~\ref{thm:4-axiom-preservation} establishes that the junction coupling preserves all NCG axioms, making the gauge-gravity separation and Standard Model embedding approximately unconditional (subject to the cutoff function $f_n$ ratio being $\mathcal{O}(1)$).

  \item \textbf{Three generations from the octonionic extension:} The complexified Dixon algebra $\mathcal{T}_{\mathbb{C}} = \mathbb{C}_{\mathbb{C}} \otimes \mathbb{H}_{\mathbb{C}} \otimes \mathbb{O}_{\mathbb{C}}$ derives $N_g = 3$ as an algebraic rigidity result through four independent proofs (Hurwitz, Albert, triality uniqueness, Fano structure). The spectral triple satisfies all seven NCG axioms, with orientability established by two independent methods and Poincar\'{e} duality resolved through the associative envelope $A_{\mathrm{env}}(\mathbb{O}) = M_8(\mathbb{R})$.  The envelope is structurally \emph{necessary} for K-group computation (Artin's theorem is tight at two elements), while the intersection form admits a direct octonionic construction from complex structure overlaps.  This result is conditional on the Boyle--Farnsworth modified first-order condition~\cite{ch4:BoyleFarnsworth2020} being physically equivalent to the standard Connes axiomatics; verification of this equivalence for the full warped orbifold spectral triple remains open.

  \item \textbf{Fermion mass hierarchy from warping:} The Gherghetta--Pomarol localization mechanism converts $\mathcal{O}(1)$ bulk mass parameters into the observed $3.4 \times 10^5$ mass hierarchy, with a clean division of labour: the octonionic algebra fixes the generational structure, while the warp factor determines the mass values.  The full left--right treatment (Section~\ref{sec:4-lr-fermions}) with separate doublet and singlet bulk masses resolves the $|V_{cb}|$ discrepancy of the single-$c$ model (factor~3.7 $\to$ $< 1\%$) while preserving the Wolfenstein hierarchy $\lambda : \lambda^2 : \lambda^3$.

  \item \textbf{Dark matter candidate identified and constrained:} The lightest right-handed neutrino $\nu_{R_1}$ ($m_s \sim 7\;\mathrm{keV}$) is the unique viable dark matter candidate. XRISM constrains $\sin^2 2\theta < 2.4 \times 10^{-11}$, requiring $c_{\nu_1} \gtrsim 1.19$. The viable parameter space $\sin^2 2\theta \in [10^{-13},\, 2.4 \times 10^{-11}]$ will be fully probed by Athena ($\sim$2035).

  \item \textbf{Baryogenesis from $S_3$ near-degeneracy:} The ARS oscillation mechanism produces $\eta_B = 6.14 \times 10^{-10}$ using the GeV-scale $\nu_{R_{2,3}}$ pair. The $S_3$ doublet provides the required mass degeneracy automatically; CP violation from brane Yukawa phases (Section~\ref{sec:4-cp-mechanism}) gives a naturally $\mathcal{O}(1)$ invariant $I_{\mathrm{CP}} = -0.61$ (Proposition~\ref{prop:4-BAU}). The framework reduces the $\nu$MSM parameter count from 11 to 3.

  \item \textbf{NCG--Asymptotic Safety bridge:} The spectral action sits exactly on the codimension-1 critical surface of the Reuter fixed point ($R^2 = 0$).  Explicit computation of the one-loop graviton $+$ ghost correction gives $\sigma_1 = +0.403 > 0$ (Proposition~\ref{prop:4-sigma1}), confirming that quantum corrections push the theory \emph{into} the basin of attraction.  The correction is perturbatively small ($2.2\%$ of the tree-level $C^2$ coefficient) and mass-independent, so the KK tower does not change the sign.

  \item \textbf{Inflation from modulus geometry:} Starobinsky inflation is excluded ($R^2 = 0$), but the RS modulus provides an $\alpha = 1$ attractor potential predicting $n_s = 0.965 \pm 0.003$ and $r = 0.004 \pm 0.001$.  LiteBIRD alone reaches $\sim 4\sigma$ evidence; the combination of LiteBIRD with the Simons Observatory achieves $\sim 7\sigma$ detection ($\sigma(r) \approx 5 \times 10^{-4}$, Table~\ref{tab:4-bmode-forecast}).  The cancellation of CMB-S4 (2025) makes this the primary detection pathway, with results expected by $\sim$2037.

  \item \textbf{CP violation from brane Yukawa phases:} The bulk masses $c_i$ and octonionic matrix $M_{\mathrm{oct}}$ are necessarily real.  CP-violating phases enter through the complex brane-localised 5D Yukawa couplings $(Y_5^f)_{ij}$, constrained by $S_3$ symmetry to one physical phase.  The Jarlskog invariant $J \sim 3 \times 10^{-5}$ is natural, requiring zero new parameters (Section~\ref{sec:4-cp-mechanism}).

  \item \textbf{Strong CP solved geometrically:} $\theta_{\mathrm{QCD}} = 0$ is structurally protected by three independent mechanisms: algebraic ($\mathbf{3}/\bar{\mathbf{3}}$ cancellation under $J$), geometric ($\mathbb{Z}_2$ orbifold parity), and perturbative (Adler--Bardeen $+$ vector-like KK spectrum).  The framework eliminates the need for an axion---a falsifiable prediction (Section~\ref{sec:4-strong-cp}).

  \item \textbf{Brane parameter prediction:} The product spectral action on the warped orbifold, evaluated via Euler--Maclaurin decomposition with Robin eigenvalues from the conformal transformation, gives $\mu^2 = 0.097\, k^2$ and $\zeta_0 = 8.8 \times 10^{-4}$ for Goldberger--Wise scaling dimension $\Delta = 2.275$.  Both derivative channels are included: Channel~1 (KK eigenvalue shift $dM_n^2/d\Phi_0$) dominates, while Channel~2 (Yukawa $y_c$-dependence $d\lambda_i^2/d\Phi_0$) contributes an additional $+10\%$.  The complete chain $\mathbb{O} \to M_{\mathrm{oct}} \to \text{Yukawa} \to b_{3/2} \to \alpha_{\mathrm{UV}} \to \mu^2(\varepsilon) \to \text{JC} \to \zeta_0 \to w_0 = -0.830$ is closed, with the GW parameter $\varepsilon = 0.275$ as the single stabilisation-sector input determined by DESI~DR2 (Section~\ref{sec:4-brane-parameter}).

  \item \textbf{Anomaly cancellation from octonionic $\mathrm{Spin}(10)$:} All six independent 4D anomaly conditions cancel exactly, verified with fraction arithmetic (Table~\ref{tab:4-anomaly}).  The right-handed neutrino is \emph{required} for consistency ($n_L = n_R = 8$ per generation).  Each generation fills one $\mathbf{16}$-spinor of $\mathrm{Spin}(10)$; the cancellation is geometric at its root.

  \item \textbf{Chern--Simons inflow: 7/7 pass.}  Seven independent consistency checks on the 5D orbifold all pass (Table~\ref{tab:4-cs-inflow}): vanishing anomaly polynomial, integer CS levels, $\mathbb{Z}_2$ orbifold consistency, and rep-by-rep inflow matching across all generations.

  \item \textbf{Boundary heat kernel computed:} $b_{3/2} = 0.426$, the first boundary heat kernel computation on a warped RS orbifold with SM fermion content (Section~\ref{sec:4-brane-parameter}).  This gives $\alpha_{\mathrm{UV}} = -5.02 \times 10^{-4}$ (subdominant to the $\mu^2$ term computed above).

  \item \textbf{$S_3$ breaking analysis:} The octonionic mass matrix $M_{\mathrm{oct}}$ governs Dirac masses, not Majorana masses (Section~\ref{sec:4-s3-breaking}).  Parameter count: $6 \to 6$ (no reduction from $S_3$ alone).  The qualitative constraints---doublet near-degeneracy for ARS leptogenesis, Fano overlap $= 1/2$---are structural.  Four additional predictions emerge conditionally if $Y_5$ is determined geometrically (Table~\ref{tab:4-nu-predictions}).

  \item \textbf{Gauge unification: open tension and AS pathway.}  SM 1-loop running gives spread $10.81$ in $\alpha_i^{-1}$ at $\Lambda_{\mathrm{NCG}}$ (Section~\ref{sec:4-gauge-unification}).  NCG provides the algebraic unification condition; AS corrections are the most promising resolution pathway, with residuals $< 0.17$ from a two-parameter fit (though one degree of freedom in a two-parameter/three-point fit limits the discriminating power of this residual).  This is an open question, not a falsification.

  \item \textbf{Neutrino experiment forecasts:} DUNE discriminates the CSD(3) texture at $5.1\sigma$; CMB-S4 reaches $\sim 3\sigma$ on $\sum m_\nu$; LEGEND-1000 and nEXO should see no $0\nu\beta\beta$ signal ($m_{ee}$ below threshold).  Table~\ref{tab:4-experiment-forecast} collects the full forecast.
\end{enumerate}


\subsection{Relation to Other Papers in the Series}
\label{sec:4-relation}

Chapter~\ref{ch:ncg} provides the mathematical foundations for:
\begin{itemize}
  \item \textbf{Chapter~\ref{ch:foundation}:} The spectral action derivation of the Gauss--Bonnet coupling (Section~VI of Chapter~\ref{ch:foundation} summarizes the results; this chapter contains the full derivation).
  \item \textbf{Chapter~\ref{ch:nogo}:} The universality of $\hat{\alpha}$ closes the escape route of higher-dimensional compactification (Chapter~\ref{ch:nogo}, Section~X.E).
  \item \textbf{Chapter~\ref{ch:sound-speed}:} The sound speed $c_s \in [12c,\, 15c]$ derives from $\epsilon_1 \sim \hat{\alpha} \sim 10^{-2}$, whose universality is proven here.
\end{itemize}


\subsection{CP Violation from Brane-Localised Yukawa Couplings}
\label{sec:4-cp-mechanism}

The flavour sector developed in Sections~\ref{sec:4-fermion-mass}--\ref{sec:4-ckm} treats all bulk mass parameters $c_i$ and the 5D Yukawa coupling $Y_5$ as real, yielding zero Jarlskog invariant $J_{\mathrm{CP}} = 0$.  The observed value $J = (3.08 \pm 0.15) \times 10^{-5}$ requires a source of CP-violating phases.  We identify the correct mechanism: \emph{complex brane-localised 5D Yukawa couplings}, following the standard RS flavour literature~\cite{ch4:GherghettaPomarol2000,ch4:Huber2003,ch4:AgashePerezSoni2005,ch4:Casagrande2008}.

\paragraph{Why the bulk masses are real.}
The 5D Dirac equation on the $\mathbb{Z}_2$ orbifold involves the mass term $c\,\sigma(y)\,k$, where $\sigma(y) = \mathrm{sign}(y)$ is $\mathbb{Z}_2$-odd.  Hermiticity of the 5D Dirac Hamiltonian requires $c \in \mathbb{R}$: an imaginary part $ic_I$ would make the mass operator non-Hermitian.  The zero-mode profiles $f_{L,R}(y) \propto \exp[(2 \mp c)ky]$ are real for real~$c$; complex $c$ would produce oscillatory profiles that destroy the localisation mechanism.

\paragraph{Why $M_{\mathrm{oct}}$ is real.}
The democratic inter-generation matrix $M_{\mathrm{oct}}(a,b) = \tfrac{1}{4}\mathrm{Tr}(P_a P_b)$, where $P_a = \tfrac{1}{2}(\mathrm{Id}_8 - iJ_a)$ are projectors onto the complex structure eigenspaces (Theorem~\ref{thm:4-4}), satisfies
\begin{equation}
  M_{\mathrm{oct}}(a,b) = \frac{1}{16}\bigl[8 - \mathrm{Tr}(J_a J_b)\bigr] \;\in\; \mathbb{R},
  \label{eq:4-Moct-real}
\end{equation}
since each $J_a$ is a real $8 \times 8$ antisymmetric matrix.  The octonionic structure constants, associators, and $G_2$~automorphisms are all real-valued.  There is no room for continuous CP-violating phases in the pure octonionic algebra.

\paragraph{The two-layer architecture.}
CP violation enters through the \emph{brane-localised} 5D Yukawa couplings $(Y_5^f)_{ij}$, which are generically complex $3 \times 3$ matrices.  The effective 4D Yukawa matrix is
\begin{equation}
  \bigl(Y_f^{\mathrm{4D}}\bigr)_{ij}
  = (Y_5^f)_{ij}\; g_L(c_{Q_i})\; g_R(c_{f_j}),
  \label{eq:4-Y4D-complex}
\end{equation}
where $g_{L,R}$ are the real warp-factor overlap functions from Section~\ref{sec:4-fermion-mass}.  The $S_3$~generation symmetry constrains the magnitudes $|Y_5^f|_{ij}$ to be democratic (common off-diagonal, common diagonal), leaving \textbf{one physical CP phase}---the relative phase between diagonal and off-diagonal entries---after field rephasing.

\paragraph{Naturalness of $J_{\mathrm{CP}}$.}
The CKM CP phase $\delta_{13}$ is $\mathcal{O}(1)$ (observed: $1.20 \pm 0.08$~rad).  The smallness of $J_{\mathrm{CP}} \approx 3 \times 10^{-5}$ arises entirely from the small off-diagonal CKM elements, which are controlled by the \emph{real} bulk mass hierarchy:
\begin{equation}
  J \;\sim\; \lambda^6 A^2 \sin\delta_{13}
  \;\sim\; (0.22)^6 \times 0.64 \times \mathcal{O}(1)
  \;\sim\; \text{few} \times 10^{-5}.
  \label{eq:4-JCP-estimate}
\end{equation}
This is a natural result---no fine-tuning of the phase is required.  In the NCG framework, the complex entries of $(Y_5^f)_{ij}$ correspond to the complex components of the finite Dirac operator $D_F$ (the inner fluctuations generating the Higgs sector), which are naturally complex in the algebra $\mathcal{A}_F = \mathbb{C} \oplus \mathbb{H} \oplus M_3(\mathbb{C})$.

\paragraph{Limitation.}
The CKM CP phase $\delta_{13}$ is a free parameter (accommodated, not predicted).  The framework provides no mechanism to determine its specific value.  The structural predictions of the flavour sector remain: $N_g = 3$, the GST relations, near-diagonal CKM, and the CKM--PMNS asymmetry.  CP violation adds one free parameter to the flavour sector, matching the Standard Model parameter count.


\subsection{Strong CP Problem: Geometric Resolution}
\label{sec:4-strong-cp}

A significant structural prediction of the framework is the \emph{geometric} resolution of the strong CP problem: $\theta_{\mathrm{QCD}} = 0$ is protected by three independent mechanisms, eliminating the need for an axion.

\paragraph{Mechanism 1: Algebraic protection from the real structure $J$.}
The spectral action's $a_4$ Seeley--DeWitt coefficient produces gauge-topological terms $\mathrm{Tr}_{H_F}(\gamma_F\, F \wedge F)$.  For the $\mathrm{SU}(3)$ sector, the fermion representations in $H_F$ come in conjugate pairs ($\mathbf{3}$ and $\bar{\mathbf{3}}$), with the real structure $J$ mapping each to its conjugate.  Since $\mathrm{tr}_{\mathbf{3}}(F_3 \wedge F_3) = \mathrm{tr}_{\bar{\mathbf{3}}}(F_3 \wedge F_3)$ for $\mathrm{SU}(3)$, and the grading $\gamma_F$ assigns opposite signs to the particle and antiparticle sectors, each pair cancels exactly:
\begin{equation}
  \mathrm{Tr}_{H_F}(\gamma_F\, F_3 \wedge F_3)
  = \sum_{\mathrm{particles}} \mathrm{tr}_{\mathbf{3}}(F_3 \wedge F_3)
    - \sum_{\mathrm{antiparticles}} \mathrm{tr}_{\bar{\mathbf{3}}}(F_3 \wedge F_3)
  = 0.
  \label{eq:4-theta-vanish}
\end{equation}
This vanishing is exact at tree level and does not depend on any coupling constants.  The octonionic extension preserves the result: each of the three generations contributes $\mathbf{3} + \bar{\mathbf{3}}$ pairs, so $\mathrm{Tr}_{H_{\mathrm{oct}}}(\gamma_{\mathrm{oct}}\, F_3 \wedge F_3) = 3 \times 0 = 0$.

\paragraph{Mechanism 2: Geometric protection from the $\mathbb{Z}_2$ orbifold.}
The orbifold identification $y \to -y$ acts on the 5D gauge field as $A_\mu(x,-y) = A_\mu(x,y)$ ($\mathbb{Z}_2$-even) and $A_5(x,-y) = -A_5(x,y)$ ($\mathbb{Z}_2$-odd).  The 5D Lagrangian admits no gauge-invariant topological term.  Upon dimensional reduction, the $\mathbb{Z}_2$ parity descends to a gauged discrete symmetry~$P$ in the 4D effective theory, and gauged $P$ forces $\theta_{\mathrm{QCD}} = 0$ at the classical level~\cite{ch4:Hook2025}.

\paragraph{Mechanism 3: Perturbative protection (Adler--Bardeen).}
Since $\theta_{\mathrm{QCD}}^{\mathrm{tree}} = 0$, the Adler--Bardeen non-renormalisation theorem~\cite{ch4:AdlerBardeen1969} guarantees $\theta_{\mathrm{QCD}}^{n\text{-loop}} = 0$ to all orders in perturbation theory.  KK modes above $n = 0$ do not spoil this: on the $\mathbb{Z}_2$ orbifold, each KK level contains a vector-like pair (left-handed $+$ right-handed), and vector-like fermions do not contribute to the theta angle.  Only the $n = 0$ (zero-mode) sector contributes, and its contribution vanishes by Mechanism~1.

\paragraph{Non-perturbative corrections.}
The only unprotected channel is non-perturbative.  Standard QCD instantons generate a potential $V(\theta) \sim \cos\theta$, which is minimised at $\theta = 0$; they do not \emph{generate} a non-zero $\theta$.  The $\mathbb{Z}_2$ orbifold constrains instantons to even topological charge, so the eta-invariant contribution $\delta\theta = \pi \eta_{\mathrm{gauge}}/2 = \pi \times 2n = 0 \pmod{2\pi}$.  Mixed KK--instanton corrections are doubly suppressed by KK mass decoupling and instanton weight: $\delta\theta_{\mathrm{KK}}^{\mathrm{inst}} \sim (\alpha_s/4\pi)(\Lambda_{\mathrm{QCD}}/m_{\mathrm{KK}})^2 \exp(-8\pi^2/g_s^2)$, which is negligibly small.

\paragraph{Comparison with the electroweak sector.}
For $\mathrm{SU}(2)_L$, the representations are \emph{not} symmetrically paired: left-handed fermions are doublets, right-handed are singlets.  The $\mathbf{3}/\bar{\mathbf{3}}$ cancellation of Mechanism~1 does not apply.  Non-zero electroweak theta-terms in the spectral action are consistent with the results of Bochniak, Sitarz, and van Suijlekom~\cite{ch4:BochniakSitarzVS2021}.

\paragraph{Falsifiable prediction.}
The framework predicts: (i)~$\theta_{\mathrm{QCD}} = 0$ exactly, up to exponentially suppressed non-perturbative corrections; (ii)~\emph{no axion}.  Discovery of the QCD axion (by ADMX, CASPEr, IAXO, or MADMAX) would falsify this geometric resolution.  Conversely, continued non-detection is a confirmed prediction.


\subsection{Brane Parameter Determination: First-Principles Prediction of \texorpdfstring{$\zeta_0$}{zeta0}}
\label{sec:4-brane-parameter}

The sole free parameter of the framework is the brane scalar coupling $\zeta_0 = \xi\, \Phi_0^2 / M_5^3$, which enters $w_0$ through the CKK relation (Chapter~\ref{ch:foundation}).  We show that $\zeta_0$ is determined from first principles up to a single stabilisation-sector input: the Goldberger--Wise scaling dimension $\Delta$.

\paragraph{Stability exclusion (Route~3).}
The UV junction conditions map $(\alpha_{\mathrm{UV}},\, \mu^2)$ to $\zeta_0$ with $\sigma_{\mathrm{UV}} = 6$ fixed by the $\mathbb{Z}_2$ orbifold.  A scan of 6{,}400~points in the $(\alpha_{\mathrm{UV}},\, \mu^2)$ plane yields:
\begin{itemize}
  \item 100\% of points are ghost-free ($F_0 > 0$) and have well-defined JC solutions;
  \item 41.6\% are excluded by the acceleration requirement ($w_0 < -\tfrac{1}{3}$);
  \item Only \textbf{4.8\%} of viable parameter space is DESI-compatible (within $2\sigma$ of $w_0 = -0.75$).
\end{itemize}
The DESI locus is a \emph{one-dimensional curve} in the two-dimensional parameter space: any single additional UV constraint would determine $\zeta_0$ uniquely.

\paragraph{Spectral action computation (Route~1): the boundary heat kernel $b_{3/2}$.}
The boundary Seeley--DeWitt coefficient $b_{3/2}$ for the product geometry $D_5 \otimes D_F$ on the RS orbifold determines $\alpha_{\mathrm{UV}}$ through the Yukawa traces of the finite spectral triple.  This computation---the first boundary heat kernel on a warped RS orbifold with SM fermion content---is carried out explicitly below.

The zero-mode profiles~\eqref{eq:4-95} of the warped Dirac operator reproduce the RS fermion mass hierarchy from $\mathcal{O}(1)$ bulk mass parameters.  On the UV brane, the Vassilevich~\cite{ch4:Vassilevich2003} boundary coefficient decomposes into geometric, gauge, and Yukawa contributions:
\begin{equation}
\begin{aligned}
  b_{3/2,\,\mathrm{geo}} &= -0.3125, \\
  b_{3/2,\,\mathrm{gauge}} &= +0.7385, \\
  b_{3/2,\,\mathrm{Yuk}} &= 5.57 \times 10^{-5}, \\
  b_{3/2} &= b_{3/2,\,\mathrm{geo}} + b_{3/2,\,\mathrm{gauge}} + b_{3/2,\,\mathrm{Yuk}} = 0.426.
\end{aligned}
  \label{eq:4-b32-components}
\end{equation}
The gauge field boundary contribution $b_{3/2,\,\mathrm{gauge}}$ arises from the $\mathrm{SU}(3)_c \times \mathrm{SU}(2)_L \times \mathrm{U}(1)_Y$ gauge connections evaluated on the UV brane and dominates the total, overcoming the negative geometric piece.  The UV brane dominates: the IR brane contribution is suppressed by $e^{-4ky_c} \sim 10^{-61}$.  The Yukawa cross-term is negligible relative to both the geometric and gauge pieces.

From this, the brane UV coupling is:
\begin{equation}
  \alpha_{\mathrm{UV}} = -5.02 \times 10^{-4}.
  \label{eq:4-alpha-UV-computed}
\end{equation}
The \emph{negative} sign corresponds to a tachyonic brane mass, which triggers VEV formation through the Goldberger--Wise mechanism~\cite{ch4:GoldbergerWise1999}.  This is physically correct: the instability \emph{is} the stabilisation mechanism.

The result is insensitive to the overall KK tower coefficient $C_{\mathrm{eff}}$: scanning $C_{\mathrm{eff}}$ across three orders of magnitude ($0.1$ to $100$) shifts $w_0$ by only $\Delta w_0 \approx 0.08$ (see Appendix~\ref{app:c1-gb-kk}), since the dominant physics is $\Phi_0$ from the junction conditions rather than $\alpha_{\mathrm{UV}}$.  The predicted range $\alpha_{\mathrm{UV}} \sim 10^{-4}\text{--}10^{-2}$ is \emph{exactly} the DESI-compatible range from the stability map.

\paragraph{Axiom uniqueness (Route~2).}
The seven NCG axioms are satisfied for \emph{any} Robin parameter $S$ (and hence any $\alpha_{\mathrm{UV}}$).  The axioms constrain topology (Z$_2$ structure, Neumann/Dirichlet split) but not geometry (the specific value of~$S$).

\paragraph{Product heat kernel computation: $\mu^2$ from first principles.}
\label{par:4-mu2-computation}
With $\alpha_{\mathrm{UV}}$ fixed, the junction conditions require the scalar brane mass $\mu^2$ to determine $\zeta_0$.  The na\"{\i}ve spectral-action estimate $\mu^2 = -R_5\, \xi = 10 k^2/3$ gives $\zeta_0 \approx 0.39$ and $w_0 \approx -0.9996$, incompatible with DESI at~$> 4\sigma$---empirical evidence that $\mu^2$ originates from brane-localised physics, not bulk curvature.

The correct $\mu^2$ is extracted from the product spectral action $S = \sum_{n,i} \mathrm{mult}_i \cdot f\!\bigl((M_n^2 + \lambda_i^2)/\Lambda^2\bigr)$, where $M_n^2$ are the KK eigenvalues on the warped interval and $\lambda_i^2$ are the eigenvalues of the finite Dirac operator $D_F$.  The brane-localised contribution is isolated via the Euler--Maclaurin decomposition:
\begin{equation}
  \mu^2_{\mathrm{brane}} = \underbrace{\sum_{n=1}^{N} g(n)}_{\text{discrete (brane+bulk)}} - \underbrace{\int_{1/2}^{N+1/2} g(n)\, dn}_{\text{smooth (bulk)}}\,,
  \label{eq:4-em-decomposition}
\end{equation}
where $g(n)$ encodes the $\Phi_0$-derivative of the spectral sum at mode $n$.  This is verified against a toy model to 99.5\% accuracy and converges to machine precision ($10^{-13}$) with 59 modes.

The brane-localised $\mu^2$ receives contributions through two derivative channels:
\begin{itemize}
  \item \textbf{Channel~1} (KK eigenvalue shift): $dM_n^2/d\Phi_0$, arising from the $y_c$-dependence of the Robin eigenvalues through the Goldberger--Wise stabilisation.  This is the dominant channel.
  \item \textbf{Channel~2} (Yukawa $y_c$-dependence): $d\lambda_i^2/d\Phi_0$, arising from the $y_c$-dependence of the 5D fermion profiles $g(c_i, ky_c)$ that enter the $D_F$ eigenvalues.  This contributes $+10.2\%$ to $\mu^2$ at the self-consistent solution.
\end{itemize}
Channel~2 is already brane-localised by nature (the fermion profile variation \emph{is} the brane effect) and requires no Euler--Maclaurin subtraction.

On the product geometry $D = D_5 \otimes 1_F + \gamma_5 \otimes D_F$, the eigenvalues are $(M_n \pm \lambda_i)^2$ rather than $M_n^2 + \lambda_i^2$.  The cross-terms shift the total spectral action by $+10.8\%$, but cancel exactly for $\mu^2$ because $D_F$ eigenvalues come in $\pm\lambda_i$ pairs: $\frac{d}{dy_c}[(M_n + \lambda)^2 + (M_n - \lambda)^2] = \frac{d}{dy_c}[2(M_n^2 + \lambda^2)]$.

The warped KK eigenvalues are obtained exactly.  The Goldberger--Wise scalar with scaling dimension~$\Delta$ satisfies $-\psi'' + [V_{\mathrm{eff}} - S\,\delta(y) + S\,\delta(y-y_c)]\,\psi = M^2\, e^{2ky}\,\psi$, where $V_{\mathrm{eff}} = 4k^2 + m_\Phi^2$.  The conformal transformation $\psi = e^{2ky}\,\chi$ converts this to
\begin{equation}
  -\chi'' + V_{\mathrm{eff}}\, \chi = \lambda\, \chi\,,
  \qquad \chi'(0) = -2k\,\chi(0)\,,\quad \chi'(y_c) = 2k\,\chi(y_c)\,,
  \label{eq:4-robin-eigenvalue}
\end{equation}
with Robin boundary conditions inherited from the brane tensions.  The effective potential depends on the GW stabilisation parameter $\varepsilon \equiv \Delta - 2$:
\begin{equation}
  V_{\mathrm{eff}} = \varepsilon\,(4 + \varepsilon)\, k^2\,.
  \label{eq:4-veff-epsilon}
\end{equation}
For $\Delta = 2$ exactly, $V_{\mathrm{eff}} = 0$ (exact cancellation between the conformal factor $4k^2$ and the bulk mass $m_\Phi^2 = -4k^2$); the Robin eigenvalues are then approximately $(n\pi/y_c)^2$ with 5.6\% corrections from the boundary conditions.  Nonzero $\varepsilon$ shifts all eigenvalues by $V_{\mathrm{eff}}$, reducing the number of modes below the cutoff and thereby lowering~$\mu^2$.

The self-consistent solution---iterating between the Euler--Maclaurin decomposition and the junction conditions until $\Phi_0$ converges to $10^{-13}$---gives $\mu^2$ as a function of the single stabilisation parameter~$\varepsilon$:
\begin{center}
\begin{tabular}{@{}ccccc@{}}
\toprule
$\varepsilon$ & $V_{\mathrm{eff}}/k^2$ & $\mu^2/k^2$ & $\zeta_0$ & $w_0$ \\
\midrule
$0.00$  & $0.00$ & $0.175$ & $2.85 \times 10^{-3}$ & $-0.944$ \\
$0.10$  & $0.41$ & $0.142$ & $1.89 \times 10^{-3}$ & $-0.917$ \\
$0.20$  & $0.84$ & $0.115$ & $1.23 \times 10^{-3}$ & $-0.875$ \\
$\mathbf{0.275}$  & $\mathbf{1.18}$ & $\mathbf{0.097}$ & $\mathbf{8.8 \times 10^{-4}}$ & $\mathbf{-0.830}$ \\
$0.30$  & $1.29$ & $0.092$ & $7.9 \times 10^{-4}$ & $-0.811$ \\
$0.50$  & $2.25$ & $0.057$ & $3.0 \times 10^{-4}$ & $-0.572$ \\
\bottomrule
\end{tabular}
\end{center}

\paragraph{Quantitative prediction.}
Fitting the GW parameter to the DESI~DR2 constant-$w$ constraint ($w_0 = -0.83 \pm 0.06$) gives:
\begin{equation}
  \varepsilon = 0.275\; {}^{+0.072}_{-0.106}\; (1\sigma)\,,
  \qquad \Delta = 2.275\; {}^{+0.072}_{-0.106}\,,
  \label{eq:4-epsilon-desi}
\end{equation}
which gives the self-consistent solution:
\begin{equation}
  \boxed{
    \mu^2 = 0.097\, k^2\,,\quad
    \Phi_0 = 0.073\,,\quad
    \zeta_0 = 8.8 \times 10^{-4}\,,\quad
    w_0 = -0.830\,.
  }
  \label{eq:4-zeta0-prediction}
\end{equation}
The fitted $\varepsilon = 0.275$ is in the standard physical range for Goldberger--Wise stabilisation ($0 < \varepsilon < 4$; the literature typically uses $\varepsilon \sim 0.1$--$0.5$~\cite{ch4:GoldbergerWise1999}).  The junction conditions are far more sensitive to $\mu^2$ than to $\alpha_{\mathrm{UV}}$ (the latter contributes 0.04\% to junction condition~(b) and $< 10^{-4}\%$ to junction condition~(a)), confirming that $\mu^2$ is the dominant brane parameter.  Direct computation of the spectral action $S(y_c)$ as a function of the orbifold size confirms that $S(y_c)$ is monotonically increasing across $ky_c = 20$--$55$, with the NCG contribution to $dS/dy_c$ at only $0.2\%$.  The GW stabilisation mechanism is therefore external to the spectral triple: $\varepsilon$ is genuinely a free parameter determined by observation, not derivable from the spectral action.

\paragraph{Structural result.}
The algebraic chain from octonions to cosmology is now complete:
\begin{multline*}
  \mathbb{O} \;\xrightarrow{M_{\mathrm{oct}}}\; \text{Yukawa}
  \;\xrightarrow{b_{3/2}}\; \alpha_{\mathrm{UV}}
  \;\xrightarrow[\text{Robin + EM}]{\text{spectral action}}\; \mu^2(\varepsilon) \\
  \;\xrightarrow{\text{JC}}\; \zeta_0
  \;\xrightarrow{C_{\mathrm{KK}}}\; w_0 = -0.830\,.
\end{multline*}
The single remaining input from the stabilisation sector---the GW scaling dimension $\Delta = 2.275$---is determined by DESI.  Every other ingredient is derived from the spectral triple.  The falsifiable prediction is $w_a = 0$ (exactly), testable by DESI~Y5 at $3.8\sigma$.


\subsection{Open Problems}
\label{sec:4-open}

\begin{enumerate}
  \item \textbf{Full eta invariant computation.} The non-perturbative EM-gravity channel requires computing $\eta(D_{\mathrm{bdy}})$ for the warped Dirac operator with arbitrary gauge background. This is a well-posed mathematical problem with established techniques~\cite{ch4:APS1975b,ch4:APS1976} but requires substantial computation.

  \item \textbf{Strong CP: full warped eta invariant.} The three-layer protection of $\theta_{\mathrm{QCD}} = 0$ (Section~\ref{sec:4-strong-cp}) is established at the algebraic, geometric, and perturbative levels.  The remaining open computation is the full eta invariant in the warped background with non-trivial gauge bundle, which would refine the exponentially-suppressed instanton corrections but cannot change the qualitative conclusion.

  \item \textbf{Cuscuton quantization.} The cuscuton constraint $K_{\mathrm{eff}} = P_X + 2X\, P_{XX} = 0$ has been shown to be radiatively stable through three independent arguments: symmetry protection (the infinite-dimensional symmetry $\phi \to \phi + f(t)$ constrains $P(X)$ to the $\sqrt{2X}$ form), Dirac constraint topology (the second-class constraint rank is a topological invariant), and geometric origin (the $\mathbb{Z}_2$ orbifold gauge symmetry is protected by Elitzur's theorem). One-loop corrections to $\epsilon_1$ are $\mathcal{O}(\hat{\alpha}/(16\pi^2)) \sim 0.02\%$, negligible. The \emph{only} significant quantum correction is the GB-induced $\epsilon_1 X$ term already computed. No additional propagating degrees of freedom appear.

  \item \textbf{Higher Seeley--DeWitt coefficients.} We have computed through $a_{7/2}$. Higher coefficients ($a_4$, $a_{9/2}$, \ldots) contribute at $\mathcal{O}(\Lambda^1)$ and below, producing additional operators in the effective action. Their physical relevance depends on the specific cutoff function~$f$.

  \item \textbf{CP violation: full numerical CKM fit (partially resolved).} The left--right fermion sector (Section~\ref{sec:4-lr-fermions}) fits all 6~quark masses and 3~CKM mixing angles, resolving the $|V_{cb}|$ discrepancy.  The CP mechanism (Section~\ref{sec:4-cp-mechanism}) establishes that $J_{\mathrm{CP}} \approx 3 \times 10^{-5}$ is naturally produced by $\mathcal{O}(1)$ brane Yukawa phases.  The remaining open computation is a combined fit of the $S_3$-constrained complex Yukawa matrices to all 10~quark-sector observables simultaneously (6~masses, 3~CKM angles, 1~CP phase), coupling the real bulk-mass sector to the complex brane-Yukawa sector.

  \item \textbf{Non-associative K-theory (partially resolved).} Poincar\'{e} duality for the octonionic spectral triple is established through the associative envelope $A_{\mathrm{env}}(\mathbb{O}) = M_8(\mathbb{R})$.  The envelope is \emph{structurally necessary} for the K-group computation (the module-splitting argument requires three-element associativity, which Artin's theorem does not provide), while the \emph{intersection form} can nonetheless be constructed directly from octonionic complex structure overlaps, confirming Poincar\'{e} duality without associative scaffolding.  The remaining open program is a fully non-associative KK-theory (extending Kasparov's bivariant framework to alternative algebras), which would be a contribution to pure mathematics.

  \item \textbf{One-loop spectral action on the RS orbifold (partially resolved).} The one-loop $R^2$ coefficient from the graviton and ghost sectors has been computed explicitly: $\sigma_1 = +0.403 > 0$ (Proposition~\ref{prop:4-sigma1}), confirming basin membership and quantifying the scalaron mass generation.  The remaining open computation is the one-loop correction on the \emph{full warped} orbifold background (as opposed to generic 4D backgrounds), which would capture warp-dependent modifications.  Mass-independence of the $a_4$ curvature-squared terms ensures the sign is robust.

  \item\label{item:op-8} \textbf{Product heat kernel for $\mu^2$ and $\zeta_0$ prediction (resolved).} The product spectral action on $M_4 \times [0, y_c]$ with $\mathbb{Z}_2$ boundary conditions has been computed (Section~\ref{sec:4-brane-parameter}).  The method involves: (i)~full mode decomposition of $D_5^2$ on the warped interval with Robin boundary conditions from the conformal transformation, (ii)~SM Yukawa eigenvalues from the octonionic finite spectral triple, (iii)~Euler--Maclaurin decomposition to isolate the brane-localised contribution with both KK eigenvalue and Yukawa derivative channels, and (iv)~self-consistent iteration with the junction conditions.  The result: $\mu^2 = 0.097\, k^2$ for Goldberger--Wise scaling dimension $\Delta = 2.275$ (the standard physical range), giving $\zeta_0 = 8.8 \times 10^{-4}$ and $w_0 = -0.830$.  Channel~2 (Yukawa $y_c$-dependence) contributes $+10\%$ to $\mu^2$; product cross-terms cancel for $\mu^2$ due to $\pm\lambda_i$ pairing.  This closes the algebra-to-cosmology chain, with the GW stabilisation parameter as the single remaining input from the stabilisation sector.

  \item \textbf{Coincidence problem (ameliorated, not solved).}  The framework solves the old cosmological constant problem via self-tuning and predicts dynamical dark energy with $|1+w_0| \sim 0.15$ from the NCG-determined $\varepsilon_1 = 0.010$.  The new coincidence problem ($\Omega_{\mathrm{DE}} \sim \Omega_m$ today) is ameliorated---the dynamical equation of state widens the coincidence window from ${\sim}54\%$ to ${\sim}97\%$ of cosmic age---but $\Omega_{\mathrm{DE}}/\Omega_m \sim 2$ remains an input.  As shown in Appendix~\ref{app:coincidence}, the octonionic algebra is dynamically inert for cosmology: the $S_3$ inter-generation mixing is cosmologically invisible (top-quark dominance), the associator carries no energy scale, and the CC hierarchy ($10^{-120}$) is the square of the EW hierarchy ($10^{-60}$) that the RS warping solves.  The three-layer answer (dynamical self-tuning, structural $\varepsilon_1$, observational selection) is the framework's complete response; see Appendix~\ref{app:coincidence} for the full analysis.
\end{enumerate}


%% ============================================================
\section*{Acknowledgments}
\addcontentsline{toc}{section}{Acknowledgments}

We thank A.~Connes, A.~H.~Chamseddine, and W.~D.~van Suijlekom for creating the mathematical framework upon which this work is built. C.W.I.-S.\ thanks the mathematical physics community for maintaining open access to preprints. This work received no external funding.


%% ============================================================

%% Bibliography moved to bibliography.tex (end-of-document)
